% Options for packages loaded elsewhere
\PassOptionsToPackage{unicode}{hyperref}
\PassOptionsToPackage{hyphens}{url}
\PassOptionsToPackage{dvipsnames,svgnames,x11names}{xcolor}
%
\documentclass[
  11pt,
]{article}

\usepackage{amsmath,amssymb}
\usepackage{iftex}
\ifPDFTeX
  \usepackage[T1]{fontenc}
  \usepackage[utf8]{inputenc}
  \usepackage{textcomp} % provide euro and other symbols
\else % if luatex or xetex
  \usepackage{unicode-math}
  \defaultfontfeatures{Scale=MatchLowercase}
  \defaultfontfeatures[\rmfamily]{Ligatures=TeX,Scale=1}
\fi
\usepackage{lmodern}
\ifPDFTeX\else  
    % xetex/luatex font selection
    \setmainfont[]{Arial}
\fi
% Use upquote if available, for straight quotes in verbatim environments
\IfFileExists{upquote.sty}{\usepackage{upquote}}{}
\IfFileExists{microtype.sty}{% use microtype if available
  \usepackage[]{microtype}
  \UseMicrotypeSet[protrusion]{basicmath} % disable protrusion for tt fonts
}{}
\makeatletter
\@ifundefined{KOMAClassName}{% if non-KOMA class
  \IfFileExists{parskip.sty}{%
    \usepackage{parskip}
  }{% else
    \setlength{\parindent}{0pt}
    \setlength{\parskip}{6pt plus 2pt minus 1pt}}
}{% if KOMA class
  \KOMAoptions{parskip=half}}
\makeatother
\usepackage{xcolor}
\usepackage[a4paper, margin=0.8in]{geometry}
\setlength{\emergencystretch}{3em} % prevent overfull lines
\setcounter{secnumdepth}{5}
% Make \paragraph and \subparagraph free-standing
\makeatletter
\ifx\paragraph\undefined\else
  \let\oldparagraph\paragraph
  \renewcommand{\paragraph}{
    \@ifstar
      \xxxParagraphStar
      \xxxParagraphNoStar
  }
  \newcommand{\xxxParagraphStar}[1]{\oldparagraph*{#1}\mbox{}}
  \newcommand{\xxxParagraphNoStar}[1]{\oldparagraph{#1}\mbox{}}
\fi
\ifx\subparagraph\undefined\else
  \let\oldsubparagraph\subparagraph
  \renewcommand{\subparagraph}{
    \@ifstar
      \xxxSubParagraphStar
      \xxxSubParagraphNoStar
  }
  \newcommand{\xxxSubParagraphStar}[1]{\oldsubparagraph*{#1}\mbox{}}
  \newcommand{\xxxSubParagraphNoStar}[1]{\oldsubparagraph{#1}\mbox{}}
\fi
\makeatother

\usepackage{color}
\usepackage{fancyvrb}
\newcommand{\VerbBar}{|}
\newcommand{\VERB}{\Verb[commandchars=\\\{\}]}
\DefineVerbatimEnvironment{Highlighting}{Verbatim}{commandchars=\\\{\}}
% Add ',fontsize=\small' for more characters per line
\usepackage{framed}
\definecolor{shadecolor}{RGB}{248,248,248}
\newenvironment{Shaded}{\begin{snugshade}}{\end{snugshade}}
\newcommand{\AlertTok}[1]{\textcolor[rgb]{0.94,0.16,0.16}{#1}}
\newcommand{\AnnotationTok}[1]{\textcolor[rgb]{0.56,0.35,0.01}{\textbf{\textit{#1}}}}
\newcommand{\AttributeTok}[1]{\textcolor[rgb]{0.13,0.29,0.53}{#1}}
\newcommand{\BaseNTok}[1]{\textcolor[rgb]{0.00,0.00,0.81}{#1}}
\newcommand{\BuiltInTok}[1]{#1}
\newcommand{\CharTok}[1]{\textcolor[rgb]{0.31,0.60,0.02}{#1}}
\newcommand{\CommentTok}[1]{\textcolor[rgb]{0.56,0.35,0.01}{\textit{#1}}}
\newcommand{\CommentVarTok}[1]{\textcolor[rgb]{0.56,0.35,0.01}{\textbf{\textit{#1}}}}
\newcommand{\ConstantTok}[1]{\textcolor[rgb]{0.56,0.35,0.01}{#1}}
\newcommand{\ControlFlowTok}[1]{\textcolor[rgb]{0.13,0.29,0.53}{\textbf{#1}}}
\newcommand{\DataTypeTok}[1]{\textcolor[rgb]{0.13,0.29,0.53}{#1}}
\newcommand{\DecValTok}[1]{\textcolor[rgb]{0.00,0.00,0.81}{#1}}
\newcommand{\DocumentationTok}[1]{\textcolor[rgb]{0.56,0.35,0.01}{\textbf{\textit{#1}}}}
\newcommand{\ErrorTok}[1]{\textcolor[rgb]{0.64,0.00,0.00}{\textbf{#1}}}
\newcommand{\ExtensionTok}[1]{#1}
\newcommand{\FloatTok}[1]{\textcolor[rgb]{0.00,0.00,0.81}{#1}}
\newcommand{\FunctionTok}[1]{\textcolor[rgb]{0.13,0.29,0.53}{\textbf{#1}}}
\newcommand{\ImportTok}[1]{#1}
\newcommand{\InformationTok}[1]{\textcolor[rgb]{0.56,0.35,0.01}{\textbf{\textit{#1}}}}
\newcommand{\KeywordTok}[1]{\textcolor[rgb]{0.13,0.29,0.53}{\textbf{#1}}}
\newcommand{\NormalTok}[1]{#1}
\newcommand{\OperatorTok}[1]{\textcolor[rgb]{0.81,0.36,0.00}{\textbf{#1}}}
\newcommand{\OtherTok}[1]{\textcolor[rgb]{0.56,0.35,0.01}{#1}}
\newcommand{\PreprocessorTok}[1]{\textcolor[rgb]{0.56,0.35,0.01}{\textit{#1}}}
\newcommand{\RegionMarkerTok}[1]{#1}
\newcommand{\SpecialCharTok}[1]{\textcolor[rgb]{0.81,0.36,0.00}{\textbf{#1}}}
\newcommand{\SpecialStringTok}[1]{\textcolor[rgb]{0.31,0.60,0.02}{#1}}
\newcommand{\StringTok}[1]{\textcolor[rgb]{0.31,0.60,0.02}{#1}}
\newcommand{\VariableTok}[1]{\textcolor[rgb]{0.00,0.00,0.00}{#1}}
\newcommand{\VerbatimStringTok}[1]{\textcolor[rgb]{0.31,0.60,0.02}{#1}}
\newcommand{\WarningTok}[1]{\textcolor[rgb]{0.56,0.35,0.01}{\textbf{\textit{#1}}}}

\providecommand{\tightlist}{%
  \setlength{\itemsep}{0pt}\setlength{\parskip}{0pt}}\usepackage{longtable,booktabs,array}
\usepackage{calc} % for calculating minipage widths
% Correct order of tables after \paragraph or \subparagraph
\usepackage{etoolbox}
\makeatletter
\patchcmd\longtable{\par}{\if@noskipsec\mbox{}\fi\par}{}{}
\makeatother
% Allow footnotes in longtable head/foot
\IfFileExists{footnotehyper.sty}{\usepackage{footnotehyper}}{\usepackage{footnote}}
\makesavenoteenv{longtable}
\usepackage{graphicx}
\makeatletter
\newsavebox\pandoc@box
\newcommand*\pandocbounded[1]{% scales image to fit in text height/width
  \sbox\pandoc@box{#1}%
  \Gscale@div\@tempa{\textheight}{\dimexpr\ht\pandoc@box+\dp\pandoc@box\relax}%
  \Gscale@div\@tempb{\linewidth}{\wd\pandoc@box}%
  \ifdim\@tempb\p@<\@tempa\p@\let\@tempa\@tempb\fi% select the smaller of both
  \ifdim\@tempa\p@<\p@\scalebox{\@tempa}{\usebox\pandoc@box}%
  \else\usebox{\pandoc@box}%
  \fi%
}
% Set default figure placement to htbp
\def\fps@figure{htbp}
\makeatother

\makeatletter
\@ifpackageloaded{caption}{}{\usepackage{caption}}
\AtBeginDocument{%
\ifdefined\contentsname
  \renewcommand*\contentsname{Table of contents}
\else
  \newcommand\contentsname{Table of contents}
\fi
\ifdefined\listfigurename
  \renewcommand*\listfigurename{List of Figures}
\else
  \newcommand\listfigurename{List of Figures}
\fi
\ifdefined\listtablename
  \renewcommand*\listtablename{List of Tables}
\else
  \newcommand\listtablename{List of Tables}
\fi
\ifdefined\figurename
  \renewcommand*\figurename{Figure}
\else
  \newcommand\figurename{Figure}
\fi
\ifdefined\tablename
  \renewcommand*\tablename{Table}
\else
  \newcommand\tablename{Table}
\fi
}
\@ifpackageloaded{float}{}{\usepackage{float}}
\floatstyle{ruled}
\@ifundefined{c@chapter}{\newfloat{codelisting}{h}{lop}}{\newfloat{codelisting}{h}{lop}[chapter]}
\floatname{codelisting}{Listing}
\newcommand*\listoflistings{\listof{codelisting}{List of Listings}}
\makeatother
\makeatletter
\makeatother
\makeatletter
\@ifpackageloaded{caption}{}{\usepackage{caption}}
\@ifpackageloaded{subcaption}{}{\usepackage{subcaption}}
\makeatother

\usepackage{bookmark}

\IfFileExists{xurl.sty}{\usepackage{xurl}}{} % add URL line breaks if available
\urlstyle{same} % disable monospaced font for URLs
\hypersetup{
  pdftitle={05\_Microbiome\_analysis},
  pdfauthor={Henry Janse van Rensburg},
  colorlinks=true,
  linkcolor={blue},
  filecolor={Maroon},
  citecolor={Blue},
  urlcolor={Blue},
  pdfcreator={LaTeX via pandoc}}


\title{05\_Microbiome\_analysis}
\author{Henry Janse van Rensburg}
\date{2025-01-31}

\begin{document}
\maketitle

\renewcommand*\contentsname{Table of contents}
{
\hypersetup{linkcolor=}
\setcounter{tocdepth}{2}
\tableofcontents
}

\section{Background}\label{background}

We conducted this experiment to study the effect of mmf1 mutants on
Arabidopsis root microbiome composition and plant growth. We
hypothesized that \emph{mmf1-2} mutants would show altered root
microbiome composition and growth compared to wild-type plants. We grew
Arabidopsis plants on BXplus and BXminus soils and phenotyped them for
rosette area and weighed them at the end of the experiment. We extracted
the phenotypes using ARADEEPOPSIS and analyzed the root microbiome
composition using amplicon sequencing for 16S rRNA, ITS fungal and ITS
oomyctes. We used hamPCR according to @lundberg2021.

\section{Materials and methods}\label{materials-and-methods}

Plant lines used: WT (Col-0), \emph{mmf1-1}, \emph{mmf1-2}

\begin{itemize}
\tightlist
\item
  \textbf{Purpose}: To study the effect of mmf1 mutants on Arabidopsis
  root microbiome composition and plant growth.
\end{itemize}

\section{Important results}\label{important-results}

\begin{itemize}
\item
  \textbf{Plant growth}: The \emph{mmf1-2} mutants showed a significant
  reduction in rosette area compared to the wild-type plants across the
  soils. \emph{mmf1-2} mutants fail to show the BX-growth promotion
  phenotype.
\item
  \textbf{Microbiome composition}: The root microbiome composition of
  \emph{mmf1-2} mutants was significantly different from the wild-type
  plants for bacteria on the BXplus microbiome. The \emph{mmf1-2}
  mutants did not show a significant reduction in bacterial diversity
  compared to the wild-type plants. The fungal and oomycetes diversity
  was not significantly different between the plant lines.
\end{itemize}

\section{Setup}\label{setup}

Purpose: Data Cleaning \& Setup: The script begins by setting working
directories, and loading necessary functions and packages. Path Setup:
Defines a list of file paths for the input data and function files,
making it easier to reference them later. knitr Settings: Ensures that
code outputs are displayed and figures are centered in the generated
document. Reproducibility: Sets a seed for reproducibility.

\begin{Shaded}
\begin{Highlighting}[]
\CommentTok{\# Clear all objects from memory}
\FunctionTok{rm}\NormalTok{(}\AttributeTok{list =} \FunctionTok{ls}\NormalTok{())}

\CommentTok{\# Set working directory to where the data is located}
\CommentTok{\#setwd("/Users/Henry/Desktop/Desktop {-} pseudomonas/Data/2022\_09\_23\_GWAS\_mutant\_screening\_E2/Microbiome\_analysis/BS09\_data/")}

\CommentTok{\# knitr settings: ensures code is echoed and figures are centered}
\NormalTok{knitr}\SpecialCharTok{::}\NormalTok{opts\_chunk}\SpecialCharTok{$}\FunctionTok{set}\NormalTok{(}\AttributeTok{echo =} \ConstantTok{TRUE}\NormalTok{, }\AttributeTok{fig.align =} \StringTok{"center"}\NormalTok{)}
\FunctionTok{options}\NormalTok{(}\AttributeTok{tinytex.verbose =} \ConstantTok{TRUE}\NormalTok{)}

\CommentTok{\# Set seed for reproducibility}
\FunctionTok{set.seed}\NormalTok{(}\DecValTok{100}\NormalTok{)}

\CommentTok{\# Define paths to the different data files and function scripts}
\NormalTok{paths }\OtherTok{\textless{}{-}} \FunctionTok{list}\NormalTok{(}
  \AttributeTok{DESIGN =} \StringTok{"Input/01\_dadapipe/01\_design/design.xlsx"}\NormalTok{,}
  \AttributeTok{bcount =} \StringTok{"Input/01\_dadapipe/13\_output/bacteria/100/bacteria\_COUNT100.tab"}\NormalTok{,}
  \AttributeTok{btaxa =} \StringTok{"Input/01\_dadapipe/13\_output/bacteria/100/bacteria\_TAXA100.tab"}\NormalTok{,}
  \AttributeTok{bseq =} \StringTok{"Input/01\_dadapipe/13\_output/bacteria/100/bacteria\_SEQ100.fasta"}\NormalTok{,}
  \AttributeTok{fcount =} \StringTok{"Input/01\_dadapipe/13\_output/fungi/100/fungi\_COUNT100.tab"}\NormalTok{,}
  \AttributeTok{ftaxa =} \StringTok{"Input/01\_dadapipe/13\_output/fungi/100/fungi\_TAXA100.tab"}\NormalTok{,}
  \AttributeTok{otaxa =} \StringTok{"Input/02\_blast/oomycetes\_TAXA100\_NCBI.tab"}\NormalTok{,}
  \AttributeTok{fseq =} \StringTok{"Input/01\_dadapipe/13\_output/fungi/100/fungi\_SEQ100.fasta"}\NormalTok{,}
  \AttributeTok{fprimers =} \StringTok{"Input/03\_primer\_tracing/03\_R/fOTU\_primer.RDS"}\NormalTok{,}
  \AttributeTok{host\_amplicon =} \StringTok{"Input/At.GI\_amplicon.fasta"}\NormalTok{,}
  \AttributeTok{functions =} \StringTok{"Input/functions/"}
\NormalTok{)}

\CommentTok{\# Set source file location and load required functions}
\ControlFlowTok{if}\NormalTok{ (}\SpecialCharTok{!}\FunctionTok{require}\NormalTok{(rstudioapi)) \{ }\FunctionTok{install.packages}\NormalTok{(}\StringTok{"rstudioapi"}\NormalTok{) \}}
\FunctionTok{library}\NormalTok{(rstudioapi)}
\CommentTok{\#setwd(dirname(getActiveDocumentContext()$path))}

\CommentTok{\# Load custom functions from the specified path}
\FunctionTok{source}\NormalTok{(}\FunctionTok{paste0}\NormalTok{(paths}\SpecialCharTok{$}\NormalTok{functions, }\StringTok{"functions.R"}\NormalTok{))}
\FunctionTok{functions}\NormalTok{(}\AttributeTok{path =}\NormalTok{ paths}\SpecialCharTok{$}\NormalTok{functions)}

\CommentTok{\# Install (if necessary) and load required libraries}
\FunctionTok{libraries}\NormalTok{()}
\FunctionTok{library}\NormalTok{(}\StringTok{"knitr"}\NormalTok{)}
\end{Highlighting}
\end{Shaded}

\subsection{Prepare data}\label{prepare-data}

Purpose: Design Data Preparation: Imports and cleans the design file,
assigning custom colors for visualization and filtering out unnecessary
samples. Data Filtering: Filters the data to focus on bacteria, fungi,
and oomycetes taxa and ensures unique sample numbers.

\begin{Shaded}
\begin{Highlighting}[]
\CommentTok{\# Import the design file that contains experimental design details}
\NormalTok{DESIGN }\OtherTok{\textless{}{-}} \FunctionTok{read\_excel}\NormalTok{(paths}\SpecialCharTok{$}\NormalTok{DESIGN, }\AttributeTok{sheet =} \StringTok{"Sheet1"}\NormalTok{)}

\CommentTok{\# Filter out SYNCOM9 samples from the design data}
\NormalTok{DESIGN }\OtherTok{\textless{}{-}}\NormalTok{ DESIGN[DESIGN}\SpecialCharTok{$}\NormalTok{sub\_libr }\SpecialCharTok{!=} \StringTok{"syncom9"}\NormalTok{, ]}

\CommentTok{\# Ensure the \textquotesingle{}line\textquotesingle{} column is a factor and remove rows with missing \textquotesingle{}line\textquotesingle{} values}
\NormalTok{DESIGN}\SpecialCharTok{$}\NormalTok{line }\OtherTok{\textless{}{-}} \FunctionTok{factor}\NormalTok{(DESIGN}\SpecialCharTok{$}\NormalTok{line, }\AttributeTok{levels =} \FunctionTok{c}\NormalTok{(}\StringTok{"WT"}\NormalTok{, }\StringTok{"159"}\NormalTok{, }\StringTok{"318"}\NormalTok{))}
\NormalTok{DESIGN }\OtherTok{\textless{}{-}}\NormalTok{ DESIGN[}\SpecialCharTok{!}\FunctionTok{is.na}\NormalTok{(DESIGN}\SpecialCharTok{$}\NormalTok{line), ]}

\CommentTok{\# Assign custom colors to lines for visualization purposes}
\NormalTok{DESIGN}\SpecialCharTok{$}\NormalTok{cols\_line }\OtherTok{\textless{}{-}} \ConstantTok{NA}
\NormalTok{DESIGN}\SpecialCharTok{$}\NormalTok{cols\_line[DESIGN}\SpecialCharTok{$}\NormalTok{line }\SpecialCharTok{==} \StringTok{"WT"}\NormalTok{] }\OtherTok{\textless{}{-}} \StringTok{"\#01686c"}
\NormalTok{DESIGN}\SpecialCharTok{$}\NormalTok{cols\_line[DESIGN}\SpecialCharTok{$}\NormalTok{line }\SpecialCharTok{==} \StringTok{"159"}\NormalTok{] }\OtherTok{\textless{}{-}} \StringTok{"\#48afb3"}
\NormalTok{DESIGN}\SpecialCharTok{$}\NormalTok{cols\_line[DESIGN}\SpecialCharTok{$}\NormalTok{line }\SpecialCharTok{==} \StringTok{"318"}\NormalTok{] }\OtherTok{\textless{}{-}} \StringTok{"\#e3913e"}

\CommentTok{\# Create a collapsed color vector for the lines}
\NormalTok{temp }\OtherTok{\textless{}{-}} \FunctionTok{data.frame}\NormalTok{(DESIGN}\SpecialCharTok{$}\NormalTok{line, DESIGN}\SpecialCharTok{$}\NormalTok{cols\_line)}
\NormalTok{temp }\OtherTok{\textless{}{-}}\NormalTok{ plyr}\SpecialCharTok{::}\FunctionTok{ddply}\NormalTok{(temp, }\AttributeTok{.variables =} \StringTok{"DESIGN.cols\_line"}\NormalTok{, }\AttributeTok{.fun =}\NormalTok{ unique)}
\NormalTok{level\_cols\_line }\OtherTok{\textless{}{-}} \FunctionTok{as.character}\NormalTok{(temp[, }\DecValTok{2}\NormalTok{])}
\FunctionTok{names}\NormalTok{(level\_cols\_line) }\OtherTok{\textless{}{-}}\NormalTok{ temp[, }\DecValTok{1}\NormalTok{]}
\FunctionTok{rm}\NormalTok{(temp)}

\CommentTok{\# Assign colors based on experimental conditions (plus/minus)}
\NormalTok{DESIGN}\SpecialCharTok{$}\NormalTok{cols\_condition }\OtherTok{\textless{}{-}} \ConstantTok{NA}
\NormalTok{DESIGN}\SpecialCharTok{$}\NormalTok{cols\_condition[DESIGN}\SpecialCharTok{$}\NormalTok{condition }\SpecialCharTok{==} \StringTok{"minus"}\NormalTok{] }\OtherTok{\textless{}{-}} \StringTok{"gold"}
\NormalTok{DESIGN}\SpecialCharTok{$}\NormalTok{cols\_condition[DESIGN}\SpecialCharTok{$}\NormalTok{condition }\SpecialCharTok{==} \StringTok{"plus"}\NormalTok{] }\OtherTok{\textless{}{-}} \StringTok{"seagreen"}

\CommentTok{\# Collapse condition color vectors similarly}
\NormalTok{temp }\OtherTok{\textless{}{-}} \FunctionTok{data.frame}\NormalTok{(DESIGN}\SpecialCharTok{$}\NormalTok{condition, DESIGN}\SpecialCharTok{$}\NormalTok{cols\_condition)}
\NormalTok{temp }\OtherTok{\textless{}{-}}\NormalTok{ plyr}\SpecialCharTok{::}\FunctionTok{ddply}\NormalTok{(temp, }\AttributeTok{.variables =} \StringTok{"DESIGN.cols\_condition"}\NormalTok{, }\AttributeTok{.fun =}\NormalTok{ unique)}
\NormalTok{level\_cols\_condition }\OtherTok{\textless{}{-}} \FunctionTok{as.character}\NormalTok{(temp[, }\DecValTok{2}\NormalTok{])}
\FunctionTok{names}\NormalTok{(level\_cols\_condition) }\OtherTok{\textless{}{-}}\NormalTok{ temp[, }\DecValTok{1}\NormalTok{]}

\CommentTok{\# Filter out unscaled and duplicate samples}
\NormalTok{DESIGN }\OtherTok{\textless{}{-}}\NormalTok{ DESIGN[}\SpecialCharTok{!}\NormalTok{(}\FunctionTok{grepl}\NormalTok{(}\StringTok{"unscaled"}\NormalTok{, DESIGN}\SpecialCharTok{$}\NormalTok{sub\_libr, }\AttributeTok{fixed =} \ConstantTok{TRUE}\NormalTok{)), ]}
\NormalTok{DESIGN }\OtherTok{\textless{}{-}}\NormalTok{ DESIGN[}\SpecialCharTok{!}\NormalTok{(}\FunctionTok{grepl}\NormalTok{(}\StringTok{"dupl\_for\_At"}\NormalTok{, DESIGN}\SpecialCharTok{$}\NormalTok{sub\_libr, }\AttributeTok{fixed =} \ConstantTok{TRUE}\NormalTok{)), ]}

\CommentTok{\# Assign taxa and split DESIGN into bacteria, fungi, and oomycetes}
\NormalTok{DESIGN}\SpecialCharTok{$}\NormalTok{taxa[DESIGN}\SpecialCharTok{$}\NormalTok{sub\_libr }\SpecialCharTok{==} \StringTok{"oomycetes"}\NormalTok{] }\OtherTok{\textless{}{-}} \StringTok{"oomycetes"}
\NormalTok{bDESIGN }\OtherTok{\textless{}{-}}\NormalTok{ DESIGN[DESIGN}\SpecialCharTok{$}\NormalTok{taxa }\SpecialCharTok{==} \StringTok{"bacteria"}\NormalTok{, ]}
\NormalTok{fDESIGN }\OtherTok{\textless{}{-}}\NormalTok{ DESIGN[DESIGN}\SpecialCharTok{$}\NormalTok{taxa }\SpecialCharTok{==} \StringTok{"fungi"}\NormalTok{, ]}
\NormalTok{oDESIGN }\OtherTok{\textless{}{-}}\NormalTok{ DESIGN[DESIGN}\SpecialCharTok{$}\NormalTok{taxa }\SpecialCharTok{==} \StringTok{"oomycetes"}\NormalTok{, ]}

\CommentTok{\# Reassign sample numbers across taxa}
\ControlFlowTok{if}\NormalTok{ (}\FunctionTok{all}\NormalTok{(bDESIGN}\SpecialCharTok{$}\NormalTok{sample\_nr }\SpecialCharTok{==}\NormalTok{ fDESIGN}\SpecialCharTok{$}\NormalTok{sample\_nr) }\SpecialCharTok{\&}
    \FunctionTok{all}\NormalTok{(bDESIGN}\SpecialCharTok{$}\NormalTok{sample\_nr }\SpecialCharTok{==}\NormalTok{ oDESIGN}\SpecialCharTok{$}\NormalTok{sample\_nr)) \{}
\NormalTok{  bDESIGN}\SpecialCharTok{$}\NormalTok{sample\_nr }\OtherTok{\textless{}{-}} \FunctionTok{paste0}\NormalTok{(}\StringTok{"sample"}\NormalTok{, }\DecValTok{1}\SpecialCharTok{:}\FunctionTok{nrow}\NormalTok{(bDESIGN))}
\NormalTok{  fDESIGN}\SpecialCharTok{$}\NormalTok{sample\_nr }\OtherTok{\textless{}{-}} \FunctionTok{paste0}\NormalTok{(}\StringTok{"sample"}\NormalTok{, }\DecValTok{1}\SpecialCharTok{:}\FunctionTok{nrow}\NormalTok{(fDESIGN))}
\NormalTok{  oDESIGN}\SpecialCharTok{$}\NormalTok{sample\_nr }\OtherTok{\textless{}{-}} \FunctionTok{paste0}\NormalTok{(}\StringTok{"sample"}\NormalTok{, }\DecValTok{1}\SpecialCharTok{:}\FunctionTok{nrow}\NormalTok{(oDESIGN))}
\NormalTok{\}}
\end{Highlighting}
\end{Shaded}

\subsection{Process ASV tables}\label{process-asv-tables}

Purpose: ASV Import: Imports and processes ASV tables for bacteria,
fungi, and oomycetes. Data Cleanup: Adapts column/row names for
consistency and removes OTUs with zero counts.

\begin{Shaded}
\begin{Highlighting}[]
\CommentTok{\# Import ASV tables for bacteria, fungi, and oomycetes}
\NormalTok{bCOUNT }\OtherTok{\textless{}{-}} \FunctionTok{read.delim}\NormalTok{(paths}\SpecialCharTok{$}\NormalTok{bcount, }\AttributeTok{header =} \ConstantTok{TRUE}\NormalTok{, }\AttributeTok{row.names =} \DecValTok{1}\NormalTok{, }\AttributeTok{sep =} \StringTok{"}\SpecialCharTok{\textbackslash{}t}\StringTok{"}\NormalTok{)}
\NormalTok{fCOUNT }\OtherTok{\textless{}{-}} \FunctionTok{read.delim}\NormalTok{(paths}\SpecialCharTok{$}\NormalTok{fcount, }\AttributeTok{header =} \ConstantTok{TRUE}\NormalTok{, }\AttributeTok{row.names =} \DecValTok{1}\NormalTok{, }\AttributeTok{sep =} \StringTok{"}\SpecialCharTok{\textbackslash{}t}\StringTok{"}\NormalTok{)}
\NormalTok{oCOUNT }\OtherTok{\textless{}{-}} \FunctionTok{read.delim}\NormalTok{(paths}\SpecialCharTok{$}\NormalTok{fcount, }\AttributeTok{header =} \ConstantTok{TRUE}\NormalTok{, }\AttributeTok{row.names =} \DecValTok{1}\NormalTok{, }\AttributeTok{sep =} \StringTok{"}\SpecialCharTok{\textbackslash{}t}\StringTok{"}\NormalTok{)}

\CommentTok{\# Transpose the tables}
\NormalTok{bCOUNT }\OtherTok{\textless{}{-}} \FunctionTok{t}\NormalTok{(bCOUNT)}
\NormalTok{fCOUNT }\OtherTok{\textless{}{-}} \FunctionTok{t}\NormalTok{(fCOUNT)}
\NormalTok{oCOUNT }\OtherTok{\textless{}{-}} \FunctionTok{t}\NormalTok{(oCOUNT)}

\CommentTok{\# Adapt column and row names for consistency}
\FunctionTok{colnames}\NormalTok{(bCOUNT) }\OtherTok{\textless{}{-}} \FunctionTok{gsub}\NormalTok{(}\StringTok{"BS09\_"}\NormalTok{, }\StringTok{""}\NormalTok{, }\FunctionTok{colnames}\NormalTok{(bCOUNT))}
\FunctionTok{rownames}\NormalTok{(bCOUNT) }\OtherTok{\textless{}{-}} \FunctionTok{gsub}\NormalTok{(}\StringTok{"OTU"}\NormalTok{, }\StringTok{"bOTU"}\NormalTok{, }\FunctionTok{rownames}\NormalTok{(bCOUNT))}
\FunctionTok{colnames}\NormalTok{(fCOUNT) }\OtherTok{\textless{}{-}} \FunctionTok{gsub}\NormalTok{(}\StringTok{"BS09\_"}\NormalTok{, }\StringTok{""}\NormalTok{, }\FunctionTok{colnames}\NormalTok{(fCOUNT))}
\FunctionTok{rownames}\NormalTok{(fCOUNT) }\OtherTok{\textless{}{-}} \FunctionTok{gsub}\NormalTok{(}\StringTok{"OTU"}\NormalTok{, }\StringTok{"fOTU"}\NormalTok{, }\FunctionTok{rownames}\NormalTok{(fCOUNT))}
\FunctionTok{colnames}\NormalTok{(oCOUNT) }\OtherTok{\textless{}{-}} \FunctionTok{gsub}\NormalTok{(}\StringTok{"BS09\_"}\NormalTok{, }\StringTok{""}\NormalTok{, }\FunctionTok{colnames}\NormalTok{(oCOUNT))}
\FunctionTok{rownames}\NormalTok{(oCOUNT) }\OtherTok{\textless{}{-}} \FunctionTok{gsub}\NormalTok{(}\StringTok{"OTU"}\NormalTok{, }\StringTok{"oOTU"}\NormalTok{, }\FunctionTok{rownames}\NormalTok{(oCOUNT))}

\CommentTok{\# Filter to match design file and remove zero{-}count OTUs}
\NormalTok{bCOUNT }\OtherTok{\textless{}{-}} \FunctionTok{as.data.frame}\NormalTok{(bCOUNT[, }\FunctionTok{match}\NormalTok{(bDESIGN}\SpecialCharTok{$}\NormalTok{bc\_ID, }\FunctionTok{colnames}\NormalTok{(bCOUNT))])}
\NormalTok{bCOUNT }\OtherTok{\textless{}{-}}\NormalTok{ bCOUNT[}\FunctionTok{rowSums}\NormalTok{(bCOUNT) }\SpecialCharTok{\textgreater{}} \DecValTok{0}\NormalTok{, ]}
\NormalTok{fCOUNT }\OtherTok{\textless{}{-}} \FunctionTok{as.data.frame}\NormalTok{(fCOUNT[, }\FunctionTok{match}\NormalTok{(fDESIGN}\SpecialCharTok{$}\NormalTok{bc\_ID, }\FunctionTok{colnames}\NormalTok{(fCOUNT))])}
\NormalTok{fCOUNT }\OtherTok{\textless{}{-}}\NormalTok{ fCOUNT[}\FunctionTok{rowSums}\NormalTok{(fCOUNT) }\SpecialCharTok{\textgreater{}} \DecValTok{0}\NormalTok{, ]}
\NormalTok{oCOUNT }\OtherTok{\textless{}{-}} \FunctionTok{as.data.frame}\NormalTok{(oCOUNT[, }\FunctionTok{match}\NormalTok{(oDESIGN}\SpecialCharTok{$}\NormalTok{bc\_ID, }\FunctionTok{colnames}\NormalTok{(oCOUNT))])}
\NormalTok{oCOUNT }\OtherTok{\textless{}{-}}\NormalTok{ oCOUNT[}\FunctionTok{rowSums}\NormalTok{(oCOUNT) }\SpecialCharTok{\textgreater{}} \DecValTok{0}\NormalTok{, ]}
\end{Highlighting}
\end{Shaded}

\subsection{Data normalization}\label{data-normalization}

Purpose: Bacterial Taxonomy Import: Imports and cleans the bacterial
taxonomy table, removing unwanted OTUs (Cyanobacteria, Mitochondria,
etc.). Label Assignment: Defines labels for bacterial phyla and sets
them as factors.

\begin{Shaded}
\begin{Highlighting}[]
\CommentTok{\# Import and clean bacterial taxonomy table}
\NormalTok{bTAXA }\OtherTok{\textless{}{-}} \FunctionTok{read.table}\NormalTok{(paths}\SpecialCharTok{$}\NormalTok{btaxa, }\AttributeTok{row.names =} \DecValTok{1}\NormalTok{, }\AttributeTok{sep =} \StringTok{"}\SpecialCharTok{\textbackslash{}t}\StringTok{"}\NormalTok{, }\AttributeTok{blank.lines.skip =} \ConstantTok{FALSE}\NormalTok{)}
\FunctionTok{colnames}\NormalTok{(bTAXA) }\OtherTok{\textless{}{-}} \FunctionTok{tolower}\NormalTok{(}\FunctionTok{colnames}\NormalTok{(bTAXA))}
\NormalTok{bTAXA[}\FunctionTok{is.na}\NormalTok{(bTAXA)] }\OtherTok{\textless{}{-}} \StringTok{"unassigned"}
\FunctionTok{rownames}\NormalTok{(bTAXA) }\OtherTok{\textless{}{-}} \FunctionTok{gsub}\NormalTok{(}\StringTok{"OTU"}\NormalTok{, }\StringTok{"bOTU"}\NormalTok{, }\FunctionTok{rownames}\NormalTok{(bTAXA))}
\NormalTok{bTAXA }\OtherTok{\textless{}{-}}\NormalTok{ bTAXA[}\FunctionTok{rownames}\NormalTok{(bTAXA) }\SpecialCharTok{\%in\%} \FunctionTok{rownames}\NormalTok{(bCOUNT), ]}

\CommentTok{\# Remove OTUs of certain taxonomic groups}
\NormalTok{ASVs\_to\_remove }\OtherTok{\textless{}{-}} \FunctionTok{c}\NormalTok{(}\SpecialCharTok{{-}}\FunctionTok{which}\NormalTok{(bTAXA}\SpecialCharTok{$}\NormalTok{phylum }\SpecialCharTok{==} \StringTok{"Cyanobacteria"}\NormalTok{),}
                    \SpecialCharTok{{-}}\FunctionTok{which}\NormalTok{(bTAXA}\SpecialCharTok{$}\NormalTok{family }\SpecialCharTok{==} \StringTok{"Mitochondria"}\NormalTok{),}
                    \SpecialCharTok{{-}}\FunctionTok{which}\NormalTok{(}\FunctionTok{is.na}\NormalTok{(bTAXA)))}

\CommentTok{\# Keep relevant ASVs only and align ASV IDs}
\NormalTok{bTAXA }\OtherTok{\textless{}{-}}\NormalTok{ bTAXA[ASVs\_to\_remove, ]}
\NormalTok{bCOUNT }\OtherTok{\textless{}{-}}\NormalTok{ bCOUNT[}\FunctionTok{rownames}\NormalTok{(bTAXA), ]}
\NormalTok{bTAXA}\SpecialCharTok{$}\NormalTok{ASV\_ID }\OtherTok{\textless{}{-}} \FunctionTok{rownames}\NormalTok{(bTAXA)}

\CommentTok{\# Define labels for Proteobacteria and set labels as factors}
\NormalTok{bTAXA}\SpecialCharTok{$}\NormalTok{labels }\OtherTok{\textless{}{-}} \FunctionTok{as.character}\NormalTok{(bTAXA}\SpecialCharTok{$}\NormalTok{phylum)}
\NormalTok{bTAXA}\SpecialCharTok{$}\NormalTok{labels[bTAXA}\SpecialCharTok{$}\NormalTok{kingdom }\SpecialCharTok{==} \StringTok{"Eukaryota"}\NormalTok{] }\OtherTok{\textless{}{-}} \StringTok{"Host"}
\NormalTok{bTAXA}\SpecialCharTok{$}\NormalTok{labels }\OtherTok{\textless{}{-}} \FunctionTok{as.factor}\NormalTok{(bTAXA}\SpecialCharTok{$}\NormalTok{labels)}
\end{Highlighting}
\end{Shaded}

Purpose: Fungi Taxonomy Import: Imports and cleans the fungi taxonomy
table. Label Assignment: Defines labels for fungi phyla and sets them as
factors.

\begin{Shaded}
\begin{Highlighting}[]
\CommentTok{\#import taxonomy table}
\NormalTok{fTAXA }\OtherTok{\textless{}{-}} \FunctionTok{read.table}\NormalTok{(paths}\SpecialCharTok{$}\NormalTok{ftaxa, }\AttributeTok{row.names=}\DecValTok{1}\NormalTok{, }\AttributeTok{sep=}\StringTok{"}\SpecialCharTok{\textbackslash{}t}\StringTok{"}\NormalTok{, }\AttributeTok{blank.lines.skip =} \ConstantTok{FALSE}\NormalTok{)}

\CommentTok{\#rename}
\FunctionTok{colnames}\NormalTok{(fTAXA) }\OtherTok{\textless{}{-}} \FunctionTok{tolower}\NormalTok{(}\FunctionTok{colnames}\NormalTok{(fTAXA))}
\NormalTok{fTAXA[}\FunctionTok{is.na}\NormalTok{(fTAXA)] }\OtherTok{\textless{}{-}} \StringTok{"unassigned"}
\FunctionTok{rownames}\NormalTok{(fTAXA) }\OtherTok{\textless{}{-}} \FunctionTok{gsub}\NormalTok{(}\StringTok{"OTU"}\NormalTok{,}\StringTok{"fOTU"}\NormalTok{, }\FunctionTok{rownames}\NormalTok{(fTAXA))}

\CommentTok{\#remove OTU with 0 seqs and not fungal ASVs}
\NormalTok{fTAXA }\OtherTok{\textless{}{-}}\NormalTok{ fTAXA[}\FunctionTok{rownames}\NormalTok{(fTAXA) }\SpecialCharTok{\%in\%} \FunctionTok{rownames}\NormalTok{(fCOUNT),]}

\CommentTok{\# define ASVs for removal}
\NormalTok{r1 }\OtherTok{\textless{}{-}} \SpecialCharTok{{-}}\FunctionTok{which}\NormalTok{(fTAXA}\SpecialCharTok{$}\NormalTok{kingdom}\SpecialCharTok{==}\StringTok{"Protista"}\NormalTok{)}
\NormalTok{r2 }\OtherTok{\textless{}{-}} \SpecialCharTok{{-}}\FunctionTok{which}\NormalTok{(fTAXA}\SpecialCharTok{$}\NormalTok{kingdom}\SpecialCharTok{==}\StringTok{"Plantae"}\NormalTok{)}
\NormalTok{r3 }\OtherTok{\textless{}{-}} \SpecialCharTok{{-}}\FunctionTok{which}\NormalTok{(fTAXA}\SpecialCharTok{$}\NormalTok{kingdom}\SpecialCharTok{==}\StringTok{"Protozoa"}\NormalTok{)}
\NormalTok{r4 }\OtherTok{\textless{}{-}} \SpecialCharTok{{-}}\FunctionTok{which}\NormalTok{(fTAXA}\SpecialCharTok{$}\NormalTok{kingdom}\SpecialCharTok{==}\StringTok{"Animalia"}\NormalTok{)}
\NormalTok{r5 }\OtherTok{\textless{}{-}} \SpecialCharTok{{-}}\FunctionTok{which}\NormalTok{(}\FunctionTok{is.na}\NormalTok{(fTAXA)) }\CommentTok{\#rows all na}
\NormalTok{ASVs\_to\_remove }\OtherTok{\textless{}{-}} \FunctionTok{c}\NormalTok{(r1,r2,r3,r4, r5)}
  \ControlFlowTok{if}\NormalTok{(}\FunctionTok{length}\NormalTok{(ASVs\_to\_remove)}\SpecialCharTok{\textgreater{}}\DecValTok{0}\NormalTok{)\{}
\NormalTok{ fTAXA }\OtherTok{\textless{}{-}}\NormalTok{ fTAXA[ASVs\_to\_remove ,]}
\NormalTok{ fCOUNT }\OtherTok{\textless{}{-}}\NormalTok{ fCOUNT[}\FunctionTok{rownames}\NormalTok{(fTAXA),]}
\NormalTok{\}}
\FunctionTok{rm}\NormalTok{(r1,r2,r3,r4,r5,ASVs\_to\_remove)}

\CommentTok{\#add ASV\_ID}
\NormalTok{fTAXA}\SpecialCharTok{$}\NormalTok{ASV\_ID }\OtherTok{\textless{}{-}} \FunctionTok{rownames}\NormalTok{(fTAXA)}
\CommentTok{\#defining ASV colors by phylum (using the bTAXAonomy file)}
\NormalTok{fTAXA}\SpecialCharTok{$}\NormalTok{labels }\OtherTok{\textless{}{-}} \FunctionTok{as.character}\NormalTok{(fTAXA}\SpecialCharTok{$}\NormalTok{phylum)}
\CommentTok{\#Defining fASV colors by phylum (using the taxonomy file)}
\NormalTok{fTAXA}\SpecialCharTok{$}\NormalTok{labels }\OtherTok{\textless{}{-}} \FunctionTok{as.character}\NormalTok{(fTAXA}\SpecialCharTok{$}\NormalTok{phylum)}

\CommentTok{\#add host to labels}
\NormalTok{Host\_OTUs }\OtherTok{\textless{}{-}} \FunctionTok{read.delim}\NormalTok{(paths}\SpecialCharTok{$}\NormalTok{otaxa) }\SpecialCharTok{\%\textgreater{}\%} \FunctionTok{filter}\NormalTok{(genus }\SpecialCharTok{==} \StringTok{"Arabidopsis"}\NormalTok{) }\SpecialCharTok{\%\textgreater{}\%} \FunctionTok{rownames}\NormalTok{() }\CommentTok{\#use NCBI to get At classification}
\NormalTok{fTAXA}\SpecialCharTok{$}\NormalTok{labels[}\FunctionTok{rownames}\NormalTok{(fTAXA)}\SpecialCharTok{==}\FunctionTok{paste0}\NormalTok{(}\StringTok{"f"}\NormalTok{,Host\_OTUs)] }\OtherTok{\textless{}{-}} \StringTok{"Host"}

\CommentTok{\#use NCBI taxonomy}
\NormalTok{Host\_taxonomy }\OtherTok{\textless{}{-}} \FunctionTok{read.delim}\NormalTok{(paths}\SpecialCharTok{$}\NormalTok{otaxa) }\SpecialCharTok{\%\textgreater{}\%} \FunctionTok{filter}\NormalTok{(genus }\SpecialCharTok{==} \StringTok{"Arabidopsis"}\NormalTok{)}
\NormalTok{fTAXA[fTAXA}\SpecialCharTok{$}\NormalTok{labels}\SpecialCharTok{==}\StringTok{"Host"}\NormalTok{,}\DecValTok{1}\SpecialCharTok{:}\DecValTok{8}\NormalTok{] }\OtherTok{\textless{}{-}}\NormalTok{ Host\_taxonomy}

\CommentTok{\#factor}
\NormalTok{fTAXA}\SpecialCharTok{$}\NormalTok{phylum }\OtherTok{\textless{}{-}} \FunctionTok{as.factor}\NormalTok{(fTAXA}\SpecialCharTok{$}\NormalTok{phylum)}
\NormalTok{fTAXA}\SpecialCharTok{$}\NormalTok{labels }\OtherTok{\textless{}{-}} \FunctionTok{as.factor}\NormalTok{(fTAXA}\SpecialCharTok{$}\NormalTok{labels)}
\end{Highlighting}
\end{Shaded}

Purpose: Fungi Taxonomy Import: Imports and cleans the fungi taxonomy
table. Label Assignment: Defines labels for fungi phyla and sets them as
factors.

\begin{Shaded}
\begin{Highlighting}[]
\CommentTok{\#import taxonomy table}
\NormalTok{oTAXA }\OtherTok{\textless{}{-}} \FunctionTok{read.table}\NormalTok{(paths}\SpecialCharTok{$}\NormalTok{otaxa, }\AttributeTok{row.names=}\DecValTok{1}\NormalTok{, }\AttributeTok{sep=}\StringTok{"}\SpecialCharTok{\textbackslash{}t}\StringTok{"}\NormalTok{, }\AttributeTok{blank.lines.skip =} \ConstantTok{FALSE}\NormalTok{)}

\CommentTok{\#rename}
\FunctionTok{colnames}\NormalTok{(oTAXA) }\OtherTok{\textless{}{-}} \FunctionTok{tolower}\NormalTok{(}\FunctionTok{colnames}\NormalTok{(oTAXA))}
\NormalTok{oTAXA[}\FunctionTok{is.na}\NormalTok{(oTAXA)] }\OtherTok{\textless{}{-}} \StringTok{"unassigned"}
\FunctionTok{rownames}\NormalTok{(oTAXA) }\OtherTok{\textless{}{-}} \FunctionTok{gsub}\NormalTok{(}\StringTok{"OTU"}\NormalTok{,}\StringTok{"oOTU"}\NormalTok{, }\FunctionTok{rownames}\NormalTok{(oTAXA))}

\CommentTok{\#remove OTU with 0 seqs and not fungal ASVs}
\NormalTok{oTAXA }\OtherTok{\textless{}{-}}\NormalTok{ oTAXA[}\FunctionTok{rownames}\NormalTok{(oTAXA) }\SpecialCharTok{\%in\%} \FunctionTok{rownames}\NormalTok{(oCOUNT),]}

\CommentTok{\#add ASV\_ID}
\NormalTok{oTAXA}\SpecialCharTok{$}\NormalTok{ASV\_ID }\OtherTok{\textless{}{-}} \FunctionTok{rownames}\NormalTok{(oTAXA)}
\CommentTok{\#defining ASV colors by phylum (using the bTAXAonomy file)}
\NormalTok{oTAXA}\SpecialCharTok{$}\NormalTok{labels }\OtherTok{\textless{}{-}} \FunctionTok{as.character}\NormalTok{(oTAXA}\SpecialCharTok{$}\NormalTok{phylum)}
\CommentTok{\#Defining fASV colors by phylum (using the taxonomy file)}
\NormalTok{oTAXA}\SpecialCharTok{$}\NormalTok{labels }\OtherTok{\textless{}{-}} \FunctionTok{as.character}\NormalTok{(oTAXA}\SpecialCharTok{$}\NormalTok{phylum)}

\CommentTok{\#add host to labels}
\NormalTok{oTAXA}\SpecialCharTok{$}\NormalTok{labels[oTAXA}\SpecialCharTok{$}\NormalTok{genus}\SpecialCharTok{==}\StringTok{"Arabidopsis"}\NormalTok{] }\OtherTok{\textless{}{-}} \StringTok{"Host"}

\CommentTok{\#factor}
\NormalTok{oTAXA}\SpecialCharTok{$}\NormalTok{phylum }\OtherTok{\textless{}{-}} \FunctionTok{as.factor}\NormalTok{(oTAXA}\SpecialCharTok{$}\NormalTok{phylum)}
\NormalTok{oTAXA}\SpecialCharTok{$}\NormalTok{labels }\OtherTok{\textless{}{-}} \FunctionTok{as.factor}\NormalTok{(oTAXA}\SpecialCharTok{$}\NormalTok{labels)}
\end{Highlighting}
\end{Shaded}

\section{Bacterial sequence
filtering}\label{bacterial-sequence-filtering}

Purpose: Filtering: Removes samples with fewer than 60,000 sequences and
ASVs with zero sequences. Summary: Provides a summary of the remaining
bacterial ASVs in terms of total sequences, min, max, median, and number
of OTUs.

\begin{Shaded}
\begin{Highlighting}[]
\CommentTok{\# Initialize a dataframe to store summary statistics}
\NormalTok{df }\OtherTok{\textless{}{-}} \FunctionTok{data.frame}\NormalTok{(}\AttributeTok{sum =} \ConstantTok{NA}\NormalTok{, }\AttributeTok{min =} \ConstantTok{NA}\NormalTok{, }\AttributeTok{max =} \ConstantTok{NA}\NormalTok{, }\AttributeTok{median =} \ConstantTok{NA}\NormalTok{, }\AttributeTok{OTUs =} \ConstantTok{NA}\NormalTok{)}

\CommentTok{\# Define the threshold for low sequence numbers}
\NormalTok{threshold }\OtherTok{\textless{}{-}} \DecValTok{60000}

\CommentTok{\# Identify and remove samples with low sequence counts (below threshold)}
\NormalTok{samples\_to\_remove }\OtherTok{\textless{}{-}} \FunctionTok{names}\NormalTok{(}\FunctionTok{which}\NormalTok{(}\FunctionTok{colSums}\NormalTok{(bCOUNT) }\SpecialCharTok{\textless{}}\NormalTok{ threshold))}
\FunctionTok{print}\NormalTok{(}\FunctionTok{paste}\NormalTok{(}\FunctionTok{length}\NormalTok{(samples\_to\_remove), }\StringTok{"samples removed due to low number of sequences (threshold \textless{}"}\NormalTok{, threshold, }\StringTok{")"}\NormalTok{))}
\end{Highlighting}
\end{Shaded}

\begin{verbatim}
[1] "0 samples removed due to low number of sequences (threshold < 60000 )"
\end{verbatim}

\begin{Shaded}
\begin{Highlighting}[]
\CommentTok{\# Remove these samples from the bDESIGN and bCOUNT tables}
\NormalTok{bDESIGN }\OtherTok{\textless{}{-}}\NormalTok{ bDESIGN[}\SpecialCharTok{!}\NormalTok{(bDESIGN}\SpecialCharTok{$}\NormalTok{bc\_ID }\SpecialCharTok{\%in\%}\NormalTok{ samples\_to\_remove), ]}
\NormalTok{bCOUNT }\OtherTok{\textless{}{-}}\NormalTok{ bCOUNT[, }\SpecialCharTok{!}\NormalTok{(}\FunctionTok{colnames}\NormalTok{(bCOUNT) }\SpecialCharTok{\%in\%}\NormalTok{ samples\_to\_remove)]}

\CommentTok{\# Remove ASVs with zero counts across all samples}
\NormalTok{ASVs\_to\_remove }\OtherTok{\textless{}{-}} \FunctionTok{names}\NormalTok{(}\FunctionTok{which}\NormalTok{(}\FunctionTok{rowSums}\NormalTok{(bCOUNT) }\SpecialCharTok{==} \DecValTok{0}\NormalTok{))}
\NormalTok{bCOUNT }\OtherTok{\textless{}{-}}\NormalTok{ bCOUNT[}\SpecialCharTok{!}\NormalTok{(}\FunctionTok{rownames}\NormalTok{(bCOUNT) }\SpecialCharTok{\%in\%}\NormalTok{ ASVs\_to\_remove), ]}
\NormalTok{bTAXA }\OtherTok{\textless{}{-}}\NormalTok{ bTAXA[}\SpecialCharTok{!}\NormalTok{(}\FunctionTok{rownames}\NormalTok{(bTAXA) }\SpecialCharTok{\%in\%}\NormalTok{ ASVs\_to\_remove), ]}

\CommentTok{\# Summarize the remaining sequence data}
\NormalTok{df}\SpecialCharTok{$}\NormalTok{sum }\OtherTok{\textless{}{-}} \FunctionTok{sum}\NormalTok{(}\FunctionTok{colSums}\NormalTok{(bCOUNT))}
\NormalTok{df}\SpecialCharTok{$}\NormalTok{min }\OtherTok{\textless{}{-}} \FunctionTok{min}\NormalTok{(}\FunctionTok{colSums}\NormalTok{(bCOUNT))}
\NormalTok{df}\SpecialCharTok{$}\NormalTok{max }\OtherTok{\textless{}{-}} \FunctionTok{max}\NormalTok{(}\FunctionTok{colSums}\NormalTok{(bCOUNT))}
\NormalTok{df}\SpecialCharTok{$}\NormalTok{median }\OtherTok{\textless{}{-}} \FunctionTok{median}\NormalTok{(}\FunctionTok{colSums}\NormalTok{(bCOUNT))}
\NormalTok{df}\SpecialCharTok{$}\NormalTok{OTUs }\OtherTok{\textless{}{-}} \FunctionTok{nrow}\NormalTok{(bCOUNT)}

\CommentTok{\# Output the summary}
\CommentTok{\#pander(df)}
\FunctionTok{kable}\NormalTok{(df, }\AttributeTok{caption =} \StringTok{"Read count: Bacteria"}\NormalTok{,}
      \AttributeTok{col.names =} \FunctionTok{c}\NormalTok{(}\StringTok{"Total"}\NormalTok{, }\StringTok{"Min"}\NormalTok{, }\StringTok{"Max"}\NormalTok{, }\StringTok{"Median"}\NormalTok{, }\StringTok{"OTUs"}\NormalTok{))}
\end{Highlighting}
\end{Shaded}

\begin{longtable}[]{@{}rrrrr@{}}
\caption{Read count: Bacteria}\tabularnewline
\toprule\noalign{}
Total & Min & Max & Median & OTUs \\
\midrule\noalign{}
\endfirsthead
\toprule\noalign{}
Total & Min & Max & Median & OTUs \\
\midrule\noalign{}
\endhead
\bottomrule\noalign{}
\endlastfoot
5633798 & 60228 & 142079 & 95708.5 & 3665 \\
\end{longtable}

\begin{Shaded}
\begin{Highlighting}[]
\CommentTok{\# Clean up}
\FunctionTok{rm}\NormalTok{(df, ASVs\_to\_remove, samples\_to\_remove, threshold)}
\end{Highlighting}
\end{Shaded}

\section{Fungi sequence filtering}\label{fungi-sequence-filtering}

Purpose: Filtering: Removes samples with fewer than 800 sequences and
ASVs with zero sequences. Summary: Similar to bacteria, provides the
sequence statistics for fungi.

\begin{Shaded}
\begin{Highlighting}[]
\CommentTok{\# Initialize a dataframe to store summary statistics}
\NormalTok{df }\OtherTok{\textless{}{-}} \FunctionTok{data.frame}\NormalTok{(}\AttributeTok{sum =} \ConstantTok{NA}\NormalTok{, }\AttributeTok{min =} \ConstantTok{NA}\NormalTok{, }\AttributeTok{max =} \ConstantTok{NA}\NormalTok{, }\AttributeTok{median =} \ConstantTok{NA}\NormalTok{, }\AttributeTok{OTUs =} \ConstantTok{NA}\NormalTok{)}

\CommentTok{\# Define the threshold for low sequence numbers}
\NormalTok{threshold }\OtherTok{\textless{}{-}} \DecValTok{800}

\CommentTok{\# Identify and remove samples with low sequence counts (below threshold)}
\NormalTok{samples\_to\_remove }\OtherTok{\textless{}{-}} \FunctionTok{names}\NormalTok{(}\FunctionTok{which}\NormalTok{(}\FunctionTok{colSums}\NormalTok{(fCOUNT) }\SpecialCharTok{\textless{}}\NormalTok{ threshold))}
\FunctionTok{print}\NormalTok{(}\FunctionTok{paste}\NormalTok{(}\FunctionTok{length}\NormalTok{(samples\_to\_remove), }\StringTok{"samples removed due to low number of sequences (threshold \textless{}"}\NormalTok{, threshold, }\StringTok{")"}\NormalTok{))}
\end{Highlighting}
\end{Shaded}

\begin{verbatim}
[1] "9 samples removed due to low number of sequences (threshold < 800 )"
\end{verbatim}

\begin{Shaded}
\begin{Highlighting}[]
\CommentTok{\# Save the total sequence depth for plotting}
\NormalTok{depth\_fun\_all }\OtherTok{\textless{}{-}} \FunctionTok{colSums}\NormalTok{(fCOUNT)}

\CommentTok{\# Remove these samples from the fDESIGN and fCOUNT tables}
\NormalTok{fDESIGN }\OtherTok{\textless{}{-}}\NormalTok{ fDESIGN[}\SpecialCharTok{!}\NormalTok{(fDESIGN}\SpecialCharTok{$}\NormalTok{bc\_ID }\SpecialCharTok{\%in\%}\NormalTok{ samples\_to\_remove), ]}
\NormalTok{fCOUNT }\OtherTok{\textless{}{-}}\NormalTok{ fCOUNT[, }\SpecialCharTok{!}\NormalTok{(}\FunctionTok{colnames}\NormalTok{(fCOUNT) }\SpecialCharTok{\%in\%}\NormalTok{ samples\_to\_remove)]}

\CommentTok{\# Remove ASVs with zero counts across all samples}
\NormalTok{ASVs\_to\_remove }\OtherTok{\textless{}{-}} \FunctionTok{names}\NormalTok{(}\FunctionTok{which}\NormalTok{(}\FunctionTok{rowSums}\NormalTok{(fCOUNT) }\SpecialCharTok{==} \DecValTok{0}\NormalTok{))}
\NormalTok{fCOUNT }\OtherTok{\textless{}{-}}\NormalTok{ fCOUNT[}\SpecialCharTok{!}\NormalTok{(}\FunctionTok{rownames}\NormalTok{(fCOUNT) }\SpecialCharTok{\%in\%}\NormalTok{ ASVs\_to\_remove), ]}
\NormalTok{fTAXA }\OtherTok{\textless{}{-}}\NormalTok{ fTAXA[}\SpecialCharTok{!}\NormalTok{(}\FunctionTok{rownames}\NormalTok{(fTAXA) }\SpecialCharTok{\%in\%}\NormalTok{ ASVs\_to\_remove), ]}

\CommentTok{\# Summarize the remaining sequence data}
\NormalTok{df}\SpecialCharTok{$}\NormalTok{sum }\OtherTok{\textless{}{-}} \FunctionTok{sum}\NormalTok{(}\FunctionTok{colSums}\NormalTok{(fCOUNT))}
\NormalTok{df}\SpecialCharTok{$}\NormalTok{min }\OtherTok{\textless{}{-}} \FunctionTok{min}\NormalTok{(}\FunctionTok{colSums}\NormalTok{(fCOUNT))}
\NormalTok{df}\SpecialCharTok{$}\NormalTok{max }\OtherTok{\textless{}{-}} \FunctionTok{max}\NormalTok{(}\FunctionTok{colSums}\NormalTok{(fCOUNT))}
\NormalTok{df}\SpecialCharTok{$}\NormalTok{median }\OtherTok{\textless{}{-}} \FunctionTok{median}\NormalTok{(}\FunctionTok{colSums}\NormalTok{(fCOUNT))}
\NormalTok{df}\SpecialCharTok{$}\NormalTok{OTUs }\OtherTok{\textless{}{-}} \FunctionTok{nrow}\NormalTok{(fCOUNT)}

\CommentTok{\# Output the summary}
\CommentTok{\#pander(df)}
\FunctionTok{kable}\NormalTok{(df, }\AttributeTok{caption =} \StringTok{"Read count: Fungi"}\NormalTok{,}
      \AttributeTok{col.names =} \FunctionTok{c}\NormalTok{(}\StringTok{"Total"}\NormalTok{, }\StringTok{"Min"}\NormalTok{, }\StringTok{"Max"}\NormalTok{, }\StringTok{"Median"}\NormalTok{, }\StringTok{"OTUs"}\NormalTok{))}
\end{Highlighting}
\end{Shaded}

\begin{longtable}[]{@{}rrrrr@{}}
\caption{Read count: Fungi}\tabularnewline
\toprule\noalign{}
Total & Min & Max & Median & OTUs \\
\midrule\noalign{}
\endfirsthead
\toprule\noalign{}
Total & Min & Max & Median & OTUs \\
\midrule\noalign{}
\endhead
\bottomrule\noalign{}
\endlastfoot
134267 & 805 & 7870 & 1920 & 223 \\
\end{longtable}

\begin{Shaded}
\begin{Highlighting}[]
\CommentTok{\# Clean up}
\FunctionTok{rm}\NormalTok{(df, ASVs\_to\_remove, samples\_to\_remove, threshold)}
\end{Highlighting}
\end{Shaded}

\section{Oomycetes sequence
filtering}\label{oomycetes-sequence-filtering}

Purpose: Filtering: Removes samples with fewer than 5,700 sequences and
ASVs with zero sequences. Summary: Outputs statistics on the oomycetes
sequence data.

\begin{Shaded}
\begin{Highlighting}[]
\CommentTok{\# Initialize a dataframe to store summary statistics}
\NormalTok{df }\OtherTok{\textless{}{-}} \FunctionTok{data.frame}\NormalTok{(}\AttributeTok{sum =} \ConstantTok{NA}\NormalTok{, }\AttributeTok{min =} \ConstantTok{NA}\NormalTok{, }\AttributeTok{max =} \ConstantTok{NA}\NormalTok{, }\AttributeTok{median =} \ConstantTok{NA}\NormalTok{, }\AttributeTok{OTUs =} \ConstantTok{NA}\NormalTok{)}

\CommentTok{\# Define the threshold for low sequence numbers}
\NormalTok{threshold }\OtherTok{\textless{}{-}} \DecValTok{5700}

\CommentTok{\# Identify and remove samples with low sequence counts (below threshold)}
\NormalTok{samples\_to\_remove }\OtherTok{\textless{}{-}} \FunctionTok{names}\NormalTok{(}\FunctionTok{which}\NormalTok{(}\FunctionTok{colSums}\NormalTok{(oCOUNT) }\SpecialCharTok{\textless{}}\NormalTok{ threshold))}
\FunctionTok{print}\NormalTok{(}\FunctionTok{paste}\NormalTok{(}\FunctionTok{length}\NormalTok{(samples\_to\_remove), }\StringTok{"samples removed due to low number of sequences (threshold \textless{}"}\NormalTok{, threshold, }\StringTok{")"}\NormalTok{))}
\end{Highlighting}
\end{Shaded}

\begin{verbatim}
[1] "0 samples removed due to low number of sequences (threshold < 5700 )"
\end{verbatim}

\begin{Shaded}
\begin{Highlighting}[]
\CommentTok{\# Save the total sequence depth for plotting}
\NormalTok{depth\_oom\_all }\OtherTok{\textless{}{-}} \FunctionTok{colSums}\NormalTok{(oCOUNT)}

\CommentTok{\# Remove these samples from the oDESIGN and oCOUNT tables}
\NormalTok{oDESIGN }\OtherTok{\textless{}{-}}\NormalTok{ oDESIGN[}\SpecialCharTok{!}\NormalTok{(oDESIGN}\SpecialCharTok{$}\NormalTok{bc\_ID }\SpecialCharTok{\%in\%}\NormalTok{ samples\_to\_remove), ]}
\NormalTok{oCOUNT }\OtherTok{\textless{}{-}}\NormalTok{ oCOUNT[, }\SpecialCharTok{!}\NormalTok{(}\FunctionTok{colnames}\NormalTok{(oCOUNT) }\SpecialCharTok{\%in\%}\NormalTok{ samples\_to\_remove)]}

\CommentTok{\# Remove ASVs with zero counts across all samples}
\NormalTok{ASVs\_to\_remove }\OtherTok{\textless{}{-}} \FunctionTok{names}\NormalTok{(}\FunctionTok{which}\NormalTok{(}\FunctionTok{rowSums}\NormalTok{(oCOUNT) }\SpecialCharTok{==} \DecValTok{0}\NormalTok{))}
\NormalTok{oCOUNT }\OtherTok{\textless{}{-}}\NormalTok{ oCOUNT[}\SpecialCharTok{!}\NormalTok{(}\FunctionTok{rownames}\NormalTok{(oCOUNT) }\SpecialCharTok{\%in\%}\NormalTok{ ASVs\_to\_remove), ]}
\NormalTok{oTAXA }\OtherTok{\textless{}{-}}\NormalTok{ oTAXA[}\SpecialCharTok{!}\NormalTok{(}\FunctionTok{rownames}\NormalTok{(oTAXA) }\SpecialCharTok{\%in\%}\NormalTok{ ASVs\_to\_remove), ]}

\CommentTok{\# Summarize the remaining sequence data}
\NormalTok{df}\SpecialCharTok{$}\NormalTok{sum }\OtherTok{\textless{}{-}} \FunctionTok{sum}\NormalTok{(}\FunctionTok{colSums}\NormalTok{(oCOUNT))}
\NormalTok{df}\SpecialCharTok{$}\NormalTok{min }\OtherTok{\textless{}{-}} \FunctionTok{min}\NormalTok{(}\FunctionTok{colSums}\NormalTok{(oCOUNT))}
\NormalTok{df}\SpecialCharTok{$}\NormalTok{max }\OtherTok{\textless{}{-}} \FunctionTok{max}\NormalTok{(}\FunctionTok{colSums}\NormalTok{(oCOUNT))}
\NormalTok{df}\SpecialCharTok{$}\NormalTok{median }\OtherTok{\textless{}{-}} \FunctionTok{median}\NormalTok{(}\FunctionTok{colSums}\NormalTok{(oCOUNT))}
\NormalTok{df}\SpecialCharTok{$}\NormalTok{OTUs }\OtherTok{\textless{}{-}} \FunctionTok{nrow}\NormalTok{(oCOUNT)}

\CommentTok{\# Output the summary}
\CommentTok{\#pander(df)}
\FunctionTok{kable}\NormalTok{(df, }\AttributeTok{caption =} \StringTok{"Read count: Oomycetes"}\NormalTok{,}
      \AttributeTok{col.names =} \FunctionTok{c}\NormalTok{(}\StringTok{"Total"}\NormalTok{, }\StringTok{"Min"}\NormalTok{, }\StringTok{"Max"}\NormalTok{, }\StringTok{"Median"}\NormalTok{, }\StringTok{"OTUs"}\NormalTok{))}
\end{Highlighting}
\end{Shaded}

\begin{longtable}[]{@{}rrrrr@{}}
\caption{Read count: Oomycetes}\tabularnewline
\toprule\noalign{}
Total & Min & Max & Median & OTUs \\
\midrule\noalign{}
\endfirsthead
\toprule\noalign{}
Total & Min & Max & Median & OTUs \\
\midrule\noalign{}
\endhead
\bottomrule\noalign{}
\endlastfoot
646360 & 5750 & 16490 & 10536 & 90 \\
\end{longtable}

\begin{Shaded}
\begin{Highlighting}[]
\CommentTok{\# Clean up}
\FunctionTok{rm}\NormalTok{(df, ASVs\_to\_remove, samples\_to\_remove, threshold)}
\end{Highlighting}
\end{Shaded}

\section{Remove Eukaryotic
Contaminants}\label{remove-eukaryotic-contaminants}

\subsection{Bacteria}\label{bacteria}

Purpose: Sequence Mapping: Compares bacterial sequences to host
sequences to detect and remove contaminant OTUs (eukaryotic).

\begin{Shaded}
\begin{Highlighting}[]
\CommentTok{\# Import At.GI amplicon sequences (host sequences)}
\NormalTok{AtSEQ }\OtherTok{\textless{}{-}}\NormalTok{ Biostrings}\SpecialCharTok{::}\FunctionTok{readDNAStringSet}\NormalTok{(paths}\SpecialCharTok{$}\NormalTok{host\_amplicon)}
\NormalTok{AtSEQ }\OtherTok{\textless{}{-}} \FunctionTok{as.data.frame}\NormalTok{(AtSEQ)}
\FunctionTok{colnames}\NormalTok{(AtSEQ) }\OtherTok{\textless{}{-}} \StringTok{"seqs"}

\CommentTok{\# Import bacterial sequences}
\NormalTok{bSEQ }\OtherTok{\textless{}{-}}\NormalTok{ Biostrings}\SpecialCharTok{::}\FunctionTok{readDNAStringSet}\NormalTok{(paths}\SpecialCharTok{$}\NormalTok{bseq)}
\NormalTok{bSEQ }\OtherTok{\textless{}{-}} \FunctionTok{as.data.frame}\NormalTok{(bSEQ)}
\FunctionTok{colnames}\NormalTok{(bSEQ) }\OtherTok{\textless{}{-}} \StringTok{"seqs"}
\FunctionTok{rownames}\NormalTok{(bSEQ) }\OtherTok{\textless{}{-}} \FunctionTok{gsub}\NormalTok{(}\StringTok{"OTU"}\NormalTok{, }\StringTok{"bOTU"}\NormalTok{, }\FunctionTok{rownames}\NormalTok{(bSEQ))}

\CommentTok{\# Identify eukaryotic OTUs in bacterial data}
\NormalTok{eOTU }\OtherTok{\textless{}{-}} \FunctionTok{rownames}\NormalTok{(bTAXA)[bTAXA}\SpecialCharTok{$}\NormalTok{labels }\SpecialCharTok{==} \StringTok{"Host"}\NormalTok{]}
\NormalTok{eSEQ }\OtherTok{\textless{}{-}}\NormalTok{ bSEQ[}\FunctionTok{rownames}\NormalTok{(bSEQ) }\SpecialCharTok{\%in\%}\NormalTok{ eOTU, , drop }\OtherTok{=} \ConstantTok{FALSE}\NormalTok{]}

\CommentTok{\# Initialize a distance matrix for mapping sequences}
\NormalTok{distances }\OtherTok{\textless{}{-}} \FunctionTok{data.frame}\NormalTok{(}\FunctionTok{matrix}\NormalTok{(}\ConstantTok{NA}\NormalTok{, }\AttributeTok{nrow =} \FunctionTok{nrow}\NormalTok{(AtSEQ), }\AttributeTok{ncol =} \FunctionTok{nrow}\NormalTok{(eSEQ)), }\AttributeTok{row.names =} \FunctionTok{rownames}\NormalTok{(AtSEQ))}
\FunctionTok{colnames}\NormalTok{(distances) }\OtherTok{\textless{}{-}} \FunctionTok{rownames}\NormalTok{(eSEQ)}

\CommentTok{\# Function to map sequences and calculate distances}
\NormalTok{map }\OtherTok{\textless{}{-}} \ControlFlowTok{function}\NormalTok{(AtSEQ) \{}
\NormalTok{  eSEQs\_to\_map }\OtherTok{\textless{}{-}} \FunctionTok{data.frame}\NormalTok{(}\AttributeTok{seqs =} \FunctionTok{c}\NormalTok{(}\FunctionTok{as.character}\NormalTok{(AtSEQ), }\FunctionTok{as.character}\NormalTok{(eSEQ}\SpecialCharTok{$}\NormalTok{seqs[i])))}
\NormalTok{  eSEQs\_to\_map }\OtherTok{\textless{}{-}}\NormalTok{ Biostrings}\SpecialCharTok{::}\FunctionTok{DNAStringSet}\NormalTok{(eSEQs\_to\_map}\SpecialCharTok{$}\NormalTok{seqs)}
\NormalTok{  eSEQs\_aln }\OtherTok{\textless{}{-}}\NormalTok{ DECIPHER}\SpecialCharTok{::}\FunctionTok{AlignSeqs}\NormalTok{(eSEQs\_to\_map, }\AttributeTok{processors =} \DecValTok{1}\NormalTok{, }\AttributeTok{verbose =} \ConstantTok{FALSE}\NormalTok{)}
\NormalTok{  dist\_all }\OtherTok{\textless{}{-}}\NormalTok{ DECIPHER}\SpecialCharTok{::}\FunctionTok{DistanceMatrix}\NormalTok{(eSEQs\_aln, }\AttributeTok{processors =} \DecValTok{1}\NormalTok{, }\AttributeTok{verbose =} \ConstantTok{FALSE}\NormalTok{)}
  \FunctionTok{return}\NormalTok{(dist\_all[}\DecValTok{1}\NormalTok{, }\DecValTok{2}\NormalTok{])}
\NormalTok{\}}

\CommentTok{\# Calculate distances for all eukaryotic sequences}
\ControlFlowTok{for}\NormalTok{ (i }\ControlFlowTok{in} \DecValTok{1}\SpecialCharTok{:}\FunctionTok{nrow}\NormalTok{(eSEQ)) \{}
\NormalTok{  distances[, i] }\OtherTok{\textless{}{-}} \FunctionTok{apply}\NormalTok{(AtSEQ, }\DecValTok{1}\NormalTok{, map)}
\NormalTok{\}}

\CommentTok{\# Calculate sequence similarities}
\NormalTok{similarities }\OtherTok{\textless{}{-}} \DecValTok{1} \SpecialCharTok{{-}}\NormalTok{ distances}

\CommentTok{\# Add OTU abundance to the similarity matrix and output}
\NormalTok{similarities }\OtherTok{\textless{}{-}} \FunctionTok{rbind}\NormalTok{(similarities, }\AttributeTok{n\_seqs =} \FunctionTok{rowSums}\NormalTok{(bCOUNT[}\FunctionTok{rownames}\NormalTok{(bCOUNT) }\SpecialCharTok{\%in\%} \FunctionTok{colnames}\NormalTok{(distances), ]))}
\FunctionTok{pander}\NormalTok{(}\FunctionTok{t}\NormalTok{(similarities))}
\end{Highlighting}
\end{Shaded}

\begin{longtable}[]{@{}
  >{\centering\arraybackslash}p{(\linewidth - 4\tabcolsep) * \real{0.1944}}
  >{\centering\arraybackslash}p{(\linewidth - 4\tabcolsep) * \real{0.2222}}
  >{\centering\arraybackslash}p{(\linewidth - 4\tabcolsep) * \real{0.1250}}@{}}
\toprule\noalign{}
\begin{minipage}[b]{\linewidth}\centering
~
\end{minipage} & \begin{minipage}[b]{\linewidth}\centering
A.thaliana\_GI
\end{minipage} & \begin{minipage}[b]{\linewidth}\centering
n\_seqs
\end{minipage} \\
\midrule\noalign{}
\endhead
\bottomrule\noalign{}
\endlastfoot
\textbf{bOTU1} & 1 & 971818 \\
\textbf{bOTU27} & 0.9966 & 24852 \\
\textbf{bOTU188} & 0.9966 & 2314 \\
\textbf{bOTU191} & 0.9966 & 2554 \\
\end{longtable}

\begin{Shaded}
\begin{Highlighting}[]
\FunctionTok{kable}\NormalTok{(}\FunctionTok{t}\NormalTok{(similarities), }\AttributeTok{caption =} \StringTok{"OTUs mapping to GIGANTEA"}\NormalTok{,}
      \AttributeTok{col.names =} \FunctionTok{c}\NormalTok{(}\StringTok{"OTUs"}\NormalTok{, }\StringTok{"Similarity"}\NormalTok{, }\StringTok{"n\_seqs"}\NormalTok{))}
\end{Highlighting}
\end{Shaded}

\begin{longtable}[]{@{}lrr@{}}
\caption{OTUs mapping to GIGANTEA}\tabularnewline
\toprule\noalign{}
OTUs & Similarity & n\_seqs \\
\midrule\noalign{}
\endfirsthead
\toprule\noalign{}
OTUs & Similarity & n\_seqs \\
\midrule\noalign{}
\endhead
\bottomrule\noalign{}
\endlastfoot
bOTU1 & 1.0000000 & 971818 \\
bOTU27 & 0.9966216 & 24852 \\
bOTU188 & 0.9966216 & 2314 \\
bOTU191 & 0.9966216 & 2554 \\
\end{longtable}

\begin{Shaded}
\begin{Highlighting}[]
\CommentTok{\# Remove OTUs with similarity below 99\%}
\NormalTok{OTUs\_to\_remove }\OtherTok{\textless{}{-}} \FunctionTok{colnames}\NormalTok{(similarities)[similarities[}\DecValTok{1}\NormalTok{, ] }\SpecialCharTok{\textless{}} \FloatTok{0.99}\NormalTok{]}
\NormalTok{bCOUNT }\OtherTok{\textless{}{-}}\NormalTok{ bCOUNT[}\SpecialCharTok{!}\NormalTok{(}\FunctionTok{rownames}\NormalTok{(bCOUNT) }\SpecialCharTok{\%in\%}\NormalTok{ OTUs\_to\_remove), ]}
\NormalTok{bTAXA }\OtherTok{\textless{}{-}}\NormalTok{ bTAXA[}\SpecialCharTok{!}\NormalTok{(}\FunctionTok{rownames}\NormalTok{(bTAXA) }\SpecialCharTok{\%in\%}\NormalTok{ OTUs\_to\_remove), ]}

\CommentTok{\# Remove OTUs with zero abundance after filtering}
\NormalTok{zeroOTUs\_to\_remove }\OtherTok{\textless{}{-}} \FunctionTok{names}\NormalTok{(}\FunctionTok{which}\NormalTok{(}\FunctionTok{colSums}\NormalTok{(bCOUNT) }\SpecialCharTok{==} \DecValTok{0}\NormalTok{))}
\NormalTok{bCOUNT }\OtherTok{\textless{}{-}}\NormalTok{ bCOUNT[, }\SpecialCharTok{!}\NormalTok{(}\FunctionTok{colnames}\NormalTok{(bCOUNT) }\SpecialCharTok{\%in\%}\NormalTok{ zeroOTUs\_to\_remove)]}
\NormalTok{bTAXA }\OtherTok{\textless{}{-}}\NormalTok{ bTAXA[}\SpecialCharTok{!}\NormalTok{(}\FunctionTok{rownames}\NormalTok{(bTAXA) }\SpecialCharTok{\%in\%}\NormalTok{ zeroOTUs\_to\_remove), ]}

\CommentTok{\# Report removed OTUs}
\ControlFlowTok{if}\NormalTok{ (}\FunctionTok{length}\NormalTok{(OTUs\_to\_remove) }\SpecialCharTok{\textgreater{}} \DecValTok{0}\NormalTok{) \{}
  \FunctionTok{print}\NormalTok{(}\FunctionTok{paste0}\NormalTok{(}\StringTok{"The following eukaryotic OTUs were removed: "}\NormalTok{, }\FunctionTok{paste}\NormalTok{(OTUs\_to\_remove, }\AttributeTok{collapse =} \StringTok{", "}\NormalTok{)))}
\NormalTok{\} }\ControlFlowTok{else}\NormalTok{ \{}
  \FunctionTok{print}\NormalTok{(}\StringTok{"No eukaryotic OTUs were removed"}\NormalTok{)}
\NormalTok{\}}
\end{Highlighting}
\end{Shaded}

\begin{verbatim}
[1] "No eukaryotic OTUs were removed"
\end{verbatim}

\begin{Shaded}
\begin{Highlighting}[]
\CommentTok{\# Clean up}
\FunctionTok{rm}\NormalTok{(eSEQ, AtSEQ, similarities, OTUs\_to\_remove, zeroOTUs\_to\_remove, distances)}
\end{Highlighting}
\end{Shaded}

\subsection{Fungi}\label{fungi}

Purpose:

Importing Sequences: Loads host and fungal sequences from specified file
paths. Filtering: Retains only those fungal sequences present in the
taxonomy data (fTAXA). Identifying Contaminants: Detects eukaryotic OTUs
(considered as potential contaminants) within the fungal dataset.
Mapping and Similarity Calculation: Aligns each host sequence with
fungal sequences to calculate pairwise distances and convert them into
similarity scores. OTU Removal: Eliminates fungal OTUs that have less
than 99\% similarity to host sequences, effectively removing potential
eukaryotic contaminants. Cleaning Up: Removes any OTUs with zero
abundance post-filtering to ensure data integrity.

\begin{Shaded}
\begin{Highlighting}[]
\CommentTok{\# Import Host Amplicon Sequences}
\NormalTok{AtSEQ }\OtherTok{\textless{}{-}}\NormalTok{ Biostrings}\SpecialCharTok{::}\FunctionTok{readDNAStringSet}\NormalTok{(paths}\SpecialCharTok{$}\NormalTok{host\_amplicon)}
\NormalTok{AtSEQ }\OtherTok{\textless{}{-}} \FunctionTok{as.data.frame}\NormalTok{(AtSEQ)}
\FunctionTok{colnames}\NormalTok{(AtSEQ) }\OtherTok{\textless{}{-}} \StringTok{"seqs"}

\CommentTok{\# Import Fungal Sequences}
\NormalTok{fSEQ }\OtherTok{\textless{}{-}}\NormalTok{ Biostrings}\SpecialCharTok{::}\FunctionTok{readDNAStringSet}\NormalTok{(paths}\SpecialCharTok{$}\NormalTok{fseq)}
\NormalTok{fSEQ }\OtherTok{\textless{}{-}} \FunctionTok{as.data.frame}\NormalTok{(fSEQ)}
\FunctionTok{colnames}\NormalTok{(fSEQ) }\OtherTok{\textless{}{-}} \StringTok{"seqs"}
\FunctionTok{rownames}\NormalTok{(fSEQ) }\OtherTok{\textless{}{-}} \FunctionTok{gsub}\NormalTok{(}\StringTok{"OTU"}\NormalTok{, }\StringTok{"fOTU"}\NormalTok{, }\FunctionTok{rownames}\NormalTok{(fSEQ))  }\CommentTok{\# Prefix OTU IDs with \textquotesingle{}f\textquotesingle{} for Fungi}

\CommentTok{\# Filter to Keep Only Fungal Sequences Present in Taxonomy Data}
\NormalTok{fSEQ }\OtherTok{\textless{}{-}}\NormalTok{ fSEQ[}\FunctionTok{rownames}\NormalTok{(fSEQ) }\SpecialCharTok{\%in\%} \FunctionTok{rownames}\NormalTok{(fTAXA), , drop }\OtherTok{=} \ConstantTok{FALSE}\NormalTok{]}

\CommentTok{\# Identify Eukaryotic OTUs (Hosts) in Fungal Data}
\NormalTok{eOTU }\OtherTok{\textless{}{-}} \FunctionTok{rownames}\NormalTok{(fTAXA)[fTAXA}\SpecialCharTok{$}\NormalTok{labels }\SpecialCharTok{==} \StringTok{"Host"}\NormalTok{]}

\CommentTok{\# Extract Eukaryotic Sequences}
\NormalTok{eSEQ }\OtherTok{\textless{}{-}}\NormalTok{ fSEQ[}\FunctionTok{rownames}\NormalTok{(fSEQ) }\SpecialCharTok{\%in\%}\NormalTok{ eOTU, , drop }\OtherTok{=} \ConstantTok{FALSE}\NormalTok{]}

\CommentTok{\# Clean Up Unnecessary Variables}
\FunctionTok{rm}\NormalTok{(fSEQ, eOTU)}

\DocumentationTok{\#\#\# MAPPING }\AlertTok{\#\#\#}

\CommentTok{\# Initialize a Data Frame to Store Pairwise Distances}
\NormalTok{distances }\OtherTok{\textless{}{-}} \FunctionTok{data.frame}\NormalTok{(}\FunctionTok{matrix}\NormalTok{(}\ConstantTok{NA}\NormalTok{, }\AttributeTok{nrow =} \FunctionTok{nrow}\NormalTok{(AtSEQ), }\AttributeTok{ncol =} \FunctionTok{nrow}\NormalTok{(eSEQ)), }\AttributeTok{row.names =} \FunctionTok{rownames}\NormalTok{(AtSEQ))}
\FunctionTok{colnames}\NormalTok{(distances) }\OtherTok{\textless{}{-}} \FunctionTok{rownames}\NormalTok{(eSEQ)}

\CommentTok{\# Define a Mapping Function to Calculate Distance Between Host and Fungal Sequences}
\NormalTok{map }\OtherTok{\textless{}{-}} \ControlFlowTok{function}\NormalTok{(AtSEQ\_entry) \{}
  \CommentTok{\# Combine Host Sequence and Current Fungal Sequence}
\NormalTok{  eSEQs\_to\_map }\OtherTok{\textless{}{-}} \FunctionTok{data.frame}\NormalTok{(}\AttributeTok{seqs =} \FunctionTok{c}\NormalTok{(}\FunctionTok{as.character}\NormalTok{(AtSEQ\_entry), }\FunctionTok{as.character}\NormalTok{(eSEQ}\SpecialCharTok{$}\NormalTok{seqs[i])))}
  
  \CommentTok{\# Convert to DNAStringSet for Alignment}
\NormalTok{  eSEQs\_to\_map }\OtherTok{\textless{}{-}}\NormalTok{ Biostrings}\SpecialCharTok{::}\FunctionTok{DNAStringSet}\NormalTok{(eSEQs\_to\_map}\SpecialCharTok{$}\NormalTok{seqs)}
  
  \CommentTok{\# Align Sequences}
\NormalTok{  eSEQs\_aln }\OtherTok{\textless{}{-}}\NormalTok{ DECIPHER}\SpecialCharTok{::}\FunctionTok{AlignSeqs}\NormalTok{(eSEQs\_to\_map, }\AttributeTok{processors =} \DecValTok{1}\NormalTok{, }\AttributeTok{verbose =} \ConstantTok{FALSE}\NormalTok{)}
  
  \CommentTok{\# Compute Distance Matrix}
\NormalTok{  dist\_all }\OtherTok{\textless{}{-}}\NormalTok{ DECIPHER}\SpecialCharTok{::}\FunctionTok{DistanceMatrix}\NormalTok{(eSEQs\_aln, }\AttributeTok{processors =} \DecValTok{1}\NormalTok{, }\AttributeTok{verbose =} \ConstantTok{FALSE}\NormalTok{)}
  
  \CommentTok{\# Return the Distance Between Host and Fungal Sequence}
  \FunctionTok{return}\NormalTok{(dist\_all[}\DecValTok{1}\NormalTok{, }\DecValTok{2}\NormalTok{])}
\NormalTok{\}}

\CommentTok{\# Calculate Distances for All Eukaryotic Sequences}
\ControlFlowTok{for}\NormalTok{ (i }\ControlFlowTok{in} \DecValTok{1}\SpecialCharTok{:}\FunctionTok{nrow}\NormalTok{(eSEQ)) \{}
\NormalTok{  distances[, i] }\OtherTok{\textless{}{-}} \FunctionTok{apply}\NormalTok{(AtSEQ, }\DecValTok{1}\NormalTok{, map)}
\NormalTok{\}}

\CommentTok{\# Compute Similarity Scores}
\NormalTok{similarities }\OtherTok{\textless{}{-}} \DecValTok{1} \SpecialCharTok{{-}}\NormalTok{ distances}

\CommentTok{\# Append OTU Abundance to Similarity Matrix}
\NormalTok{similarities }\OtherTok{\textless{}{-}} \FunctionTok{rbind}\NormalTok{(similarities, }\AttributeTok{n\_seqs =} \FunctionTok{rowSums}\NormalTok{(fCOUNT[}\FunctionTok{rownames}\NormalTok{(fCOUNT) }\SpecialCharTok{\%in\%} \FunctionTok{colnames}\NormalTok{(distances), ]))}

\CommentTok{\# Display the Similarity Matrix}
\CommentTok{\#pander(t(similarities))}
\FunctionTok{kable}\NormalTok{(}\FunctionTok{t}\NormalTok{(similarities), }\AttributeTok{caption =} \StringTok{"OTUs mapping to GIGANTEA"}\NormalTok{,}
      \AttributeTok{col.names =} \FunctionTok{c}\NormalTok{(}\StringTok{"OTUs"}\NormalTok{, }\StringTok{"Similarity"}\NormalTok{, }\StringTok{"n\_seqs"}\NormalTok{))}
\end{Highlighting}
\end{Shaded}

\begin{longtable}[]{@{}lrr@{}}
\caption{OTUs mapping to GIGANTEA}\tabularnewline
\toprule\noalign{}
OTUs & Similarity & n\_seqs \\
\midrule\noalign{}
\endfirsthead
\toprule\noalign{}
OTUs & Similarity & n\_seqs \\
\midrule\noalign{}
\endhead
\bottomrule\noalign{}
\endlastfoot
fOTU1 & 1 & 87461 \\
\end{longtable}

\begin{Shaded}
\begin{Highlighting}[]
\CommentTok{\# Identify and Remove OTUs with Similarity Below 99\%}
\NormalTok{OTUs\_to\_remove }\OtherTok{\textless{}{-}} \FunctionTok{colnames}\NormalTok{(similarities)[similarities[}\DecValTok{1}\NormalTok{, ] }\SpecialCharTok{\textless{}} \FloatTok{0.99}\NormalTok{]}
\NormalTok{fCOUNT }\OtherTok{\textless{}{-}}\NormalTok{ fCOUNT[}\SpecialCharTok{!}\NormalTok{(}\FunctionTok{rownames}\NormalTok{(fCOUNT) }\SpecialCharTok{\%in\%}\NormalTok{ OTUs\_to\_remove), ]}
\NormalTok{fTAXA }\OtherTok{\textless{}{-}}\NormalTok{ fTAXA[}\SpecialCharTok{!}\NormalTok{(}\FunctionTok{rownames}\NormalTok{(fTAXA) }\SpecialCharTok{\%in\%}\NormalTok{ OTUs\_to\_remove), ]}

\CommentTok{\# Further Remove OTUs with Zero Abundance After Filtering}
\NormalTok{zeroOTUs\_to\_remove }\OtherTok{\textless{}{-}} \FunctionTok{names}\NormalTok{(}\FunctionTok{which}\NormalTok{(}\FunctionTok{colSums}\NormalTok{(fCOUNT) }\SpecialCharTok{==} \DecValTok{0}\NormalTok{))}
\NormalTok{fCOUNT }\OtherTok{\textless{}{-}}\NormalTok{ fCOUNT[, }\SpecialCharTok{!}\NormalTok{(}\FunctionTok{colnames}\NormalTok{(fCOUNT) }\SpecialCharTok{\%in\%}\NormalTok{ zeroOTUs\_to\_remove)]}
\NormalTok{fTAXA }\OtherTok{\textless{}{-}}\NormalTok{ fTAXA[}\SpecialCharTok{!}\NormalTok{(}\FunctionTok{rownames}\NormalTok{(fTAXA) }\SpecialCharTok{\%in\%}\NormalTok{ zeroOTUs\_to\_remove), ]}

\CommentTok{\# Report Removed OTUs}
\ControlFlowTok{if}\NormalTok{ (}\FunctionTok{length}\NormalTok{(OTUs\_to\_remove) }\SpecialCharTok{\textgreater{}} \DecValTok{0}\NormalTok{) \{}
  \FunctionTok{print}\NormalTok{(}\FunctionTok{paste0}\NormalTok{(}\StringTok{"The following eukaryotic OTUs were removed: "}\NormalTok{, }\FunctionTok{paste}\NormalTok{(OTUs\_to\_remove, }\AttributeTok{collapse =} \StringTok{", "}\NormalTok{)))}
\NormalTok{\} }\ControlFlowTok{else}\NormalTok{ \{}
  \FunctionTok{print}\NormalTok{(}\StringTok{"No eukaryotic OTUs were removed"}\NormalTok{)}
\NormalTok{\}}
\end{Highlighting}
\end{Shaded}

\begin{verbatim}
[1] "No eukaryotic OTUs were removed"
\end{verbatim}

\begin{Shaded}
\begin{Highlighting}[]
\CommentTok{\# Clean Up Temporary Variables}
\FunctionTok{rm}\NormalTok{(eSEQ, AtSEQ, similarities, OTUs\_to\_remove, zeroOTUs\_to\_remove, distances)}
\end{Highlighting}
\end{Shaded}

\subsection{Oomycetes}\label{oomycetes}

Purpose:

Importing Sequences: Loads host sequences and reuses the fungal dataset
(paths\$fseq) for Oomycetes analysis. Filtering: Retains only those
oomycetal sequences present in the taxonomy data (oTAXA). Identifying
Contaminants: Detects eukaryotic OTUs (considered as potential
contaminants) within the oomycetal dataset. Mapping and Similarity
Calculation: Aligns each host sequence with oomycetal sequences to
calculate pairwise distances and convert them into similarity scores.
OTU Removal: Eliminates oomycetal OTUs that have less than 99\%
similarity to host sequences, effectively removing potential eukaryotic
contaminants. Cleaning Up: Removes any OTUs with zero abundance
post-filtering to ensure data integrity.

\begin{Shaded}
\begin{Highlighting}[]
\CommentTok{\# Import Host Amplicon Sequences}
\NormalTok{AtSEQ }\OtherTok{\textless{}{-}}\NormalTok{ Biostrings}\SpecialCharTok{::}\FunctionTok{readDNAStringSet}\NormalTok{(paths}\SpecialCharTok{$}\NormalTok{host\_amplicon)}
\NormalTok{AtSEQ }\OtherTok{\textless{}{-}} \FunctionTok{as.data.frame}\NormalTok{(AtSEQ)}
\FunctionTok{colnames}\NormalTok{(AtSEQ) }\OtherTok{\textless{}{-}} \StringTok{"seqs"}

\CommentTok{\# Import Oomycetal Sequences (Same Dataset as Fungi)}
\NormalTok{oSEQ }\OtherTok{\textless{}{-}}\NormalTok{ Biostrings}\SpecialCharTok{::}\FunctionTok{readDNAStringSet}\NormalTok{(paths}\SpecialCharTok{$}\NormalTok{fseq)  }\CommentTok{\# Reusing the fungal dataset for Oomycetes}
\NormalTok{oSEQ }\OtherTok{\textless{}{-}} \FunctionTok{as.data.frame}\NormalTok{(oSEQ)}
\FunctionTok{colnames}\NormalTok{(oSEQ) }\OtherTok{\textless{}{-}} \StringTok{"seqs"}
\FunctionTok{rownames}\NormalTok{(oSEQ) }\OtherTok{\textless{}{-}} \FunctionTok{gsub}\NormalTok{(}\StringTok{"OTU"}\NormalTok{, }\StringTok{"oOTU"}\NormalTok{, }\FunctionTok{rownames}\NormalTok{(oSEQ))  }\CommentTok{\# Prefix OTU IDs with \textquotesingle{}o\textquotesingle{} for Oomycetes}

\CommentTok{\# Filter to Keep Only Oomycetal Sequences Present in Taxonomy Data}
\NormalTok{oSEQ }\OtherTok{\textless{}{-}}\NormalTok{ oSEQ[}\FunctionTok{rownames}\NormalTok{(oSEQ) }\SpecialCharTok{\%in\%} \FunctionTok{rownames}\NormalTok{(oTAXA), , drop }\OtherTok{=} \ConstantTok{FALSE}\NormalTok{]}

\CommentTok{\# Identify Eukaryotic OTUs (Hosts) in Oomycetal Data}
\NormalTok{eOTU }\OtherTok{\textless{}{-}} \FunctionTok{rownames}\NormalTok{(oTAXA)[oTAXA}\SpecialCharTok{$}\NormalTok{labels }\SpecialCharTok{==} \StringTok{"Host"}\NormalTok{]}

\CommentTok{\# Extract Eukaryotic Sequences}
\NormalTok{eSEQ }\OtherTok{\textless{}{-}}\NormalTok{ oSEQ[}\FunctionTok{rownames}\NormalTok{(oSEQ) }\SpecialCharTok{\%in\%}\NormalTok{ eOTU, , drop }\OtherTok{=} \ConstantTok{FALSE}\NormalTok{]}

\CommentTok{\# Clean Up Unnecessary Variables}
\FunctionTok{rm}\NormalTok{(oSEQ, eOTU)}

\DocumentationTok{\#\#\# MAPPING }\AlertTok{\#\#\#}

\CommentTok{\# Initialize a Data Frame to Store Pairwise Distances}
\NormalTok{distances }\OtherTok{\textless{}{-}} \FunctionTok{data.frame}\NormalTok{(}\FunctionTok{matrix}\NormalTok{(}\ConstantTok{NA}\NormalTok{, }\AttributeTok{nrow =} \FunctionTok{nrow}\NormalTok{(AtSEQ), }\AttributeTok{ncol =} \FunctionTok{nrow}\NormalTok{(eSEQ)), }\AttributeTok{row.names =} \FunctionTok{rownames}\NormalTok{(AtSEQ))}
\FunctionTok{colnames}\NormalTok{(distances) }\OtherTok{\textless{}{-}} \FunctionTok{rownames}\NormalTok{(eSEQ)}

\CommentTok{\# Define a Mapping Function to Calculate Distance Between Host and Oomycetal Sequences}
\NormalTok{map }\OtherTok{\textless{}{-}} \ControlFlowTok{function}\NormalTok{(AtSEQ\_entry) \{}
  \CommentTok{\# Combine Host Sequence and Current Oomycetal Sequence}
\NormalTok{  eSEQs\_to\_map }\OtherTok{\textless{}{-}} \FunctionTok{data.frame}\NormalTok{(}\AttributeTok{seqs =} \FunctionTok{c}\NormalTok{(}\FunctionTok{as.character}\NormalTok{(AtSEQ\_entry), }\FunctionTok{as.character}\NormalTok{(eSEQ}\SpecialCharTok{$}\NormalTok{seqs[i])))}
  
  \CommentTok{\# Convert to DNAStringSet for Alignment}
\NormalTok{  eSEQs\_to\_map }\OtherTok{\textless{}{-}}\NormalTok{ Biostrings}\SpecialCharTok{::}\FunctionTok{DNAStringSet}\NormalTok{(eSEQs\_to\_map}\SpecialCharTok{$}\NormalTok{seqs)}
  
  \CommentTok{\# Align Sequences}
\NormalTok{  eSEQs\_aln }\OtherTok{\textless{}{-}}\NormalTok{ DECIPHER}\SpecialCharTok{::}\FunctionTok{AlignSeqs}\NormalTok{(eSEQs\_to\_map, }\AttributeTok{processors =} \DecValTok{1}\NormalTok{, }\AttributeTok{verbose =} \ConstantTok{FALSE}\NormalTok{)}
  
  \CommentTok{\# Compute Distance Matrix}
\NormalTok{  dist\_all }\OtherTok{\textless{}{-}}\NormalTok{ DECIPHER}\SpecialCharTok{::}\FunctionTok{DistanceMatrix}\NormalTok{(eSEQs\_aln, }\AttributeTok{processors =} \DecValTok{1}\NormalTok{, }\AttributeTok{verbose =} \ConstantTok{FALSE}\NormalTok{)}
  
  \CommentTok{\# Return the Distance Between Host and Oomycetal Sequence}
  \FunctionTok{return}\NormalTok{(dist\_all[}\DecValTok{1}\NormalTok{, }\DecValTok{2}\NormalTok{])}
\NormalTok{\}}

\CommentTok{\# Calculate Distances for All Eukaryotic Sequences}
\ControlFlowTok{for}\NormalTok{ (i }\ControlFlowTok{in} \DecValTok{1}\SpecialCharTok{:}\FunctionTok{nrow}\NormalTok{(eSEQ)) \{}
\NormalTok{  distances[, i] }\OtherTok{\textless{}{-}} \FunctionTok{apply}\NormalTok{(AtSEQ, }\DecValTok{1}\NormalTok{, map)}
\NormalTok{\}}

\CommentTok{\# Compute Similarity Scores}
\NormalTok{similarities }\OtherTok{\textless{}{-}} \DecValTok{1} \SpecialCharTok{{-}}\NormalTok{ distances}

\CommentTok{\# Append OTU Abundance to Similarity Matrix}
\NormalTok{similarities }\OtherTok{\textless{}{-}} \FunctionTok{rbind}\NormalTok{(similarities, }\AttributeTok{n\_seqs =} \FunctionTok{rowSums}\NormalTok{(oCOUNT[}\FunctionTok{rownames}\NormalTok{(oCOUNT) }\SpecialCharTok{\%in\%} \FunctionTok{colnames}\NormalTok{(distances), ]))}

\CommentTok{\# Display the Similarity Matrix}
\CommentTok{\#pander(t(similarities))}
\FunctionTok{kable}\NormalTok{(}\FunctionTok{t}\NormalTok{(similarities), }\AttributeTok{caption =} \StringTok{"OTUs mapping to GIGANTEA"}\NormalTok{,}
      \AttributeTok{col.names =} \FunctionTok{c}\NormalTok{(}\StringTok{"OTUs"}\NormalTok{, }\StringTok{"Similarity"}\NormalTok{, }\StringTok{"n\_seqs"}\NormalTok{))}
\end{Highlighting}
\end{Shaded}

\begin{longtable}[]{@{}lrr@{}}
\caption{OTUs mapping to GIGANTEA}\tabularnewline
\toprule\noalign{}
OTUs & Similarity & n\_seqs \\
\midrule\noalign{}
\endfirsthead
\toprule\noalign{}
OTUs & Similarity & n\_seqs \\
\midrule\noalign{}
\endhead
\bottomrule\noalign{}
\endlastfoot
oOTU1 & 1 & 354385 \\
\end{longtable}

\begin{Shaded}
\begin{Highlighting}[]
\CommentTok{\# Identify and Remove OTUs with Similarity Below 99\%}
\NormalTok{OTUs\_to\_remove }\OtherTok{\textless{}{-}} \FunctionTok{colnames}\NormalTok{(similarities)[similarities[}\DecValTok{1}\NormalTok{, ] }\SpecialCharTok{\textless{}} \FloatTok{0.99}\NormalTok{]}
\NormalTok{oCOUNT }\OtherTok{\textless{}{-}}\NormalTok{ oCOUNT[}\SpecialCharTok{!}\NormalTok{(}\FunctionTok{rownames}\NormalTok{(oCOUNT) }\SpecialCharTok{\%in\%}\NormalTok{ OTUs\_to\_remove), ]}
\NormalTok{oTAXA }\OtherTok{\textless{}{-}}\NormalTok{ oTAXA[}\SpecialCharTok{!}\NormalTok{(}\FunctionTok{rownames}\NormalTok{(oTAXA) }\SpecialCharTok{\%in\%}\NormalTok{ OTUs\_to\_remove), ]}

\CommentTok{\# Further Remove OTUs with Zero Abundance After Filtering}
\NormalTok{zeroOTUs\_to\_remove }\OtherTok{\textless{}{-}} \FunctionTok{names}\NormalTok{(}\FunctionTok{which}\NormalTok{(}\FunctionTok{colSums}\NormalTok{(oCOUNT) }\SpecialCharTok{==} \DecValTok{0}\NormalTok{))}
\NormalTok{oCOUNT }\OtherTok{\textless{}{-}}\NormalTok{ oCOUNT[, }\SpecialCharTok{!}\NormalTok{(}\FunctionTok{colnames}\NormalTok{(oCOUNT) }\SpecialCharTok{\%in\%}\NormalTok{ zeroOTUs\_to\_remove)]}
\NormalTok{oTAXA }\OtherTok{\textless{}{-}}\NormalTok{ oTAXA[}\SpecialCharTok{!}\NormalTok{(}\FunctionTok{rownames}\NormalTok{(oTAXA) }\SpecialCharTok{\%in\%}\NormalTok{ zeroOTUs\_to\_remove), ]}

\CommentTok{\# Report Removed OTUs}
\ControlFlowTok{if}\NormalTok{ (}\FunctionTok{length}\NormalTok{(OTUs\_to\_remove) }\SpecialCharTok{\textgreater{}} \DecValTok{0}\NormalTok{) \{}
  \FunctionTok{print}\NormalTok{(}\FunctionTok{paste0}\NormalTok{(}\StringTok{"The following eukaryotic OTUs were removed: "}\NormalTok{, }\FunctionTok{paste}\NormalTok{(OTUs\_to\_remove, }\AttributeTok{collapse =} \StringTok{", "}\NormalTok{)))}
\NormalTok{\} }\ControlFlowTok{else}\NormalTok{ \{}
  \FunctionTok{print}\NormalTok{(}\StringTok{"No eukaryotic OTUs were removed"}\NormalTok{)}
\NormalTok{\}}
\end{Highlighting}
\end{Shaded}

\begin{verbatim}
[1] "No eukaryotic OTUs were removed"
\end{verbatim}

\begin{Shaded}
\begin{Highlighting}[]
\CommentTok{\# Clean Up Temporary Variables}
\FunctionTok{rm}\NormalTok{(eSEQ, AtSEQ, similarities, OTUs\_to\_remove, zeroOTUs\_to\_remove, distances)}
\end{Highlighting}
\end{Shaded}

\section{Sequencing depth}\label{sequencing-depth}

\subsection{Function for Boxplot of Sequencing
Depth}\label{function-for-boxplot-of-sequencing-depth}

Purpose: \texttt{seq\_nr\_plot}` This function creates a boxplot for
sequencing depth across conditions and lines, displayed on a log2 scale.
It uses jitter to display individual points. samples\_plot: This
function creates a barplot for sequencing depth per sample, grouped by
condition and line.

\begin{Shaded}
\begin{Highlighting}[]
\NormalTok{seq\_nr\_plot }\OtherTok{\textless{}{-}} \ControlFlowTok{function}\NormalTok{(df, }\AttributeTok{title =} \StringTok{""}\NormalTok{)\{}
\NormalTok{  seq\_nr }\OtherTok{\textless{}{-}} \FunctionTok{ggplot}\NormalTok{(}\AttributeTok{data =}\NormalTok{ df, }\FunctionTok{aes}\NormalTok{(}\AttributeTok{x =}\NormalTok{ condition, }\AttributeTok{y =}\NormalTok{ col\_sum, }\AttributeTok{fill =}\NormalTok{ line)) }\SpecialCharTok{+} 
              \FunctionTok{geom\_boxplot}\NormalTok{(}\AttributeTok{position =} \FunctionTok{position\_dodge2}\NormalTok{(}\AttributeTok{width =} \FloatTok{0.75}\NormalTok{, }\AttributeTok{preserve =} \StringTok{"single"}\NormalTok{), }\AttributeTok{outlier.colour =} \ConstantTok{NA}\NormalTok{) }\SpecialCharTok{+}
              \FunctionTok{geom\_jitter}\NormalTok{(}\FunctionTok{aes}\NormalTok{(}\AttributeTok{col =}\NormalTok{ line, }\AttributeTok{group =}\NormalTok{ line), }\AttributeTok{size =} \DecValTok{1}\NormalTok{, }\AttributeTok{position =} \FunctionTok{position\_jitterdodge}\NormalTok{(}\AttributeTok{jitter.width =} \FloatTok{0.2}\NormalTok{, }\AttributeTok{dodge.width =} \FloatTok{0.75}\NormalTok{)) }\SpecialCharTok{+}
              \FunctionTok{facet\_grid}\NormalTok{(. }\SpecialCharTok{\textasciitilde{}}\NormalTok{ condition, }\AttributeTok{scales =} \StringTok{\textquotesingle{}free\_x\textquotesingle{}}\NormalTok{, }\AttributeTok{space =} \StringTok{"free"}\NormalTok{) }\SpecialCharTok{+}
              \FunctionTok{theme\_bw}\NormalTok{() }\SpecialCharTok{+}
              \FunctionTok{theme}\NormalTok{(}\AttributeTok{legend.position =} \StringTok{"bottom"}\NormalTok{, }\AttributeTok{axis.text.x =} \FunctionTok{element\_text}\NormalTok{(}\AttributeTok{size =} \DecValTok{5}\NormalTok{)) }\SpecialCharTok{+}
              \FunctionTok{ylab}\NormalTok{(}\StringTok{"log2(seq)"}\NormalTok{) }\SpecialCharTok{+}
              \FunctionTok{scale\_fill\_manual}\NormalTok{(}\AttributeTok{values =}\NormalTok{ level\_cols\_line) }\SpecialCharTok{+}
              \FunctionTok{scale\_color\_manual}\NormalTok{(}\AttributeTok{values =}\NormalTok{ level\_cols\_line) }\SpecialCharTok{+}
              \FunctionTok{scale\_y\_continuous}\NormalTok{(}\AttributeTok{trans =} \StringTok{\textquotesingle{}log2\textquotesingle{}}\NormalTok{) }\SpecialCharTok{+}
              \FunctionTok{ggtitle}\NormalTok{(title)}
  \FunctionTok{return}\NormalTok{(seq\_nr)}
\NormalTok{\}}

\CommentTok{\# Function for Barplot of Sequencing Depth by Sample}
\NormalTok{samples\_plot }\OtherTok{\textless{}{-}} \ControlFlowTok{function}\NormalTok{(df)\{}
\NormalTok{  samples }\OtherTok{\textless{}{-}} \FunctionTok{ggplot}\NormalTok{(df, }\FunctionTok{aes}\NormalTok{(}\AttributeTok{x =}\NormalTok{ x, }\AttributeTok{y =}\NormalTok{ col\_sum, }\AttributeTok{fill =}\NormalTok{ line, }\AttributeTok{color =}\NormalTok{ line)) }\SpecialCharTok{+}
               \FunctionTok{geom\_bar}\NormalTok{(}\AttributeTok{width =} \DecValTok{1}\NormalTok{, }\AttributeTok{linewidth =} \FloatTok{0.1}\NormalTok{, }\AttributeTok{stat =} \StringTok{"identity"}\NormalTok{, }\AttributeTok{col =} \StringTok{"black"}\NormalTok{, }\AttributeTok{lty =} \FloatTok{0.5}\NormalTok{) }\SpecialCharTok{+}
               \FunctionTok{scale\_fill\_manual}\NormalTok{(}\AttributeTok{values =}\NormalTok{ level\_cols\_line) }\SpecialCharTok{+}
               \FunctionTok{scale\_colour\_manual}\NormalTok{(}\AttributeTok{values =}\NormalTok{ level\_cols\_line) }\SpecialCharTok{+}
               \FunctionTok{xlab}\NormalTok{(}\StringTok{"samples"}\NormalTok{) }\SpecialCharTok{+}
               \FunctionTok{ylab}\NormalTok{(}\StringTok{"seq / samples"}\NormalTok{) }\SpecialCharTok{+}
               \FunctionTok{facet\_grid}\NormalTok{(. }\SpecialCharTok{\textasciitilde{}}\NormalTok{ condition, }\AttributeTok{scales =} \StringTok{\textquotesingle{}free\_x\textquotesingle{}}\NormalTok{, }\AttributeTok{space =} \StringTok{"free"}\NormalTok{) }\SpecialCharTok{+}
               \FunctionTok{theme\_bw}\NormalTok{() }\SpecialCharTok{+}
               \FunctionTok{theme}\NormalTok{(}\AttributeTok{axis.text.x =} \FunctionTok{element\_blank}\NormalTok{(), }\AttributeTok{axis.ticks.x =} \FunctionTok{element\_blank}\NormalTok{(),}
                     \AttributeTok{strip.background =} \FunctionTok{element\_blank}\NormalTok{(), }\AttributeTok{strip.text.x =} \FunctionTok{element\_blank}\NormalTok{(),}
                     \AttributeTok{legend.position =} \StringTok{"none"}\NormalTok{)}
  \FunctionTok{return}\NormalTok{(samples)}
\NormalTok{\}}
\end{Highlighting}
\end{Shaded}

\subsection{Figure - Bacteria Sequencing Depth
Plots}\label{figure---bacteria-sequencing-depth-plots}

Purpose: Boxplot: Shows the distribution of sequencing depth for
bacteria across different lines and conditions. Barplot: Displays
sequencing depth per sample, ordered by condition and sequence depth.

\subsection{Figure - Fungi Sequencing Depth
Plots}\label{figure---fungi-sequencing-depth-plots}

Purpose: Boxplot: Similar to bacteria, shows the distribution of
sequencing depth for fungi. Barplot: Displays sequencing depth per
sample for fungi.

\subsection{Figure - Oomycetes Sequencing Depth
Plots}\label{figure---oomycetes-sequencing-depth-plots}

Purpose: Boxplot: Shows the sequencing depth distribution for oomycetes.
Barplot: Displays sequencing depth per sample for oomycetes.

\begin{center}
\pandocbounded{\includegraphics[keepaspectratio]{05_Microbiome_analysis_files/figure-pdf/Sequencing depth-1.pdf}}
\end{center}

\section{Rarefaction Analysis}\label{rarefaction-analysis}

\subsection{Bacteria Sequencing Depth
Test}\label{bacteria-sequencing-depth-test}

Explanation: Group Factor: Groups the bacteria samples by combining the
line and condition columns from bDESIGN. Sequencing Depth Calculation:
Computes the total number of sequences for each sample. Kruskal-Wallis
Test: Tests if sequencing depth differs significantly between the
groups. Rarefaction Threshold: Determines the minimum sequencing depth
for rarefaction, rounded down to the nearest 100.

\begin{Shaded}
\begin{Highlighting}[]
\CommentTok{\# Bacteria Sequencing Depth Test}

\CommentTok{\# Create group factor by concatenating line and condition for each sample}
\NormalTok{group }\OtherTok{\textless{}{-}} \FunctionTok{paste}\NormalTok{(bDESIGN}\SpecialCharTok{$}\NormalTok{line, bDESIGN}\SpecialCharTok{$}\NormalTok{condition, }\AttributeTok{sep =} \StringTok{"\_"}\NormalTok{) }\SpecialCharTok{\%\textgreater{}\%}\NormalTok{ as.factor}

\CommentTok{\# Calculate sequencing depth for each sample}
\NormalTok{sample\_depth }\OtherTok{\textless{}{-}} \FunctionTok{colSums}\NormalTok{(bCOUNT)}

\CommentTok{\# Perform Kruskal{-}Wallis Test for differences in sequencing depth between groups}
\NormalTok{DAT\_krusk }\OtherTok{\textless{}{-}}\NormalTok{ coin}\SpecialCharTok{::}\FunctionTok{kruskal\_test}\NormalTok{(sample\_depth }\SpecialCharTok{\textasciitilde{}}\NormalTok{ group)}
\NormalTok{DAT\_krusk  }\CommentTok{\# Display test results}
\end{Highlighting}
\end{Shaded}

\begin{verbatim}

    Asymptotic Kruskal-Wallis Test

data:  sample_depth by
     group (159_minus, 159_plus, 318_minus, 318_plus, WT_minus, WT_plus)
chi-squared = 12.281, df = 5, p-value = 0.03114
\end{verbatim}

\begin{Shaded}
\begin{Highlighting}[]
\CommentTok{\# Determine rarefaction threshold (minimum sequencing depth rounded down to nearest 100)}
\NormalTok{brare\_threshold }\OtherTok{\textless{}{-}} \FunctionTok{floor}\NormalTok{(}\FunctionTok{min}\NormalTok{(sample\_depth) }\SpecialCharTok{/} \DecValTok{100}\NormalTok{) }\SpecialCharTok{*} \DecValTok{100}
\NormalTok{brare\_threshold  }\CommentTok{\# Display rarefaction threshold}
\end{Highlighting}
\end{Shaded}

\begin{verbatim}
[1] 60200
\end{verbatim}

\begin{Shaded}
\begin{Highlighting}[]
\CommentTok{\# Clean up}
\FunctionTok{rm}\NormalTok{(group, sample\_depth, DAT\_krusk)}
\end{Highlighting}
\end{Shaded}

\subsection{Fungi Sequencing Depth
Test}\label{fungi-sequencing-depth-test}

Explanation: Similar to the bacteria section, this test is applied to
the fungi dataset to check for differences in sequencing depth between
groups. Rarefaction Threshold: The rarefaction threshold for fungi is
calculated based on the minimum sequencing depth.

\begin{Shaded}
\begin{Highlighting}[]
\CommentTok{\# Fungi Sequencing Depth Test}

\CommentTok{\# Create group factor by concatenating line and condition for each sample}
\NormalTok{group }\OtherTok{\textless{}{-}} \FunctionTok{paste}\NormalTok{(fDESIGN}\SpecialCharTok{$}\NormalTok{line, fDESIGN}\SpecialCharTok{$}\NormalTok{condition, }\AttributeTok{sep =} \StringTok{"\_"}\NormalTok{) }\SpecialCharTok{\%\textgreater{}\%}\NormalTok{ as.factor}

\CommentTok{\# Calculate sequencing depth for each sample}
\NormalTok{sample\_depth }\OtherTok{\textless{}{-}} \FunctionTok{colSums}\NormalTok{(fCOUNT)}

\CommentTok{\# Perform Kruskal{-}Wallis Test for differences in sequencing depth between groups}
\NormalTok{DAT\_krusk }\OtherTok{\textless{}{-}}\NormalTok{ coin}\SpecialCharTok{::}\FunctionTok{kruskal\_test}\NormalTok{(sample\_depth }\SpecialCharTok{\textasciitilde{}}\NormalTok{ group)}
\NormalTok{DAT\_krusk  }\CommentTok{\# Display test results}
\end{Highlighting}
\end{Shaded}

\begin{verbatim}

    Asymptotic Kruskal-Wallis Test

data:  sample_depth by
     group (159_minus, 159_plus, 318_minus, 318_plus, WT_minus, WT_plus)
chi-squared = 8.7855, df = 5, p-value = 0.1179
\end{verbatim}

\begin{Shaded}
\begin{Highlighting}[]
\CommentTok{\# Determine rarefaction threshold (minimum sequencing depth rounded down to nearest 100)}
\NormalTok{frare\_threshold }\OtherTok{\textless{}{-}} \FunctionTok{floor}\NormalTok{(}\FunctionTok{min}\NormalTok{(sample\_depth) }\SpecialCharTok{/} \DecValTok{100}\NormalTok{) }\SpecialCharTok{*} \DecValTok{100}
\NormalTok{frare\_threshold  }\CommentTok{\# Display rarefaction threshold}
\end{Highlighting}
\end{Shaded}

\begin{verbatim}
[1] 800
\end{verbatim}

\begin{Shaded}
\begin{Highlighting}[]
\CommentTok{\# Clean up}
\FunctionTok{rm}\NormalTok{(group, sample\_depth, DAT\_krusk)}
\end{Highlighting}
\end{Shaded}

\subsection{Oomycetes Sequencing Depth
Test}\label{oomycetes-sequencing-depth-test}

Explanation: The same process is applied to the oomycetes dataset,
testing for differences in sequencing depth between groups and
determining the rarefaction threshold.

\begin{Shaded}
\begin{Highlighting}[]
\CommentTok{\# Oomycetes Sequencing Depth Test}

\CommentTok{\# Create group factor by concatenating line and condition for each sample}
\NormalTok{group }\OtherTok{\textless{}{-}} \FunctionTok{paste}\NormalTok{(oDESIGN}\SpecialCharTok{$}\NormalTok{line, oDESIGN}\SpecialCharTok{$}\NormalTok{condition, }\AttributeTok{sep =} \StringTok{"\_"}\NormalTok{) }\SpecialCharTok{\%\textgreater{}\%}\NormalTok{ as.factor}

\CommentTok{\# Calculate sequencing depth for each sample}
\NormalTok{sample\_depth }\OtherTok{\textless{}{-}} \FunctionTok{colSums}\NormalTok{(oCOUNT)}

\CommentTok{\# Perform Kruskal{-}Wallis Test for differences in sequencing depth between groups}
\NormalTok{DAT\_krusk }\OtherTok{\textless{}{-}}\NormalTok{ coin}\SpecialCharTok{::}\FunctionTok{kruskal\_test}\NormalTok{(sample\_depth }\SpecialCharTok{\textasciitilde{}}\NormalTok{ group)}
\NormalTok{DAT\_krusk  }\CommentTok{\# Display test results}
\end{Highlighting}
\end{Shaded}

\begin{verbatim}

    Asymptotic Kruskal-Wallis Test

data:  sample_depth by
     group (159_minus, 159_plus, 318_minus, 318_plus, WT_minus, WT_plus)
chi-squared = 8.4511, df = 5, p-value = 0.1331
\end{verbatim}

\begin{Shaded}
\begin{Highlighting}[]
\CommentTok{\# Determine rarefaction threshold (minimum sequencing depth rounded down to nearest 100)}
\NormalTok{orare\_threshold }\OtherTok{\textless{}{-}} \FunctionTok{floor}\NormalTok{(}\FunctionTok{min}\NormalTok{(sample\_depth) }\SpecialCharTok{/} \DecValTok{100}\NormalTok{) }\SpecialCharTok{*} \DecValTok{100}
\NormalTok{orare\_threshold  }\CommentTok{\# Display rarefaction threshold}
\end{Highlighting}
\end{Shaded}

\begin{verbatim}
[1] 5700
\end{verbatim}

\begin{Shaded}
\begin{Highlighting}[]
\CommentTok{\# Clean up}
\FunctionTok{rm}\NormalTok{(group, sample\_depth, DAT\_krusk)}
\end{Highlighting}
\end{Shaded}

\textbf{Conclusion:} Because there are significant differences between
the treatments in bacteria, we normalize the data by rarefication for
diversity comparisions (see Weiss et al.~(2017), Microbiome Journal).
For consistency we do it for bacteria, fungi and oomycetes. We define
the rarefaction threshold per sample to \ensuremath{6.02\times 10^{4}}
for bacteria, 800 for fungi and 5700 for oomycetes.

\subsection{Rarefaction}\label{rarefaction}

\begin{Shaded}
\begin{Highlighting}[]
\CommentTok{\# Bacteria Rarefaction}
\NormalTok{bCOUNT\_norm }\OtherTok{\textless{}{-}} \FunctionTok{t}\NormalTok{(vegan}\SpecialCharTok{::}\FunctionTok{rrarefy}\NormalTok{(}\FunctionTok{t}\NormalTok{(bCOUNT), brare\_threshold)) }\SpecialCharTok{\%\textgreater{}\%} \FunctionTok{as.data.frame}\NormalTok{()}
\NormalTok{bCOUNT\_norm }\OtherTok{\textless{}{-}}\NormalTok{ bCOUNT\_norm[}\FunctionTok{rowSums}\NormalTok{(bCOUNT\_norm) }\SpecialCharTok{\textgreater{}} \DecValTok{0}\NormalTok{, ]  }\CommentTok{\# Remove rows with 0 sequences}
\NormalTok{bTAXA\_norm }\OtherTok{\textless{}{-}}\NormalTok{ bTAXA[}\FunctionTok{rownames}\NormalTok{(bTAXA) }\SpecialCharTok{\%in\%} \FunctionTok{rownames}\NormalTok{(bCOUNT\_norm), ]}

\CommentTok{\# Fungi Rarefaction}
\NormalTok{fCOUNT\_norm }\OtherTok{\textless{}{-}} \FunctionTok{t}\NormalTok{(vegan}\SpecialCharTok{::}\FunctionTok{rrarefy}\NormalTok{(}\FunctionTok{t}\NormalTok{(fCOUNT), frare\_threshold)) }\SpecialCharTok{\%\textgreater{}\%} \FunctionTok{as.data.frame}\NormalTok{()}
\NormalTok{fCOUNT\_norm }\OtherTok{\textless{}{-}}\NormalTok{ fCOUNT\_norm[}\FunctionTok{rowSums}\NormalTok{(fCOUNT\_norm) }\SpecialCharTok{\textgreater{}} \DecValTok{0}\NormalTok{, ]  }\CommentTok{\# Remove rows with 0 sequences}
\NormalTok{fTAXA\_norm }\OtherTok{\textless{}{-}}\NormalTok{ fTAXA[}\FunctionTok{rownames}\NormalTok{(fTAXA) }\SpecialCharTok{\%in\%} \FunctionTok{rownames}\NormalTok{(fCOUNT\_norm), ]}

\CommentTok{\# Oomycetes Rarefaction}
\NormalTok{oCOUNT\_norm }\OtherTok{\textless{}{-}} \FunctionTok{t}\NormalTok{(vegan}\SpecialCharTok{::}\FunctionTok{rrarefy}\NormalTok{(}\FunctionTok{t}\NormalTok{(oCOUNT), orare\_threshold)) }\SpecialCharTok{\%\textgreater{}\%} \FunctionTok{as.data.frame}\NormalTok{()}
\NormalTok{oCOUNT\_norm }\OtherTok{\textless{}{-}}\NormalTok{ oCOUNT\_norm[}\FunctionTok{rowSums}\NormalTok{(oCOUNT\_norm) }\SpecialCharTok{\textgreater{}} \DecValTok{0}\NormalTok{, ]  }\CommentTok{\# Remove rows with 0 sequences}
\NormalTok{oTAXA\_norm }\OtherTok{\textless{}{-}}\NormalTok{ oTAXA[}\FunctionTok{rownames}\NormalTok{(oTAXA) }\SpecialCharTok{\%in\%} \FunctionTok{rownames}\NormalTok{(oCOUNT\_norm), ]}
\end{Highlighting}
\end{Shaded}

\subsection{Bacteria Rarefaction Plot}\label{bacteria-rarefaction-plot}

Rarefaction Plot: A line plot is created with sequencing depth on the
x-axis and the number of ASVs on the y-axis. A red vertical line
indicates the rarefaction threshold.

\subsection{Fungi Rarefaction Plot}\label{fungi-rarefaction-plot}

\subsection{Oomycetes Rarefaction
Plot}\label{oomycetes-rarefaction-plot}

\subsection{Combine Rarefaction Plots}\label{combine-rarefaction-plots}

\begin{center}
\pandocbounded{\includegraphics[keepaspectratio]{05_Microbiome_analysis_files/figure-pdf/Rarefaction plot-1.pdf}}
\end{center}

\section{Outlier Detection}\label{outlier-detection}

\subsection{Sample Control}\label{sample-control}

We use the method CLOUD developped by Montassier et al.~2018, which is a
non.parametric detection test for outliers. We perform the test with
Bray-Curtis distances from the rarefied data for each substrate and each
plastic treatment individually. We set the number of nearest neighbors
to 60\% of the samples size and chose an empirical outlier percentile of
0.1. We remove all outliers from our data.

\begin{Shaded}
\begin{Highlighting}[]
\CommentTok{\# Function to Detect Outliers}
\NormalTok{detect\_outliers }\OtherTok{\textless{}{-}} \ControlFlowTok{function}\NormalTok{(DESIGN, COUNT\_norm, line) \{}
  
  \CommentTok{\# Subset the samples for the specified line}
\NormalTok{  samples\_sub }\OtherTok{\textless{}{-}}\NormalTok{ DESIGN}\SpecialCharTok{$}\NormalTok{sample\_name[DESIGN}\SpecialCharTok{$}\NormalTok{line }\SpecialCharTok{\%in\%}\NormalTok{ line]}
  
  \CommentTok{\# Subset the count matrix for the selected samples}
\NormalTok{  DAT\_sub }\OtherTok{\textless{}{-}}\NormalTok{ COUNT\_norm[, }\FunctionTok{colnames}\NormalTok{(COUNT\_norm) }\SpecialCharTok{\%in\%}\NormalTok{ samples\_sub]}
  
  \CommentTok{\# Compute Bray{-}Curtis dissimilarity between the samples}
\NormalTok{  bdist }\OtherTok{\textless{}{-}} \FunctionTok{vegdist}\NormalTok{(}\FunctionTok{t}\NormalTok{(DAT\_sub), }\AttributeTok{method =} \StringTok{"bray"}\NormalTok{)}
  
  \CommentTok{\# Set the number of nearest neighbors (60\% of the sample size)}
\NormalTok{  k }\OtherTok{\textless{}{-}} \FunctionTok{floor}\NormalTok{(}\FunctionTok{ncol}\NormalTok{(DAT\_sub) }\SpecialCharTok{*} \FloatTok{0.6}\NormalTok{)}
  
  \CommentTok{\# Apply the CLOUD test from Montassier et al. 2018}
\NormalTok{  test }\OtherTok{\textless{}{-}} \FunctionTok{piecewise\_kn\_V1}\NormalTok{(}\AttributeTok{d =}\NormalTok{ bdist, }\AttributeTok{test.ix =} \FunctionTok{colnames}\NormalTok{(DAT\_sub), }\AttributeTok{k =}\NormalTok{ k)}
  
  \CommentTok{\# Identify outliers (samples with p{-}values \textless{}= 0.1)}
\NormalTok{  outliers }\OtherTok{\textless{}{-}} \FunctionTok{colnames}\NormalTok{(DAT\_sub)[test[[}\StringTok{"pvals"}\NormalTok{]] }\SpecialCharTok{\textless{}=} \FloatTok{0.1}\NormalTok{]}
  
  \FunctionTok{return}\NormalTok{(outliers)}
\NormalTok{\}}
\end{Highlighting}
\end{Shaded}

\subsection{Outlier Detection}\label{outlier-detection-1}

\begin{Shaded}
\begin{Highlighting}[]
\CommentTok{\# Detect Outliers for Bacteria, Fungi, and Oomycetes}

\CommentTok{\# Bacteria}
\NormalTok{boutliers }\OtherTok{\textless{}{-}} \FunctionTok{detect\_outliers}\NormalTok{(}\AttributeTok{DESIGN =}\NormalTok{ bDESIGN, }\AttributeTok{COUNT\_norm =}\NormalTok{ bCOUNT\_norm, }\AttributeTok{line =} \FunctionTok{unique}\NormalTok{(bDESIGN}\SpecialCharTok{$}\NormalTok{line))}

\CommentTok{\# Fungi}
\NormalTok{foutliers }\OtherTok{\textless{}{-}} \FunctionTok{detect\_outliers}\NormalTok{(}\AttributeTok{DESIGN =}\NormalTok{ fDESIGN, }\AttributeTok{COUNT\_norm =}\NormalTok{ fCOUNT\_norm, }\AttributeTok{line =} \FunctionTok{unique}\NormalTok{(fDESIGN}\SpecialCharTok{$}\NormalTok{line))}

\CommentTok{\# Oomycetes}
\NormalTok{ooutliers }\OtherTok{\textless{}{-}} \FunctionTok{detect\_outliers}\NormalTok{(}\AttributeTok{DESIGN =}\NormalTok{ oDESIGN, }\AttributeTok{COUNT\_norm =}\NormalTok{ oCOUNT\_norm, }\AttributeTok{line =} \FunctionTok{unique}\NormalTok{(oDESIGN}\SpecialCharTok{$}\NormalTok{line))}
\end{Highlighting}
\end{Shaded}

\subsection{Outlier Detection Summary}\label{outlier-detection-summary}

\begin{longtable}[]{@{}
  >{\centering\arraybackslash}p{(\linewidth - 2\tabcolsep) * \real{0.1667}}
  >{\centering\arraybackslash}p{(\linewidth - 2\tabcolsep) * \real{0.1667}}@{}}
\caption{Number of outliers}\tabularnewline
\toprule\noalign{}
\begin{minipage}[b]{\linewidth}\centering
taxa
\end{minipage} & \begin{minipage}[b]{\linewidth}\centering
outliers
\end{minipage} \\
\midrule\noalign{}
\endfirsthead
\toprule\noalign{}
\begin{minipage}[b]{\linewidth}\centering
taxa
\end{minipage} & \begin{minipage}[b]{\linewidth}\centering
outliers
\end{minipage} \\
\midrule\noalign{}
\endhead
\bottomrule\noalign{}
\endlastfoot
Bacteria & 3 \\
Fungi & 5 \\
Oomycetes & 4 \\
\end{longtable}

\begin{longtable}[]{@{}lr@{}}
\caption{Number of outliers}\tabularnewline
\toprule\noalign{}
taxa & Outliers \\
\midrule\noalign{}
\endfirsthead
\toprule\noalign{}
taxa & Outliers \\
\midrule\noalign{}
\endhead
\bottomrule\noalign{}
\endlastfoot
Bacteria & 3 \\
Fungi & 5 \\
Oomycetes & 4 \\
\end{longtable}

\begin{verbatim}
[1] "Bacteria: FLD0297, FLD0340, FLD0348"
\end{verbatim}

\begin{verbatim}
[1] "Fungi: FLD0117, FLD0131, FLD0135, FLD0146, FLD0151"
\end{verbatim}

\begin{verbatim}
[1] "Oomycetes: FLD0195, FLD0219, FLD0227, FLD0232"
\end{verbatim}

\subsection{Remove outliers for
bacteria}\label{remove-outliers-for-bacteria}

\begin{Shaded}
\begin{Highlighting}[]
\CommentTok{\# Bacteria Outlier Removal}

\CommentTok{\# Remove outliers from the design table}
\NormalTok{bDESIGN }\OtherTok{\textless{}{-}}\NormalTok{ bDESIGN[}\SpecialCharTok{!}\NormalTok{(bDESIGN}\SpecialCharTok{$}\NormalTok{sample\_name }\SpecialCharTok{\%in\%}\NormalTok{ boutliers),]}

\CommentTok{\# Remove outliers from the normalized count matrix}
\NormalTok{bCOUNT\_norm }\OtherTok{\textless{}{-}}\NormalTok{ bCOUNT\_norm[, }\SpecialCharTok{!}\NormalTok{(}\FunctionTok{colnames}\NormalTok{(bCOUNT\_norm) }\SpecialCharTok{\%in\%}\NormalTok{ boutliers)]}

\CommentTok{\# Identify and remove ASVs with zero counts across all samples}
\NormalTok{ASVs\_to\_remove }\OtherTok{\textless{}{-}} \FunctionTok{names}\NormalTok{(}\FunctionTok{which}\NormalTok{(}\FunctionTok{rowSums}\NormalTok{(bCOUNT\_norm) }\SpecialCharTok{==} \DecValTok{0}\NormalTok{))}
\NormalTok{bCOUNT\_norm }\OtherTok{\textless{}{-}}\NormalTok{ bCOUNT\_norm[}\SpecialCharTok{!}\NormalTok{(}\FunctionTok{rownames}\NormalTok{(bCOUNT\_norm) }\SpecialCharTok{\%in\%}\NormalTok{ ASVs\_to\_remove),]}

\CommentTok{\# Update the taxonomy table to remove ASVs with zero counts}
\NormalTok{bTAXA\_norm }\OtherTok{\textless{}{-}}\NormalTok{ bTAXA\_norm[}\SpecialCharTok{!}\NormalTok{(}\FunctionTok{rownames}\NormalTok{(bTAXA\_norm) }\SpecialCharTok{\%in\%}\NormalTok{ ASVs\_to\_remove),]}
\end{Highlighting}
\end{Shaded}

\subsection{Remove outliers for fungi}\label{remove-outliers-for-fungi}

\begin{Shaded}
\begin{Highlighting}[]
\CommentTok{\# Fungi Outlier Removal}

\CommentTok{\# Remove outliers from the design table}
\NormalTok{fDESIGN }\OtherTok{\textless{}{-}}\NormalTok{ fDESIGN[}\SpecialCharTok{!}\NormalTok{(fDESIGN}\SpecialCharTok{$}\NormalTok{sample\_name }\SpecialCharTok{\%in\%}\NormalTok{ foutliers),]}

\CommentTok{\# Remove outliers from the normalized count matrix}
\NormalTok{fCOUNT\_norm }\OtherTok{\textless{}{-}}\NormalTok{ fCOUNT\_norm[, }\SpecialCharTok{!}\NormalTok{(}\FunctionTok{colnames}\NormalTok{(fCOUNT\_norm) }\SpecialCharTok{\%in\%}\NormalTok{ foutliers)]}

\CommentTok{\# Identify and remove ASVs with zero counts across all samples}
\NormalTok{ASVs\_to\_remove }\OtherTok{\textless{}{-}} \FunctionTok{names}\NormalTok{(}\FunctionTok{which}\NormalTok{(}\FunctionTok{rowSums}\NormalTok{(fCOUNT\_norm) }\SpecialCharTok{==} \DecValTok{0}\NormalTok{))}
\NormalTok{fCOUNT\_norm }\OtherTok{\textless{}{-}}\NormalTok{ fCOUNT\_norm[}\SpecialCharTok{!}\NormalTok{(}\FunctionTok{rownames}\NormalTok{(fCOUNT\_norm) }\SpecialCharTok{\%in\%}\NormalTok{ ASVs\_to\_remove),]}

\CommentTok{\# Update the taxonomy table to remove ASVs with zero counts}
\NormalTok{fTAXA\_norm }\OtherTok{\textless{}{-}}\NormalTok{ fTAXA\_norm[}\SpecialCharTok{!}\NormalTok{(}\FunctionTok{rownames}\NormalTok{(fTAXA\_norm) }\SpecialCharTok{\%in\%}\NormalTok{ ASVs\_to\_remove),]}
\end{Highlighting}
\end{Shaded}

\subsection{Remove outliers for
oomycetes}\label{remove-outliers-for-oomycetes}

\begin{Shaded}
\begin{Highlighting}[]
\CommentTok{\# Oomycetes Outlier Removal}

\CommentTok{\# Remove outliers from the design table}
\NormalTok{oDESIGN }\OtherTok{\textless{}{-}}\NormalTok{ oDESIGN[}\SpecialCharTok{!}\NormalTok{(oDESIGN}\SpecialCharTok{$}\NormalTok{sample\_name }\SpecialCharTok{\%in\%}\NormalTok{ ooutliers),]}

\CommentTok{\# Remove outliers from the normalized count matrix}
\NormalTok{oCOUNT\_norm }\OtherTok{\textless{}{-}}\NormalTok{ oCOUNT\_norm[, }\SpecialCharTok{!}\NormalTok{(}\FunctionTok{colnames}\NormalTok{(oCOUNT\_norm) }\SpecialCharTok{\%in\%}\NormalTok{ ooutliers)]}

\CommentTok{\# Identify and remove ASVs with zero counts across all samples}
\NormalTok{ASVs\_to\_remove }\OtherTok{\textless{}{-}} \FunctionTok{names}\NormalTok{(}\FunctionTok{which}\NormalTok{(}\FunctionTok{rowSums}\NormalTok{(oCOUNT\_norm) }\SpecialCharTok{==} \DecValTok{0}\NormalTok{))}
\NormalTok{oCOUNT\_norm }\OtherTok{\textless{}{-}}\NormalTok{ oCOUNT\_norm[}\SpecialCharTok{!}\NormalTok{(}\FunctionTok{rownames}\NormalTok{(oCOUNT\_norm) }\SpecialCharTok{\%in\%}\NormalTok{ ASVs\_to\_remove),]}

\CommentTok{\# Update the taxonomy table to remove ASVs with zero counts}
\NormalTok{oTAXA\_norm }\OtherTok{\textless{}{-}}\NormalTok{ oTAXA\_norm[}\SpecialCharTok{!}\NormalTok{(}\FunctionTok{rownames}\NormalTok{(oTAXA\_norm) }\SpecialCharTok{\%in\%}\NormalTok{ ASVs\_to\_remove),]}
\end{Highlighting}
\end{Shaded}

\section{Sample Size}\label{sample-size}

We end up with the following number of samples per treatment for the
analysis.

\begin{longtable}[]{@{}lrr@{}}
\caption{Bacteria: Sample profile}\tabularnewline
\toprule\noalign{}
Line & BXminus & BXplus \\
\midrule\noalign{}
\endfirsthead
\toprule\noalign{}
Line & BXminus & BXplus \\
\midrule\noalign{}
\endhead
\bottomrule\noalign{}
\endlastfoot
WT & 10 & 10 \\
159 & 10 & 8 \\
318 & 9 & 10 \\
\end{longtable}

\begin{longtable}[]{@{}lrr@{}}
\caption{Fungi: Sample profile}\tabularnewline
\toprule\noalign{}
Line & BXminus & BXplus \\
\midrule\noalign{}
\endfirsthead
\toprule\noalign{}
Line & BXminus & BXplus \\
\midrule\noalign{}
\endhead
\bottomrule\noalign{}
\endlastfoot
WT & 8 & 7 \\
159 & 9 & 6 \\
318 & 7 & 9 \\
\end{longtable}

\begin{longtable}[]{@{}lrr@{}}
\caption{Oomycetes: Sample profile}\tabularnewline
\toprule\noalign{}
Line & BXminus & BXplus \\
\midrule\noalign{}
\endfirsthead
\toprule\noalign{}
Line & BXminus & BXplus \\
\midrule\noalign{}
\endhead
\bottomrule\noalign{}
\endlastfoot
WT & 9 & 8 \\
159 & 10 & 10 \\
318 & 9 & 10 \\
\end{longtable}

\section{Beta diversity}\label{beta-diversity}

\section{Background}\label{background-1}

We ge an overview of the microbial communitites between Col-0 and
\emph{mmf1-2} mutants grown on BXplus and BXminus soils. We perform
alpha and beta diversity analysis to investigate the differences in the
microbial communities between the different Arabidopsis lines and
conditions. We use the Shannon index to calculate alpha diversity. We
use the Kruskal-Wallis test to test the significance of the differences
in alpha diversity between the different Arabidopsis lines and
conditions. We use the Bray-Curtis distance metric for beta diversity
analysis. We perform PERMANOVA to test the significance of the
differences in beta diversity between the different Arabidopsis lines
and conditions. We also perform CAP analysis to investigate the effects
of line and condition on beta diversity.

\section{Setup}\label{setup-1}

\subsection{Prepare the environment}\label{prepare-the-environment}

Prepare the environment by loading the necessary libraries, setting the
seed, and loading the custom functions.

\begin{Shaded}
\begin{Highlighting}[]
\DocumentationTok{\#\#clear the object from memory}
\FunctionTok{rm}\NormalTok{(}\AttributeTok{list=}\FunctionTok{ls}\NormalTok{())}

\NormalTok{knitr}\SpecialCharTok{::}\NormalTok{opts\_chunk}\SpecialCharTok{$}\FunctionTok{set}\NormalTok{(}\AttributeTok{echo=}\ConstantTok{TRUE}\NormalTok{, }\AttributeTok{fig.align=}\StringTok{"center"}\NormalTok{)}
\FunctionTok{options}\NormalTok{(}\AttributeTok{tinytex.verbose =} \ConstantTok{TRUE}\NormalTok{)}

\CommentTok{\#set seed}
\FunctionTok{set.seed}\NormalTok{(}\DecValTok{100}\NormalTok{)}

\CommentTok{\# Set the CRAN mirror globally}
\FunctionTok{options}\NormalTok{(}\AttributeTok{repos =} \FunctionTok{c}\NormalTok{(}\AttributeTok{CRAN =} \StringTok{"https://cran.r{-}project.org"}\NormalTok{))}
\DocumentationTok{\#\# set source to file location}
\FunctionTok{library}\NormalTok{(rstudioapi)}
\CommentTok{\#setwd(dirname(getActiveDocumentContext()$path))}

\DocumentationTok{\#\# load functions}
\FunctionTok{source}\NormalTok{(}\StringTok{"Input/functions/functions.R"}\NormalTok{)}
\FunctionTok{functions}\NormalTok{(}\AttributeTok{path=}\StringTok{"Input/functions/"}\NormalTok{)}

\DocumentationTok{\#\# installs (if necessary) and loads libraries}

\ControlFlowTok{if}\NormalTok{ (}\SpecialCharTok{!}\FunctionTok{requireNamespace}\NormalTok{(}\StringTok{"BiocManager"}\NormalTok{, }\AttributeTok{quietly =} \ConstantTok{TRUE}\NormalTok{))}
    \FunctionTok{install.packages}\NormalTok{(}\StringTok{"BiocManager"}\NormalTok{)}

\CommentTok{\# Install Maaslin2}
\NormalTok{BiocManager}\SpecialCharTok{::}\FunctionTok{install}\NormalTok{(}\StringTok{"Maaslin2"}\NormalTok{)}
\FunctionTok{libraries}\NormalTok{()}
\FunctionTok{install.packages}\NormalTok{(}\StringTok{"ggnewscale"}\NormalTok{)}
\end{Highlighting}
\end{Shaded}

\begin{verbatim}

The downloaded binary packages are in
    /var/folders/wj/_hq5dng10tx61tf5cgw71trw0000gp/T//RtmpNqJfVH/downloaded_packages
\end{verbatim}

\begin{Shaded}
\begin{Highlighting}[]
\FunctionTok{library}\NormalTok{(}\StringTok{"ggnewscale"}\NormalTok{)}
\FunctionTok{library}\NormalTok{(}\StringTok{"ggpubr"}\NormalTok{)}
\FunctionTok{library}\NormalTok{(}\StringTok{"vegan"}\NormalTok{)}
\FunctionTok{library}\NormalTok{(}\StringTok{"ggpubr"}\NormalTok{)}
\FunctionTok{library}\NormalTok{(}\StringTok{"knitr"}\NormalTok{)}
\end{Highlighting}
\end{Shaded}

\subsection{Load data}\label{load-data}

Load the necessary data for the analysis and import the color palettes.

\begin{Shaded}
\begin{Highlighting}[]
\CommentTok{\# Import}

\CommentTok{\# set wd to RDS files}
\FunctionTok{setwd}\NormalTok{(}\StringTok{"Input/interim/02\_Taxa\_Analysis/"}\NormalTok{)}

\CommentTok{\#from here on we only use the normfied dataasets without host}
\CommentTok{\#for simplicity, we dont call them X\_norm\_noHost anymore}

\CommentTok{\#Import RDS files}
\NormalTok{bDESIGN }\OtherTok{\textless{}{-}} \FunctionTok{readRDS}\NormalTok{(}\StringTok{"bDESIGN.RDS"}\NormalTok{)}
\NormalTok{bCOUNT }\OtherTok{\textless{}{-}} \FunctionTok{readRDS}\NormalTok{(}\StringTok{"bCOUNT\_microbial\_load.RDS"}\NormalTok{)}
\NormalTok{bTAXA }\OtherTok{\textless{}{-}} \FunctionTok{readRDS}\NormalTok{(}\StringTok{"bTAXA\_microbial\_load.RDS"}\NormalTok{)}
\NormalTok{bPHYSEQ }\OtherTok{\textless{}{-}} \FunctionTok{readRDS}\NormalTok{(}\StringTok{"bPHYSEQ\_microbial\_load\_filt.RDS"}\NormalTok{)}

\NormalTok{fDESIGN }\OtherTok{\textless{}{-}} \FunctionTok{readRDS}\NormalTok{(}\StringTok{"fDESIGN.RDS"}\NormalTok{)}
\NormalTok{fCOUNT }\OtherTok{\textless{}{-}} \FunctionTok{readRDS}\NormalTok{(}\StringTok{"fCOUNT\_microbial\_load.RDS"}\NormalTok{)}
\NormalTok{fTAXA }\OtherTok{\textless{}{-}} \FunctionTok{readRDS}\NormalTok{(}\StringTok{"fTAXA\_microbial\_load.RDS"}\NormalTok{)}
\NormalTok{fPHYSEQ }\OtherTok{\textless{}{-}} \FunctionTok{readRDS}\NormalTok{(}\StringTok{"fPHYSEQ\_microbial\_load\_filter.RDS"}\NormalTok{)}

\NormalTok{oDESIGN }\OtherTok{\textless{}{-}} \FunctionTok{readRDS}\NormalTok{(}\StringTok{"oDESIGN.RDS"}\NormalTok{)}
\NormalTok{oCOUNT }\OtherTok{\textless{}{-}} \FunctionTok{readRDS}\NormalTok{(}\StringTok{"oCOUNT\_microbial\_load.RDS"}\NormalTok{)}
\NormalTok{oTAXA }\OtherTok{\textless{}{-}} \FunctionTok{readRDS}\NormalTok{(}\StringTok{"oTAXA\_microbial\_load.RDS"}\NormalTok{)}
\NormalTok{oPHYSEQ }\OtherTok{\textless{}{-}} \FunctionTok{readRDS}\NormalTok{(}\StringTok{"oPHYSEQ\_microbial\_load\_filter.RDS"}\NormalTok{)}

\CommentTok{\#colors}
\FunctionTok{setwd}\NormalTok{(}\StringTok{"../01\_Import\_Normalization/"}\NormalTok{)}
\NormalTok{level\_cols\_condition }\OtherTok{\textless{}{-}} \FunctionTok{readRDS}\NormalTok{(}\StringTok{"level\_cols\_condition.RDS"}\NormalTok{)}
\NormalTok{level\_cols\_line }\OtherTok{\textless{}{-}} \FunctionTok{readRDS}\NormalTok{(}\StringTok{"level\_cols\_line.RDS"}\NormalTok{)}
\end{Highlighting}
\end{Shaded}

\subsection{Modify data for analysis}\label{modify-data-for-analysis}

Filter the data to only include the WT and \emph{mmf1-2} lines.
Transpose the count tables to have samples as rows and OTUs/taxa as
columns.

\begin{Shaded}
\begin{Highlighting}[]
\CommentTok{\#Filter bacteria for only WT and 318}
\NormalTok{bDESIGN\_filt }\OtherTok{\textless{}{-}} \FunctionTok{filter}\NormalTok{(bDESIGN, line }\SpecialCharTok{\%in\%} \FunctionTok{c}\NormalTok{(}\StringTok{"WT"}\NormalTok{, }\StringTok{"318"}\NormalTok{))}
\NormalTok{bCOUNT\_t }\OtherTok{\textless{}{-}} \FunctionTok{t}\NormalTok{(bCOUNT)}
\NormalTok{bCOUNT\_filt }\OtherTok{\textless{}{-}}\NormalTok{ bCOUNT\_t[}\FunctionTok{rownames}\NormalTok{(bCOUNT\_t)}\SpecialCharTok{\%in\%}\NormalTok{bDESIGN\_filt}\SpecialCharTok{$}\NormalTok{sample\_name,]}
\NormalTok{bCOUNT\_filt }\OtherTok{\textless{}{-}} \FunctionTok{t}\NormalTok{(bCOUNT\_filt)}
\CommentTok{\#Filter bacteria for only WT and 318}
\NormalTok{fDESIGN\_filt }\OtherTok{\textless{}{-}} \FunctionTok{filter}\NormalTok{(fDESIGN, line }\SpecialCharTok{\%in\%} \FunctionTok{c}\NormalTok{(}\StringTok{"WT"}\NormalTok{, }\StringTok{"318"}\NormalTok{))}
\NormalTok{fCOUNT\_t }\OtherTok{\textless{}{-}} \FunctionTok{t}\NormalTok{(fCOUNT)}
\NormalTok{fCOUNT\_filt }\OtherTok{\textless{}{-}}\NormalTok{ fCOUNT\_t[}\FunctionTok{rownames}\NormalTok{(fCOUNT\_t)}\SpecialCharTok{\%in\%}\NormalTok{fDESIGN\_filt}\SpecialCharTok{$}\NormalTok{sample\_name,]}
\NormalTok{fCOUNT\_filt }\OtherTok{\textless{}{-}} \FunctionTok{t}\NormalTok{(fCOUNT\_filt)}
\CommentTok{\#Filter bacteria for only WT and 318}
\NormalTok{oDESIGN\_filt }\OtherTok{\textless{}{-}} \FunctionTok{filter}\NormalTok{(oDESIGN, line }\SpecialCharTok{\%in\%} \FunctionTok{c}\NormalTok{(}\StringTok{"WT"}\NormalTok{, }\StringTok{"318"}\NormalTok{))}
\NormalTok{oCOUNT\_t }\OtherTok{\textless{}{-}} \FunctionTok{t}\NormalTok{(oCOUNT)}
\NormalTok{oCOUNT\_filt }\OtherTok{\textless{}{-}}\NormalTok{ oCOUNT\_t[}\FunctionTok{rownames}\NormalTok{(oCOUNT\_t)}\SpecialCharTok{\%in\%}\NormalTok{oDESIGN\_filt}\SpecialCharTok{$}\NormalTok{sample\_name,]}
\NormalTok{oCOUNT\_filt }\OtherTok{\textless{}{-}} \FunctionTok{t}\NormalTok{(oCOUNT\_filt)}
\end{Highlighting}
\end{Shaded}

\section{Alpha Diversity}\label{alpha-diversity}

Prepare Count tables Prepare the count tables for the alpha diversity
analysis by transposing the data to have samples as rows and OTUs/taxa
as columns.

\begin{Shaded}
\begin{Highlighting}[]
\CommentTok{\# Prepare data for shannon indices}

\CommentTok{\# Transpose data to have samples as rows and OTUs/taxa as columns}
\NormalTok{bCOUNT\_filt\_t }\OtherTok{\textless{}{-}} \FunctionTok{t}\NormalTok{(bCOUNT\_filt)}
\NormalTok{fCOUNT\_filt\_t }\OtherTok{\textless{}{-}} \FunctionTok{t}\NormalTok{(fCOUNT\_filt)}
\NormalTok{oCOUNT\_filt\_t }\OtherTok{\textless{}{-}} \FunctionTok{t}\NormalTok{(oCOUNT\_filt)}
\end{Highlighting}
\end{Shaded}

\subsection{Calculate Shannon index}\label{calculate-shannon-index}

Calculate the Shannon diversity index for the bacterial, fungal, and
oomycete communities.

\begin{Shaded}
\begin{Highlighting}[]
\CommentTok{\# Calculate Shannon index using the vegan package}
\NormalTok{bSHANNON }\OtherTok{\textless{}{-}} \FunctionTok{diversity}\NormalTok{(bCOUNT\_filt\_t, }\AttributeTok{index=}\StringTok{"shannon"}\NormalTok{)}
\NormalTok{fSHANNON }\OtherTok{\textless{}{-}} \FunctionTok{diversity}\NormalTok{(fCOUNT\_filt\_t, }\AttributeTok{index=}\StringTok{"shannon"}\NormalTok{)}
\NormalTok{oSHANNON }\OtherTok{\textless{}{-}} \FunctionTok{diversity}\NormalTok{(oCOUNT\_filt\_t, }\AttributeTok{index=}\StringTok{"shannon"}\NormalTok{)}

\CommentTok{\# Print the calculated Shannon diversity indices}
\CommentTok{\#print(bSHANNON)}
\CommentTok{\#print(fSHANNON)}
\CommentTok{\#print(oSHANNON)}
\end{Highlighting}
\end{Shaded}

\subsection{Alpha diversity Bacteria}\label{alpha-diversity-bacteria}

Perform the Kruskal-Wallis test to test the significance of the
differences in bacterial alpha diversity between the different
Arabidopsis lines and conditions.

\begin{Shaded}
\begin{Highlighting}[]
\CommentTok{\# Step 1: Ensure bSHANNON is a data frame}
\ControlFlowTok{if}\NormalTok{ (}\SpecialCharTok{!}\FunctionTok{is.data.frame}\NormalTok{(bSHANNON)) \{}
\NormalTok{  bSHANNON }\OtherTok{\textless{}{-}} \FunctionTok{as.data.frame}\NormalTok{(bSHANNON)}
\NormalTok{\}}

\CommentTok{\# Step 2: Convert rownames to a column in bSHANNON}
\NormalTok{bSHANNON }\OtherTok{\textless{}{-}}\NormalTok{ bSHANNON }\SpecialCharTok{\%\textgreater{}\%}
\NormalTok{  tibble}\SpecialCharTok{::}\FunctionTok{rownames\_to\_column}\NormalTok{(}\AttributeTok{var =} \StringTok{"sample\_name"}\NormalTok{)}

\CommentTok{\# Step 3: Merge with bDESIGN\_filt}
\NormalTok{merged\_data }\OtherTok{\textless{}{-}} \FunctionTok{merge}\NormalTok{(bSHANNON, bDESIGN\_filt, }\AttributeTok{by =} \StringTok{"sample\_name"}\NormalTok{)}

\CommentTok{\# Step 4: Calculate Kruskal{-}Wallis p{-}values for each condition}
\NormalTok{p\_values }\OtherTok{\textless{}{-}}\NormalTok{ merged\_data }\SpecialCharTok{\%\textgreater{}\%}
  \FunctionTok{group\_by}\NormalTok{(condition) }\SpecialCharTok{\%\textgreater{}\%}
  \FunctionTok{summarise}\NormalTok{(}\AttributeTok{p\_value =} \FunctionTok{kruskal.test}\NormalTok{(bSHANNON }\SpecialCharTok{\textasciitilde{}}\NormalTok{ line)}\SpecialCharTok{$}\NormalTok{p.value)}

\CommentTok{\# Print the p{-}values}
\CommentTok{\#print(p\_values)}
\FunctionTok{kable}\NormalTok{(p\_values, }\AttributeTok{caption =} \StringTok{"Shannon index: Bacteria"}\NormalTok{,}
      \AttributeTok{col.names =} \FunctionTok{c}\NormalTok{(}\StringTok{"Microbiome"}\NormalTok{, }\StringTok{"p{-}value"}\NormalTok{))}
\end{Highlighting}
\end{Shaded}

\begin{longtable}[]{@{}lr@{}}
\caption{Shannon index: Bacteria}\tabularnewline
\toprule\noalign{}
Microbiome & p-value \\
\midrule\noalign{}
\endfirsthead
\toprule\noalign{}
Microbiome & p-value \\
\midrule\noalign{}
\endhead
\bottomrule\noalign{}
\endlastfoot
minus & 0.3691076 \\
plus & 0.3257514 \\
\end{longtable}

\subsection{Plot - Bacteria}\label{plot---bacteria}

\subsection{Alpha diversity for Fungi}\label{alpha-diversity-for-fungi}

Perform the Kruskal-Wallis test to test the significance of the
differences in fungal alpha diversity between the different Arabidopsis
lines and conditions.

\begin{Shaded}
\begin{Highlighting}[]
\CommentTok{\# Step 1: Ensure fSHANNON is a data frame}
\ControlFlowTok{if}\NormalTok{ (}\SpecialCharTok{!}\FunctionTok{is.data.frame}\NormalTok{(fSHANNON)) \{}
\NormalTok{  fSHANNON }\OtherTok{\textless{}{-}} \FunctionTok{as.data.frame}\NormalTok{(fSHANNON)}
\NormalTok{\}}

\CommentTok{\# Step 2: Convert rownames to a column in fSHANNON}
\NormalTok{fSHANNON }\OtherTok{\textless{}{-}}\NormalTok{ fSHANNON }\SpecialCharTok{\%\textgreater{}\%}
\NormalTok{  tibble}\SpecialCharTok{::}\FunctionTok{rownames\_to\_column}\NormalTok{(}\AttributeTok{var =} \StringTok{"sample\_name"}\NormalTok{)}

\CommentTok{\# Step 3: Merge with fDESIGN\_filt}
\NormalTok{merged\_data }\OtherTok{\textless{}{-}} \FunctionTok{merge}\NormalTok{(fSHANNON, fDESIGN\_filt, }\AttributeTok{by =} \StringTok{"sample\_name"}\NormalTok{)}
  
\CommentTok{\# Step 4: Calculate Kruskal{-}Wallis p{-}values for each condition}
\NormalTok{p\_values }\OtherTok{\textless{}{-}}\NormalTok{ merged\_data }\SpecialCharTok{\%\textgreater{}\%}
  \FunctionTok{group\_by}\NormalTok{(condition) }\SpecialCharTok{\%\textgreater{}\%}
  \FunctionTok{summarise}\NormalTok{(}\AttributeTok{p\_value =} \FunctionTok{kruskal.test}\NormalTok{(fSHANNON }\SpecialCharTok{\textasciitilde{}}\NormalTok{ line)}\SpecialCharTok{$}\NormalTok{p.value)}

\CommentTok{\# Print the p{-}values}
\CommentTok{\#print(p\_values)}
\FunctionTok{kable}\NormalTok{(p\_values, }\AttributeTok{caption =} \StringTok{"Shannon index: Fungi"}\NormalTok{,}
      \AttributeTok{col.names =} \FunctionTok{c}\NormalTok{(}\StringTok{"Microbiome"}\NormalTok{, }\StringTok{"p{-}value"}\NormalTok{))}
\end{Highlighting}
\end{Shaded}

\begin{longtable}[]{@{}lr@{}}
\caption{Shannon index: Fungi}\tabularnewline
\toprule\noalign{}
Microbiome & p-value \\
\midrule\noalign{}
\endfirsthead
\toprule\noalign{}
Microbiome & p-value \\
\midrule\noalign{}
\endhead
\bottomrule\noalign{}
\endlastfoot
minus & 0.8169613 \\
plus & 0.7110254 \\
\end{longtable}

\subsection{Plot - Fungi}\label{plot---fungi}

\subsection{Alpha diversity for
Oomyctes}\label{alpha-diversity-for-oomyctes}

Perform the Kruskal-Wallis test to test the significance of the
differences in oomycete alpha diversity between the different
Arabidopsis lines and conditions.

\begin{Shaded}
\begin{Highlighting}[]
\ControlFlowTok{if}\NormalTok{ (}\SpecialCharTok{!}\FunctionTok{is.data.frame}\NormalTok{(oSHANNON)) \{}
\NormalTok{  oSHANNON }\OtherTok{\textless{}{-}} \FunctionTok{as.data.frame}\NormalTok{(oSHANNON)}
\NormalTok{\}}

\CommentTok{\# Step 2: Convert rownames to a column in oSHANNON}
\NormalTok{oSHANNON }\OtherTok{\textless{}{-}}\NormalTok{ oSHANNON }\SpecialCharTok{\%\textgreater{}\%}
\NormalTok{  tibble}\SpecialCharTok{::}\FunctionTok{rownames\_to\_column}\NormalTok{(}\AttributeTok{var =} \StringTok{"sample\_name"}\NormalTok{)}

\CommentTok{\# Step 3: Merge with oDESIGN\_filt}
\NormalTok{merged\_data }\OtherTok{\textless{}{-}} \FunctionTok{merge}\NormalTok{(oSHANNON, oDESIGN\_filt, }\AttributeTok{by =} \StringTok{"sample\_name"}\NormalTok{)}

\CommentTok{\# Step 4: Calculate Kruskal{-}Wallis p{-}values for each condition}
\NormalTok{p\_values }\OtherTok{\textless{}{-}}\NormalTok{ merged\_data }\SpecialCharTok{\%\textgreater{}\%}
  \FunctionTok{group\_by}\NormalTok{(condition) }\SpecialCharTok{\%\textgreater{}\%}
  \FunctionTok{summarise}\NormalTok{(}\AttributeTok{p\_value =} \FunctionTok{kruskal.test}\NormalTok{(oSHANNON }\SpecialCharTok{\textasciitilde{}}\NormalTok{ line)}\SpecialCharTok{$}\NormalTok{p.value)}

\CommentTok{\# Print the p{-}values}
\CommentTok{\#print(p\_values)}
\FunctionTok{kable}\NormalTok{(p\_values, }\AttributeTok{caption =} \StringTok{"Shannon index: Oomycetes"}\NormalTok{,}
      \AttributeTok{col.names =} \FunctionTok{c}\NormalTok{(}\StringTok{"Microbiome"}\NormalTok{, }\StringTok{"p{-}value"}\NormalTok{))}
\end{Highlighting}
\end{Shaded}

\begin{longtable}[]{@{}lr@{}}
\caption{Shannon index: Oomycetes}\tabularnewline
\toprule\noalign{}
Microbiome & p-value \\
\midrule\noalign{}
\endfirsthead
\toprule\noalign{}
Microbiome & p-value \\
\midrule\noalign{}
\endhead
\bottomrule\noalign{}
\endlastfoot
minus & 0.6264941 \\
plus & 0.8589549 \\
\end{longtable}

\subsection{Plot Oomyctes}\label{plot-oomyctes}

\subsection{Alpha diversities}\label{alpha-diversities}

\begin{figure}

\centering{

\pandocbounded{\includegraphics[keepaspectratio]{05_Microbiome_analysis_files/figure-pdf/fig-SHANNON_plot-1.pdf}}

}

\caption{\label{fig-SHANNON_plot}Shannon indices.}

\end{figure}%

\section{Beta Diversity}\label{beta-diversity-1}

We get an overview by investigating the full model to see which factors
alters the beta diversity. We use the Bray-Curtis distance metric for
beta diversity analysis. We perform PERMANOVA to test the significance
of the differences in beta diversity between the different Arabidopsis
lines and conditions.

\subsection{Permanova}\label{permanova}

\begin{Shaded}
\begin{Highlighting}[]
\FunctionTok{set.seed}\NormalTok{(}\DecValTok{100010}\NormalTok{)}
\CommentTok{\#Bacteria}
\NormalTok{bdist }\OtherTok{\textless{}{-}} \FunctionTok{vegdist}\NormalTok{(}\FunctionTok{t}\NormalTok{(bCOUNT\_filt), }\AttributeTok{method=}\StringTok{"bray"}\NormalTok{)}
\NormalTok{bdist\_paov }\OtherTok{\textless{}{-}} \FunctionTok{adonis2}\NormalTok{(bdist }\SpecialCharTok{\textasciitilde{}}\NormalTok{ line }\SpecialCharTok{*}\NormalTok{ condition, }\AttributeTok{by =} \StringTok{"terms"}\NormalTok{, }\AttributeTok{data=}\NormalTok{bDESIGN\_filt)}
\CommentTok{\#pander(bdist\_paov, caption="PERMANOVA: Bacteria")}
\FunctionTok{kable}\NormalTok{(bdist\_paov, }\AttributeTok{caption =} \StringTok{"PERMANOVA: Bacteria"}\NormalTok{,}
      \AttributeTok{col.names =} \FunctionTok{c}\NormalTok{(}\StringTok{"Term"}\NormalTok{, }\StringTok{"DF"}\NormalTok{, }\StringTok{"Sums of Squares"}\NormalTok{, }\StringTok{"Mean Squares"}\NormalTok{, }\StringTok{"F{-}Model"}\NormalTok{, }\StringTok{"P{-}Value"}\NormalTok{))}
\end{Highlighting}
\end{Shaded}

\begin{longtable}[]{@{}lrrrrr@{}}
\caption{PERMANOVA: Bacteria}\tabularnewline
\toprule\noalign{}
Term & DF & Sums of Squares & Mean Squares & F-Model & P-Value \\
\midrule\noalign{}
\endfirsthead
\toprule\noalign{}
Term & DF & Sums of Squares & Mean Squares & F-Model & P-Value \\
\midrule\noalign{}
\endhead
\bottomrule\noalign{}
\endlastfoot
line & 1 & 0.1916056 & 0.0569100 & 2.310204 & 0.007 \\
condition & 1 & 0.1316189 & 0.0390930 & 1.586939 & 0.070 \\
line:condition & 1 & 0.1407344 & 0.0418004 & 1.696846 & 0.047 \\
Residual & 35 & 2.9028592 & 0.8621966 & NA & NA \\
Total & 38 & 3.3668181 & 1.0000000 & NA & NA \\
\end{longtable}

\begin{Shaded}
\begin{Highlighting}[]
\CommentTok{\#Fungi}
\NormalTok{fdist }\OtherTok{\textless{}{-}} \FunctionTok{vegdist}\NormalTok{(}\FunctionTok{t}\NormalTok{(fCOUNT\_filt), }\AttributeTok{method=}\StringTok{"bray"}\NormalTok{)}
\NormalTok{fdist\_paov }\OtherTok{\textless{}{-}} \FunctionTok{adonis2}\NormalTok{(fdist }\SpecialCharTok{\textasciitilde{}}\NormalTok{ line }\SpecialCharTok{*}\NormalTok{ condition, }\AttributeTok{by =} \StringTok{"terms"}\NormalTok{, }\AttributeTok{data=}\NormalTok{fDESIGN\_filt)}
\CommentTok{\#pander(fdist\_paov, caption="PERMANOVA: Fungi")}
\FunctionTok{kable}\NormalTok{(fdist\_paov, }\AttributeTok{caption =} \StringTok{"PERMANOVA: Fungi"}\NormalTok{,}
      \AttributeTok{col.names =} \FunctionTok{c}\NormalTok{(}\StringTok{"Term"}\NormalTok{, }\StringTok{"DF"}\NormalTok{, }\StringTok{"Sums of Squares"}\NormalTok{, }\StringTok{"Mean Squares"}\NormalTok{, }\StringTok{"F{-}Model"}\NormalTok{, }\StringTok{"P{-}Value"}\NormalTok{))}
\end{Highlighting}
\end{Shaded}

\begin{longtable}[]{@{}lrrrrr@{}}
\caption{PERMANOVA: Fungi}\tabularnewline
\toprule\noalign{}
Term & DF & Sums of Squares & Mean Squares & F-Model & P-Value \\
\midrule\noalign{}
\endfirsthead
\toprule\noalign{}
Term & DF & Sums of Squares & Mean Squares & F-Model & P-Value \\
\midrule\noalign{}
\endhead
\bottomrule\noalign{}
\endlastfoot
line & 1 & 0.4890016 & 0.0371969 & 1.116721 & 0.207 \\
condition & 1 & 0.4381551 & 0.0333291 & 1.000605 & 0.422 \\
line:condition & 1 & 0.3961077 & 0.0301307 & 0.904582 & 0.710 \\
Residual & 27 & 11.8230396 & 0.8993432 & NA & NA \\
Total & 30 & 13.1463040 & 1.0000000 & NA & NA \\
\end{longtable}

\begin{Shaded}
\begin{Highlighting}[]
\CommentTok{\#Bacteria}
\NormalTok{odist }\OtherTok{\textless{}{-}} \FunctionTok{vegdist}\NormalTok{(}\FunctionTok{t}\NormalTok{(oCOUNT\_filt), }\AttributeTok{method=}\StringTok{"bray"}\NormalTok{)}
\NormalTok{odist\_paov }\OtherTok{\textless{}{-}} \FunctionTok{adonis2}\NormalTok{(odist }\SpecialCharTok{\textasciitilde{}}\NormalTok{ line }\SpecialCharTok{*}\NormalTok{ condition, }\AttributeTok{by =} \StringTok{"terms"}\NormalTok{, }\AttributeTok{data=}\NormalTok{oDESIGN\_filt)}
\CommentTok{\#pander(odist\_paov, caption="PERMANOVA: Oomycetes")}
\FunctionTok{kable}\NormalTok{(odist\_paov, }\AttributeTok{caption =} \StringTok{"PERMANOVA: Oomycetes"}\NormalTok{,}
      \AttributeTok{col.names =} \FunctionTok{c}\NormalTok{(}\StringTok{"Term"}\NormalTok{, }\StringTok{"DF"}\NormalTok{, }\StringTok{"Sums of Squares"}\NormalTok{, }\StringTok{"Mean Squares"}\NormalTok{, }\StringTok{"F{-}Model"}\NormalTok{, }\StringTok{"P{-}Value"}\NormalTok{))}
\end{Highlighting}
\end{Shaded}

\begin{longtable}[]{@{}lrrrrr@{}}
\caption{PERMANOVA: Oomycetes}\tabularnewline
\toprule\noalign{}
Term & DF & Sums of Squares & Mean Squares & F-Model & P-Value \\
\midrule\noalign{}
\endfirsthead
\toprule\noalign{}
Term & DF & Sums of Squares & Mean Squares & F-Model & P-Value \\
\midrule\noalign{}
\endhead
\bottomrule\noalign{}
\endlastfoot
line & 1 & 0.5032322 & 0.0327083 & 1.190464 & 0.219 \\
condition & 1 & 0.8446513 & 0.0548994 & 1.998136 & 0.015 \\
line:condition & 1 & 0.5105272 & 0.0331825 & 1.207721 & 0.223 \\
Residual & 32 & 13.5270250 & 0.8792097 & NA & NA \\
Total & 35 & 15.3854357 & 1.0000000 & NA & NA \\
\end{longtable}

\textbf{Conclusion:} The bacterial communities are different between the
different Arabidopsis-lines, between the different conditions but not
their interaction. For fungi we find a difference between the
conditions, but not between lines or their interactions. For oomycetes
we only see community differences for the two conditions.

\subsection{Reduced Permanova}\label{reduced-permanova}

Compare Col-0 with mmf1 on the two soils. We use the Bray-Curtis
distance metric for beta diversity analysis. We perform PERMANOVA to
test the significance of the differences in beta diversity between Col-0
and \emph{mmf1-2} on the different soils.

\begin{Shaded}
\begin{Highlighting}[]
\CommentTok{\#Filter bacteria for only plus}
\NormalTok{bDESIGN\_plus }\OtherTok{\textless{}{-}} \FunctionTok{filter}\NormalTok{(bDESIGN\_filt, condition }\SpecialCharTok{\%in\%} \StringTok{"plus"}\NormalTok{)}
\NormalTok{bCOUNT\_t }\OtherTok{\textless{}{-}} \FunctionTok{t}\NormalTok{(bCOUNT)}
\NormalTok{bCOUNT\_plus }\OtherTok{\textless{}{-}}\NormalTok{ bCOUNT\_t[}\FunctionTok{rownames}\NormalTok{(bCOUNT\_t)}\SpecialCharTok{\%in\%}\NormalTok{bDESIGN\_plus}\SpecialCharTok{$}\NormalTok{sample\_name,]}
\NormalTok{bCOUNT\_plus }\OtherTok{\textless{}{-}} \FunctionTok{t}\NormalTok{(bCOUNT\_plus)}

\CommentTok{\#Filter bacteria for only minus}
\NormalTok{bDESIGN\_minus }\OtherTok{\textless{}{-}} \FunctionTok{filter}\NormalTok{(bDESIGN\_filt, condition }\SpecialCharTok{\%in\%} \StringTok{"minus"}\NormalTok{)}
\NormalTok{bCOUNT\_t }\OtherTok{\textless{}{-}} \FunctionTok{t}\NormalTok{(bCOUNT)}
\NormalTok{bCOUNT\_minus }\OtherTok{\textless{}{-}}\NormalTok{ bCOUNT\_t[}\FunctionTok{rownames}\NormalTok{(bCOUNT\_t)}\SpecialCharTok{\%in\%}\NormalTok{bDESIGN\_minus}\SpecialCharTok{$}\NormalTok{sample\_name,]}
\NormalTok{bCOUNT\_minus }\OtherTok{\textless{}{-}} \FunctionTok{t}\NormalTok{(bCOUNT\_minus)}

\CommentTok{\#Filter bacteria for only plus}
\NormalTok{fDESIGN\_plus }\OtherTok{\textless{}{-}} \FunctionTok{filter}\NormalTok{(fDESIGN\_filt, condition }\SpecialCharTok{\%in\%} \StringTok{"plus"}\NormalTok{)}
\NormalTok{fCOUNT\_t }\OtherTok{\textless{}{-}} \FunctionTok{t}\NormalTok{(fCOUNT)}
\NormalTok{fCOUNT\_plus }\OtherTok{\textless{}{-}}\NormalTok{ fCOUNT\_t[}\FunctionTok{rownames}\NormalTok{(fCOUNT\_t)}\SpecialCharTok{\%in\%}\NormalTok{fDESIGN\_plus}\SpecialCharTok{$}\NormalTok{sample\_name,]}
\NormalTok{fCOUNT\_plus }\OtherTok{\textless{}{-}} \FunctionTok{t}\NormalTok{(fCOUNT\_plus)}

\CommentTok{\#Filter bacteria for only minus}
\NormalTok{fDESIGN\_minus }\OtherTok{\textless{}{-}} \FunctionTok{filter}\NormalTok{(fDESIGN\_filt, condition }\SpecialCharTok{\%in\%} \StringTok{"minus"}\NormalTok{)}
\NormalTok{fCOUNT\_t }\OtherTok{\textless{}{-}} \FunctionTok{t}\NormalTok{(fCOUNT)}
\NormalTok{fCOUNT\_minus }\OtherTok{\textless{}{-}}\NormalTok{ fCOUNT\_t[}\FunctionTok{rownames}\NormalTok{(fCOUNT\_t)}\SpecialCharTok{\%in\%}\NormalTok{fDESIGN\_minus}\SpecialCharTok{$}\NormalTok{sample\_name,]}
\NormalTok{fCOUNT\_minus }\OtherTok{\textless{}{-}} \FunctionTok{t}\NormalTok{(fCOUNT\_minus)}

\CommentTok{\#Filter oomycetes for only plus}
\NormalTok{oDESIGN\_plus }\OtherTok{\textless{}{-}} \FunctionTok{filter}\NormalTok{(oDESIGN\_filt, condition }\SpecialCharTok{\%in\%} \StringTok{"plus"}\NormalTok{)}
\NormalTok{oCOUNT\_t }\OtherTok{\textless{}{-}} \FunctionTok{t}\NormalTok{(oCOUNT)}
\NormalTok{oCOUNT\_plus }\OtherTok{\textless{}{-}}\NormalTok{ oCOUNT\_t[}\FunctionTok{rownames}\NormalTok{(oCOUNT\_t)}\SpecialCharTok{\%in\%}\NormalTok{oDESIGN\_plus}\SpecialCharTok{$}\NormalTok{sample\_name,]}
\NormalTok{oCOUNT\_plus }\OtherTok{\textless{}{-}} \FunctionTok{t}\NormalTok{(oCOUNT\_plus)}

\CommentTok{\#Filter oomycetes for only minus}
\NormalTok{oDESIGN\_minus }\OtherTok{\textless{}{-}} \FunctionTok{filter}\NormalTok{(oDESIGN\_filt, condition }\SpecialCharTok{\%in\%} \StringTok{"minus"}\NormalTok{)}
\NormalTok{oCOUNT\_t }\OtherTok{\textless{}{-}} \FunctionTok{t}\NormalTok{(oCOUNT)}
\NormalTok{oCOUNT\_minus }\OtherTok{\textless{}{-}}\NormalTok{ oCOUNT\_t[}\FunctionTok{rownames}\NormalTok{(oCOUNT\_t)}\SpecialCharTok{\%in\%}\NormalTok{oDESIGN\_minus}\SpecialCharTok{$}\NormalTok{sample\_name,]}
\NormalTok{oCOUNT\_minus }\OtherTok{\textless{}{-}} \FunctionTok{t}\NormalTok{(oCOUNT\_minus)}
\end{Highlighting}
\end{Shaded}

\subsubsection{\texorpdfstring{Beta diversity between Col-0 and
\emph{mmf1-2}}{Beta diversity between Col-0 and mmf1-2}}\label{beta-diversity-between-col-0-and-mmf1-2}

We get an overview by investigating the full model to see which factors
alters the beta diversity

\begin{Shaded}
\begin{Highlighting}[]
\CommentTok{\#normalise to 100\&}
\CommentTok{\# bCOUNT \textless{}{-} hilldiv::tss(bCOUNT)*100}
\CommentTok{\# fCOUNT \textless{}{-} hilldiv::tss(fCOUNT)*100}
\CommentTok{\# oCOUNT \textless{}{-} hilldiv::tss(oCOUNT)*100}
\FunctionTok{set.seed}\NormalTok{(}\DecValTok{10001}\NormalTok{)}
\CommentTok{\#Bacteria on plus}
\NormalTok{bdist\_plus }\OtherTok{\textless{}{-}} \FunctionTok{vegdist}\NormalTok{(}\FunctionTok{t}\NormalTok{(bCOUNT\_plus), }\AttributeTok{method=}\StringTok{"bray"}\NormalTok{)}
\NormalTok{bdist\_paov\_plus }\OtherTok{\textless{}{-}} \FunctionTok{adonis2}\NormalTok{(bdist\_plus }\SpecialCharTok{\textasciitilde{}}\NormalTok{ line,}\AttributeTok{by =} \StringTok{"terms"}\NormalTok{, }\AttributeTok{data=}\NormalTok{bDESIGN\_plus)}
\CommentTok{\#pander(bdist\_paov\_plus, caption="PERMANOVA: Bacteria on BXplus")}
\FunctionTok{kable}\NormalTok{(bdist\_paov\_plus, }\AttributeTok{caption =} \StringTok{"PERMANOVA: Bacteria on BXplus"}\NormalTok{,}
      \AttributeTok{col.names =} \FunctionTok{c}\NormalTok{(}\StringTok{"Term"}\NormalTok{, }\StringTok{"DF"}\NormalTok{, }\StringTok{"Sums of Squares"}\NormalTok{, }\StringTok{"Mean Squares"}\NormalTok{, }\StringTok{"F{-}Model"}\NormalTok{, }\StringTok{"P{-}Value"}\NormalTok{))}
\end{Highlighting}
\end{Shaded}

\begin{longtable}[]{@{}lrrrrr@{}}
\caption{PERMANOVA: Bacteria on BXplus}\tabularnewline
\toprule\noalign{}
Term & DF & Sums of Squares & Mean Squares & F-Model & P-Value \\
\midrule\noalign{}
\endfirsthead
\toprule\noalign{}
Term & DF & Sums of Squares & Mean Squares & F-Model & P-Value \\
\midrule\noalign{}
\endhead
\bottomrule\noalign{}
\endlastfoot
line & 1 & 0.2606093 & 0.1582689 & 3.384502 & 0.001 \\
Residual & 18 & 1.3860141 & 0.8417311 & NA & NA \\
Total & 19 & 1.6466235 & 1.0000000 & NA & NA \\
\end{longtable}

\begin{Shaded}
\begin{Highlighting}[]
\CommentTok{\#Bacteria on minus}
\NormalTok{bdist\_minus }\OtherTok{\textless{}{-}} \FunctionTok{vegdist}\NormalTok{(}\FunctionTok{t}\NormalTok{(bCOUNT\_minus), }\AttributeTok{method=}\StringTok{"bray"}\NormalTok{)}
\NormalTok{bdist\_paov\_minus }\OtherTok{\textless{}{-}} \FunctionTok{adonis2}\NormalTok{(bdist\_minus }\SpecialCharTok{\textasciitilde{}}\NormalTok{ line,}\AttributeTok{by =} \StringTok{"terms"}\NormalTok{, }\AttributeTok{data=}\NormalTok{bDESIGN\_minus)}
\CommentTok{\#pander(bdist\_paov\_minus, caption="PERMANOVA: Bacteria on BXminus")}
\FunctionTok{kable}\NormalTok{(bdist\_paov\_minus, }\AttributeTok{caption =} \StringTok{"PERMANOVA: Bacteria on BXminus"}\NormalTok{,}
      \AttributeTok{col.names =} \FunctionTok{c}\NormalTok{(}\StringTok{"Term"}\NormalTok{, }\StringTok{"DF"}\NormalTok{, }\StringTok{"Sums of Squares"}\NormalTok{, }\StringTok{"Mean Squares"}\NormalTok{, }\StringTok{"F{-}Model"}\NormalTok{, }\StringTok{"P{-}Value"}\NormalTok{))}
\end{Highlighting}
\end{Shaded}

\begin{longtable}[]{@{}lrrrrr@{}}
\caption{PERMANOVA: Bacteria on BXminus}\tabularnewline
\toprule\noalign{}
Term & DF & Sums of Squares & Mean Squares & F-Model & P-Value \\
\midrule\noalign{}
\endfirsthead
\toprule\noalign{}
Term & DF & Sums of Squares & Mean Squares & F-Model & P-Value \\
\midrule\noalign{}
\endhead
\bottomrule\noalign{}
\endlastfoot
line & 1 & 0.0725056 & 0.0456196 & 0.8126044 & 0.636 \\
Residual & 17 & 1.5168450 & 0.9543804 & NA & NA \\
Total & 18 & 1.5893506 & 1.0000000 & NA & NA \\
\end{longtable}

\begin{Shaded}
\begin{Highlighting}[]
\CommentTok{\#Fungi on plus}
\NormalTok{fdist\_plus }\OtherTok{\textless{}{-}} \FunctionTok{vegdist}\NormalTok{(}\FunctionTok{t}\NormalTok{(fCOUNT\_plus), }\AttributeTok{method=}\StringTok{"bray"}\NormalTok{)}
\NormalTok{fdist\_paov\_plus }\OtherTok{\textless{}{-}} \FunctionTok{adonis2}\NormalTok{(fdist\_plus }\SpecialCharTok{\textasciitilde{}}\NormalTok{ line,}\AttributeTok{by =} \StringTok{"terms"}\NormalTok{,}\AttributeTok{data=}\NormalTok{fDESIGN\_plus)}
\CommentTok{\#pander(fdist\_paov\_plus, caption="PERMANOVA: Fungi on BXplus")}
\FunctionTok{kable}\NormalTok{(fdist\_paov\_plus, }\AttributeTok{caption =} \StringTok{"PERMANOVA: Fungi on BXplus"}\NormalTok{,}
      \AttributeTok{col.names =} \FunctionTok{c}\NormalTok{(}\StringTok{"Term"}\NormalTok{, }\StringTok{"DF"}\NormalTok{, }\StringTok{"Sums of Squares"}\NormalTok{, }\StringTok{"Mean Squares"}\NormalTok{, }\StringTok{"F{-}Model"}\NormalTok{, }\StringTok{"P{-}Value"}\NormalTok{))}
\end{Highlighting}
\end{Shaded}

\begin{longtable}[]{@{}lrrrrr@{}}
\caption{PERMANOVA: Fungi on BXplus}\tabularnewline
\toprule\noalign{}
Term & DF & Sums of Squares & Mean Squares & F-Model & P-Value \\
\midrule\noalign{}
\endfirsthead
\toprule\noalign{}
Term & DF & Sums of Squares & Mean Squares & F-Model & P-Value \\
\midrule\noalign{}
\endhead
\bottomrule\noalign{}
\endlastfoot
line & 1 & 0.4521943 & 0.0675697 & 1.014527 & 0.455 \\
Residual & 14 & 6.2400688 & 0.9324303 & NA & NA \\
Total & 15 & 6.6922631 & 1.0000000 & NA & NA \\
\end{longtable}

\begin{Shaded}
\begin{Highlighting}[]
\CommentTok{\#Fungi on minus}
\NormalTok{fdist\_minus }\OtherTok{\textless{}{-}} \FunctionTok{vegdist}\NormalTok{(}\FunctionTok{t}\NormalTok{(fCOUNT\_minus), }\AttributeTok{method=}\StringTok{"bray"}\NormalTok{)}
\NormalTok{fdist\_paov\_minus }\OtherTok{\textless{}{-}} \FunctionTok{adonis2}\NormalTok{(fdist\_minus }\SpecialCharTok{\textasciitilde{}}\NormalTok{ line, }\AttributeTok{by =} \StringTok{"terms"}\NormalTok{,}\AttributeTok{data=}\NormalTok{fDESIGN\_minus)}
\CommentTok{\#pander(fdist\_paov\_minus, caption="PERMANOVA: Fungi on BXminus")}
\FunctionTok{kable}\NormalTok{(fdist\_paov\_minus, }\AttributeTok{caption =} \StringTok{"PERMANOVA: Fungi on BXminus"}\NormalTok{,}
      \AttributeTok{col.names =} \FunctionTok{c}\NormalTok{(}\StringTok{"Term"}\NormalTok{, }\StringTok{"DF"}\NormalTok{, }\StringTok{"Sums of Squares"}\NormalTok{, }\StringTok{"Mean Squares"}\NormalTok{, }\StringTok{"F{-}Model"}\NormalTok{, }\StringTok{"P{-}Value"}\NormalTok{))}
\end{Highlighting}
\end{Shaded}

\begin{longtable}[]{@{}lrrrrr@{}}
\caption{PERMANOVA: Fungi on BXminus}\tabularnewline
\toprule\noalign{}
Term & DF & Sums of Squares & Mean Squares & F-Model & P-Value \\
\midrule\noalign{}
\endfirsthead
\toprule\noalign{}
Term & DF & Sums of Squares & Mean Squares & F-Model & P-Value \\
\midrule\noalign{}
\endhead
\bottomrule\noalign{}
\endlastfoot
line & 1 & 0.4345633 & 0.0722162 & 1.011885 & 0.43 \\
Residual & 13 & 5.5829708 & 0.9277838 & NA & NA \\
Total & 14 & 6.0175341 & 1.0000000 & NA & NA \\
\end{longtable}

\begin{Shaded}
\begin{Highlighting}[]
\CommentTok{\#Oomycetes on plus}
\NormalTok{odist\_plus }\OtherTok{\textless{}{-}} \FunctionTok{vegdist}\NormalTok{(}\FunctionTok{t}\NormalTok{(oCOUNT\_plus), }\AttributeTok{method=}\StringTok{"bray"}\NormalTok{)}
\NormalTok{odist\_paov\_plus }\OtherTok{\textless{}{-}} \FunctionTok{adonis2}\NormalTok{(odist\_plus }\SpecialCharTok{\textasciitilde{}}\NormalTok{ line, }\AttributeTok{by =} \StringTok{"terms"}\NormalTok{, }\AttributeTok{data=}\NormalTok{oDESIGN\_plus)}
\CommentTok{\#pander(odist\_paov\_plus, caption="PERMANOVA: Oomycetes on BXplus")}
\FunctionTok{kable}\NormalTok{(odist\_paov\_plus, }\AttributeTok{caption =} \StringTok{"PERMANOVA: Oomycetes on BXplus"}\NormalTok{,}
      \AttributeTok{col.names =} \FunctionTok{c}\NormalTok{(}\StringTok{"Term"}\NormalTok{, }\StringTok{"DF"}\NormalTok{, }\StringTok{"Sums of Squares"}\NormalTok{, }\StringTok{"Mean Squares"}\NormalTok{, }\StringTok{"F{-}Model"}\NormalTok{, }\StringTok{"P{-}Value"}\NormalTok{))}
\end{Highlighting}
\end{Shaded}

\begin{longtable}[]{@{}lrrrrr@{}}
\caption{PERMANOVA: Oomycetes on BXplus}\tabularnewline
\toprule\noalign{}
Term & DF & Sums of Squares & Mean Squares & F-Model & P-Value \\
\midrule\noalign{}
\endfirsthead
\toprule\noalign{}
Term & DF & Sums of Squares & Mean Squares & F-Model & P-Value \\
\midrule\noalign{}
\endhead
\bottomrule\noalign{}
\endlastfoot
line & 1 & 0.4362029 & 0.0626403 & 1.069221 & 0.348 \\
Residual & 16 & 6.5274127 & 0.9373597 & NA & NA \\
Total & 17 & 6.9636156 & 1.0000000 & NA & NA \\
\end{longtable}

\begin{Shaded}
\begin{Highlighting}[]
\CommentTok{\#Oomycetes on minus}
\NormalTok{odist\_minus }\OtherTok{\textless{}{-}} \FunctionTok{vegdist}\NormalTok{(}\FunctionTok{t}\NormalTok{(oCOUNT\_minus), }\AttributeTok{method=}\StringTok{"bray"}\NormalTok{)}
\NormalTok{odist\_paov\_minus }\OtherTok{\textless{}{-}} \FunctionTok{adonis2}\NormalTok{(odist\_minus }\SpecialCharTok{\textasciitilde{}}\NormalTok{ line, }\AttributeTok{by =} \StringTok{"terms"}\NormalTok{, }\AttributeTok{data=}\NormalTok{oDESIGN\_minus)}
\CommentTok{\#pander(odist\_paov\_minus, caption="PERMANOVA: Oomycetes on BXminus")}
\FunctionTok{kable}\NormalTok{(odist\_paov\_minus, }\AttributeTok{caption =} \StringTok{"PERMANOVA: Oomycetes on BXminus"}\NormalTok{,}
      \AttributeTok{col.names =} \FunctionTok{c}\NormalTok{(}\StringTok{"Term"}\NormalTok{, }\StringTok{"DF"}\NormalTok{, }\StringTok{"Sums of Squares"}\NormalTok{, }\StringTok{"Mean Squares"}\NormalTok{, }\StringTok{"F{-}Model"}\NormalTok{, }\StringTok{"P{-}Value"}\NormalTok{))}
\end{Highlighting}
\end{Shaded}

\begin{longtable}[]{@{}lrrrrr@{}}
\caption{PERMANOVA: Oomycetes on BXminus}\tabularnewline
\toprule\noalign{}
Term & DF & Sums of Squares & Mean Squares & F-Model & P-Value \\
\midrule\noalign{}
\endfirsthead
\toprule\noalign{}
Term & DF & Sums of Squares & Mean Squares & F-Model & P-Value \\
\midrule\noalign{}
\endhead
\bottomrule\noalign{}
\endlastfoot
line & 1 & 0.5667319 & 0.0749017 & 1.295459 & 0.11 \\
Residual & 16 & 6.9996123 & 0.9250983 & NA & NA \\
Total & 17 & 7.5663442 & 1.0000000 & NA & NA \\
\end{longtable}

\textbf{Conclusion:} The bacterial communities are different between
Col-0 and \emph{mmf1-2} on the BXplus soil, but not on the BXminus soil.
The fungal and oomycte communitites are not different between Col-0 and
\emph{mmf1-2} on the BXplus or BXminus soils.

\section{\texorpdfstring{\textbf{CAP - line and condition effect on beta
diversity}}{CAP - line and condition effect on beta diversity}}\label{cap---line-and-condition-effect-on-beta-diversity}

\begin{Shaded}
\begin{Highlighting}[]
\CommentTok{\#colors}
\NormalTok{bDESIGN}\SpecialCharTok{$}\NormalTok{cols\_line }\OtherTok{\textless{}{-}} \ConstantTok{NA}
\NormalTok{bDESIGN}\SpecialCharTok{$}\NormalTok{cols\_line[bDESIGN}\SpecialCharTok{$}\NormalTok{line}\SpecialCharTok{==}\StringTok{"WT"}\NormalTok{] }\OtherTok{\textless{}{-}} \StringTok{"deepskyblue4"}
\NormalTok{bDESIGN}\SpecialCharTok{$}\NormalTok{cols\_line[bDESIGN}\SpecialCharTok{$}\NormalTok{line}\SpecialCharTok{==}\StringTok{"159"}\NormalTok{] }\OtherTok{\textless{}{-}} \StringTok{"\#33BBEE"}
\NormalTok{bDESIGN}\SpecialCharTok{$}\NormalTok{cols\_line[bDESIGN}\SpecialCharTok{$}\NormalTok{line}\SpecialCharTok{==}\StringTok{"318"}\NormalTok{] }\OtherTok{\textless{}{-}} \StringTok{"orange2"}

\NormalTok{temp }\OtherTok{\textless{}{-}} \FunctionTok{data.frame}\NormalTok{(bDESIGN}\SpecialCharTok{$}\NormalTok{line, bDESIGN}\SpecialCharTok{$}\NormalTok{cols\_line)  }\CommentTok{\#collaps color vector}
\NormalTok{temp }\OtherTok{\textless{}{-}}\NormalTok{ plyr}\SpecialCharTok{::}\FunctionTok{ddply}\NormalTok{(temp, }\AttributeTok{.variables=}\StringTok{"bDESIGN.cols\_line"}\NormalTok{, }\AttributeTok{.fun=}\NormalTok{unique)}
\NormalTok{level\_cols\_line }\OtherTok{\textless{}{-}} \FunctionTok{as.character}\NormalTok{(temp[,}\DecValTok{2}\NormalTok{])}
\FunctionTok{names}\NormalTok{(level\_cols\_line) }\OtherTok{\textless{}{-}}\NormalTok{ temp[,}\DecValTok{1}\NormalTok{]}
\FunctionTok{rm}\NormalTok{(temp)}


\NormalTok{bDESIGN}\SpecialCharTok{$}\NormalTok{cols\_condition }\OtherTok{\textless{}{-}} \ConstantTok{NA}
\NormalTok{bDESIGN}\SpecialCharTok{$}\NormalTok{cols\_condition[bDESIGN}\SpecialCharTok{$}\NormalTok{condition}\SpecialCharTok{==}\StringTok{"minus"}\NormalTok{] }\OtherTok{\textless{}{-}} \StringTok{"seagreen"}
\NormalTok{bDESIGN}\SpecialCharTok{$}\NormalTok{cols\_condition[bDESIGN}\SpecialCharTok{$}\NormalTok{condition}\SpecialCharTok{==}\StringTok{"plus"}\NormalTok{] }\OtherTok{\textless{}{-}} \StringTok{"gold2"}
\NormalTok{temp }\OtherTok{\textless{}{-}} \FunctionTok{data.frame}\NormalTok{(bDESIGN}\SpecialCharTok{$}\NormalTok{condition, bDESIGN}\SpecialCharTok{$}\NormalTok{cols\_condition)}
\NormalTok{temp }\OtherTok{\textless{}{-}}\NormalTok{ plyr}\SpecialCharTok{::}\FunctionTok{ddply}\NormalTok{(temp, }\AttributeTok{.variables=}\StringTok{"bDESIGN.cols\_condition"}\NormalTok{, }\AttributeTok{.fun=}\NormalTok{unique)}
\NormalTok{level\_cols\_condition }\OtherTok{\textless{}{-}} \FunctionTok{as.character}\NormalTok{(temp[,}\DecValTok{2}\NormalTok{])}
\FunctionTok{names}\NormalTok{(level\_cols\_condition) }\OtherTok{\textless{}{-}}\NormalTok{ temp[,}\DecValTok{1}\NormalTok{]}
\end{Highlighting}
\end{Shaded}

\subsection{Prepare figure}\label{prepare-figure}

We use the Constrained Analysis of Principal coordinates (CAP) to
analyze the effects of line and condition on beta diversity for
bacterial, fungal, and oomycete communities. We use the Bray-Curtis
distance metric for beta diversity analysis.

\begin{Shaded}
\begin{Highlighting}[]
\FunctionTok{library}\NormalTok{(tidyverse, }\AttributeTok{quietly =} \ConstantTok{TRUE}\NormalTok{)}

\DocumentationTok{\#\#\# Figure 5.2 | CAP {-} line and condition effect on beta diversity}

\CommentTok{\# This script performs Constrained Analysis of Principal coordinates (CAP)}
\CommentTok{\# on bacterial, fungal, and oomycete communities, analyzing the effects of }
\CommentTok{\# line and condition on beta diversity.}

\CommentTok{\# Set plot dimensions}
\NormalTok{fig.height }\OtherTok{\textless{}{-}} \DecValTok{6}
\NormalTok{fig.width }\OtherTok{\textless{}{-}} \DecValTok{12}

\CommentTok{\# ==============================}
\CommentTok{\# Bacteria Analysis {-}{-}{-}{-}{-}{-}{-}{-}{-}{-}{-}{-}{-}{-}{-}{-}{-}{-}{-}{-}{-}{-}{-}{-}{-}{-}{-}}
\CommentTok{\# ==============================}

\CommentTok{\# CAP ordination for Bacteria}
\NormalTok{CAP\_bray }\OtherTok{\textless{}{-}} \FunctionTok{ordinate}\NormalTok{(bPHYSEQ, }\AttributeTok{method =} \StringTok{"CAP"}\NormalTok{, }\AttributeTok{formula =} \SpecialCharTok{\textasciitilde{}}\NormalTok{ line }\SpecialCharTok{*}\NormalTok{ condition, }\AttributeTok{distance =} \StringTok{"bray"}\NormalTok{)}
\NormalTok{CAP\_bray\_var\_tbl }\OtherTok{\textless{}{-}} \FunctionTok{variability\_table}\NormalTok{(CAP\_bray)  }\CommentTok{\# Extract variation details}

\CommentTok{\# Summarize abundance for the top 6 phyla}
\NormalTok{phylum.sum }\OtherTok{\textless{}{-}} \FunctionTok{tapply}\NormalTok{(}
  \FunctionTok{taxa\_sums}\NormalTok{(bPHYSEQ), }
  \FunctionTok{tax\_table}\NormalTok{(bPHYSEQ)[, }\StringTok{"labels"}\NormalTok{], }
\NormalTok{  sum, }
  \AttributeTok{na.rm =} \ConstantTok{TRUE}
\NormalTok{)}
\NormalTok{topphyla }\OtherTok{\textless{}{-}} \FunctionTok{names}\NormalTok{(}\FunctionTok{sort}\NormalTok{(phylum.sum, }\AttributeTok{decreasing =} \ConstantTok{TRUE}\NormalTok{))[}\DecValTok{1}\SpecialCharTok{:}\DecValTok{6}\NormalTok{]}
\NormalTok{phy\_sub }\OtherTok{\textless{}{-}} \FunctionTok{prune\_taxa}\NormalTok{(}
  \FunctionTok{tax\_table}\NormalTok{(bPHYSEQ)[, }\StringTok{"labels"}\NormalTok{] }\SpecialCharTok{\%in\%}\NormalTok{ topphyla, }
\NormalTok{  bPHYSEQ}
\NormalTok{)}

\CommentTok{\# Extract R2 and p{-}values for the subtitle}
\NormalTok{score }\OtherTok{\textless{}{-}} \FunctionTok{paste}\NormalTok{(}
  \StringTok{" \textasciitilde{} line * condition }\SpecialCharTok{\textbackslash{}n}\StringTok{ line: "}\NormalTok{,}
  \FunctionTok{format}\NormalTok{(bdist\_paov}\SpecialCharTok{$}\NormalTok{R2[}\FunctionTok{rownames}\NormalTok{(bdist\_paov) }\SpecialCharTok{==} \StringTok{"line"}\NormalTok{] }\SpecialCharTok{*} \DecValTok{100}\NormalTok{, }\AttributeTok{digits =} \DecValTok{2}\NormalTok{, }\AttributeTok{nsmall =} \DecValTok{1}\NormalTok{),}
  \StringTok{"\% var, P="}\NormalTok{, }\FunctionTok{format}\NormalTok{(bdist\_paov}\SpecialCharTok{$}\StringTok{\textasciigrave{}}\AttributeTok{Pr(\textgreater{}F)}\StringTok{\textasciigrave{}}\NormalTok{[}\FunctionTok{rownames}\NormalTok{(bdist\_paov) }\SpecialCharTok{==} \StringTok{"line"}\NormalTok{], }\AttributeTok{digits =} \DecValTok{2}\NormalTok{),}
  \StringTok{"}\SpecialCharTok{\textbackslash{}n}\StringTok{ condition: "}\NormalTok{,}
  \FunctionTok{format}\NormalTok{(bdist\_paov}\SpecialCharTok{$}\NormalTok{R2[}\FunctionTok{rownames}\NormalTok{(bdist\_paov) }\SpecialCharTok{==} \StringTok{"condition"}\NormalTok{] }\SpecialCharTok{*} \DecValTok{100}\NormalTok{, }\AttributeTok{digits =} \DecValTok{2}\NormalTok{, }\AttributeTok{nsmall =} \DecValTok{1}\NormalTok{),}
  \StringTok{"\% var, P="}\NormalTok{, }\FunctionTok{format}\NormalTok{(bdist\_paov}\SpecialCharTok{$}\StringTok{\textasciigrave{}}\AttributeTok{Pr(\textgreater{}F)}\StringTok{\textasciigrave{}}\NormalTok{[}\FunctionTok{rownames}\NormalTok{(bdist\_paov) }\SpecialCharTok{==} \StringTok{"condition"}\NormalTok{], }\AttributeTok{digits =} \DecValTok{2}\NormalTok{),}
  \StringTok{"}\SpecialCharTok{\textbackslash{}n}\StringTok{ line:condition: "}\NormalTok{,}
  \FunctionTok{format}\NormalTok{(bdist\_paov}\SpecialCharTok{$}\NormalTok{R2[}\FunctionTok{rownames}\NormalTok{(bdist\_paov) }\SpecialCharTok{==} \StringTok{"line:condition"}\NormalTok{] }\SpecialCharTok{*} \DecValTok{100}\NormalTok{, }\AttributeTok{digits =} \DecValTok{2}\NormalTok{, }\AttributeTok{nsmall =} \DecValTok{1}\NormalTok{),}
  \StringTok{"\% var, P="}\NormalTok{, }\FunctionTok{format}\NormalTok{(bdist\_paov}\SpecialCharTok{$}\StringTok{\textasciigrave{}}\AttributeTok{Pr(\textgreater{}F)}\StringTok{\textasciigrave{}}\NormalTok{[}\FunctionTok{rownames}\NormalTok{(bdist\_paov) }\SpecialCharTok{==} \StringTok{"line:condition"}\NormalTok{], }\AttributeTok{digits =} \DecValTok{2}\NormalTok{)}
\NormalTok{)}

\CommentTok{\# Plot ordination for Bacteria}
\NormalTok{p0 }\OtherTok{\textless{}{-}}\NormalTok{ phyloseq}\SpecialCharTok{::}\FunctionTok{plot\_ordination}\NormalTok{(phy\_sub, CAP\_bray, }\AttributeTok{title =} \StringTok{"Bacteria"}\NormalTok{, }\AttributeTok{axes =} \FunctionTok{c}\NormalTok{(}\DecValTok{1}\NormalTok{, }\DecValTok{2}\NormalTok{))}
\NormalTok{p0}\SpecialCharTok{$}\NormalTok{layers }\OtherTok{\textless{}{-}}\NormalTok{ p0}\SpecialCharTok{$}\NormalTok{layers[}\SpecialCharTok{{-}}\DecValTok{1}\NormalTok{]  }\CommentTok{\# Remove default layer (points)}
\NormalTok{p0 }\OtherTok{\textless{}{-}}\NormalTok{ p0 }\SpecialCharTok{+}
  \FunctionTok{stat\_ellipse}\NormalTok{(}
    \AttributeTok{geom =} \StringTok{"polygon"}\NormalTok{, }\AttributeTok{type =} \StringTok{"t"}\NormalTok{, }\AttributeTok{level =} \FloatTok{0.95}\NormalTok{, }\AttributeTok{size =} \FloatTok{0.4}\NormalTok{,}
    \FunctionTok{aes}\NormalTok{(}\AttributeTok{group =}\NormalTok{ condition, }\AttributeTok{color =}\NormalTok{ condition, }\AttributeTok{fill =}\NormalTok{ condition), }
    \AttributeTok{alpha =} \FloatTok{0.01}
\NormalTok{  ) }\SpecialCharTok{+}
  \FunctionTok{scale\_fill\_manual}\NormalTok{(}\AttributeTok{values =}\NormalTok{ level\_cols\_condition) }\SpecialCharTok{+}
  \FunctionTok{scale\_color\_manual}\NormalTok{(}\AttributeTok{values =}\NormalTok{ level\_cols\_condition) }\SpecialCharTok{+}
  \FunctionTok{new\_scale\_color}\NormalTok{() }\SpecialCharTok{+}
  \FunctionTok{geom\_point}\NormalTok{(}\FunctionTok{aes}\NormalTok{(}\AttributeTok{shape =} \FunctionTok{factor}\NormalTok{(line), }\AttributeTok{color =}\NormalTok{ condition), }\AttributeTok{size =} \FloatTok{1.0}\NormalTok{) }\SpecialCharTok{+}
  \FunctionTok{scale\_shape\_manual}\NormalTok{(}\AttributeTok{values =} \FunctionTok{c}\NormalTok{(}\DecValTok{16}\NormalTok{, }\DecValTok{2}\NormalTok{)) }\SpecialCharTok{+}
  \FunctionTok{scale\_color\_manual}\NormalTok{(}\AttributeTok{values =}\NormalTok{ level\_cols\_condition) }\SpecialCharTok{+}
  \FunctionTok{labs}\NormalTok{(}
    \AttributeTok{subtitle =}\NormalTok{ score,}
    \AttributeTok{color =} \StringTok{"Microbiome"}\NormalTok{, }\AttributeTok{fill =} \StringTok{"Microbiome"}\NormalTok{, }
    \AttributeTok{shape =} \StringTok{"Line"}\NormalTok{, }\AttributeTok{alpha =} \StringTok{"Line"}
\NormalTok{  ) }\SpecialCharTok{+}
  \FunctionTok{theme\_bw}\NormalTok{() }\SpecialCharTok{+}
  \FunctionTok{theme}\NormalTok{(}
    \AttributeTok{text =} \FunctionTok{element\_text}\NormalTok{(}\AttributeTok{size =} \DecValTok{7}\NormalTok{),             }\CommentTok{\# Set all text to 7pt}
    \AttributeTok{axis.title =} \FunctionTok{element\_text}\NormalTok{(}\AttributeTok{size =} \DecValTok{7}\NormalTok{),       }\CommentTok{\# Axis titles}
    \AttributeTok{axis.text =} \FunctionTok{element\_text}\NormalTok{(}\AttributeTok{size =} \DecValTok{7}\NormalTok{),        }\CommentTok{\# Axis text}
    \AttributeTok{legend.title =} \FunctionTok{element\_text}\NormalTok{(}\AttributeTok{size =} \DecValTok{7}\NormalTok{),     }\CommentTok{\# Legend titles}
    \AttributeTok{legend.text =} \FunctionTok{element\_text}\NormalTok{(}\AttributeTok{size =} \DecValTok{7}\NormalTok{),      }\CommentTok{\# Legend text}
    \AttributeTok{plot.title =} \FunctionTok{element\_text}\NormalTok{(}\AttributeTok{size =} \DecValTok{7}\NormalTok{),       }\CommentTok{\# Plot title}
    \AttributeTok{plot.subtitle =} \FunctionTok{element\_text}\NormalTok{(}\AttributeTok{size =} \DecValTok{7}\NormalTok{)     }\CommentTok{\# Plot subtitle}
\NormalTok{  )}
\NormalTok{fig\_bac }\OtherTok{\textless{}{-}}\NormalTok{ p0  }\CommentTok{\# Save the plot}

\CommentTok{\# ==============================}
\CommentTok{\# Fungi and Oomycetes Analyses {-}{-}{-}{-}{-}{-}{-}{-}{-}{-}{-}{-}{-}{-}{-}{-}{-}{-}{-}{-}{-}{-}{-}{-}{-}{-}{-}}
\CommentTok{\# ==============================}

\CommentTok{\# Repeat similar steps for Fungi and Oomycetes}
\NormalTok{process\_CAP }\OtherTok{\textless{}{-}} \ControlFlowTok{function}\NormalTok{(PHYSEQ, dist\_paov, title, colors) \{}
\NormalTok{  CAP\_bray }\OtherTok{\textless{}{-}} \FunctionTok{ordinate}\NormalTok{(PHYSEQ, }\AttributeTok{method =} \StringTok{"CAP"}\NormalTok{, }\AttributeTok{formula =} \SpecialCharTok{\textasciitilde{}}\NormalTok{ line }\SpecialCharTok{*}\NormalTok{ condition, }\AttributeTok{distance =} \StringTok{"bray"}\NormalTok{)}
\NormalTok{  phylum.sum }\OtherTok{\textless{}{-}} \FunctionTok{tapply}\NormalTok{(}
    \FunctionTok{taxa\_sums}\NormalTok{(PHYSEQ), }
    \FunctionTok{tax\_table}\NormalTok{(PHYSEQ)[, }\StringTok{"labels"}\NormalTok{], }
\NormalTok{    sum, }
    \AttributeTok{na.rm =} \ConstantTok{TRUE}
\NormalTok{  )}
\NormalTok{  topphyla }\OtherTok{\textless{}{-}} \FunctionTok{names}\NormalTok{(}\FunctionTok{sort}\NormalTok{(phylum.sum, }\AttributeTok{decreasing =} \ConstantTok{TRUE}\NormalTok{))[}\DecValTok{1}\SpecialCharTok{:}\DecValTok{6}\NormalTok{]}
\NormalTok{  phy\_sub }\OtherTok{\textless{}{-}} \FunctionTok{prune\_taxa}\NormalTok{(}\FunctionTok{tax\_table}\NormalTok{(PHYSEQ)[, }\StringTok{"labels"}\NormalTok{] }\SpecialCharTok{\%in\%}\NormalTok{ topphyla, PHYSEQ)}
\NormalTok{  score }\OtherTok{\textless{}{-}} \FunctionTok{paste}\NormalTok{(}
    \StringTok{" \textasciitilde{} line * condition }\SpecialCharTok{\textbackslash{}n}\StringTok{ line: "}\NormalTok{,}
    \FunctionTok{format}\NormalTok{(dist\_paov}\SpecialCharTok{$}\NormalTok{R2[}\FunctionTok{rownames}\NormalTok{(dist\_paov) }\SpecialCharTok{==} \StringTok{"line"}\NormalTok{] }\SpecialCharTok{*} \DecValTok{100}\NormalTok{, }\AttributeTok{digits =} \DecValTok{2}\NormalTok{, }\AttributeTok{nsmall =} \DecValTok{1}\NormalTok{),}
    \StringTok{"\% var, P="}\NormalTok{, }\FunctionTok{format}\NormalTok{(dist\_paov}\SpecialCharTok{$}\StringTok{\textasciigrave{}}\AttributeTok{Pr(\textgreater{}F)}\StringTok{\textasciigrave{}}\NormalTok{[}\FunctionTok{rownames}\NormalTok{(dist\_paov) }\SpecialCharTok{==} \StringTok{"line"}\NormalTok{], }\AttributeTok{digits =} \DecValTok{2}\NormalTok{),}
    \StringTok{"}\SpecialCharTok{\textbackslash{}n}\StringTok{ condition: "}\NormalTok{,}
    \FunctionTok{format}\NormalTok{(dist\_paov}\SpecialCharTok{$}\NormalTok{R2[}\FunctionTok{rownames}\NormalTok{(dist\_paov) }\SpecialCharTok{==} \StringTok{"condition"}\NormalTok{] }\SpecialCharTok{*} \DecValTok{100}\NormalTok{, }\AttributeTok{digits =} \DecValTok{2}\NormalTok{, }\AttributeTok{nsmall =} \DecValTok{1}\NormalTok{),}
    \StringTok{"\% var, P="}\NormalTok{, }\FunctionTok{format}\NormalTok{(dist\_paov}\SpecialCharTok{$}\StringTok{\textasciigrave{}}\AttributeTok{Pr(\textgreater{}F)}\StringTok{\textasciigrave{}}\NormalTok{[}\FunctionTok{rownames}\NormalTok{(dist\_paov) }\SpecialCharTok{==} \StringTok{"condition"}\NormalTok{], }\AttributeTok{digits =} \DecValTok{2}\NormalTok{),}
    \StringTok{"}\SpecialCharTok{\textbackslash{}n}\StringTok{ line:condition: "}\NormalTok{,}
    \FunctionTok{format}\NormalTok{(dist\_paov}\SpecialCharTok{$}\NormalTok{R2[}\FunctionTok{rownames}\NormalTok{(dist\_paov) }\SpecialCharTok{==} \StringTok{"line:condition"}\NormalTok{] }\SpecialCharTok{*} \DecValTok{100}\NormalTok{, }\AttributeTok{digits =} \DecValTok{2}\NormalTok{, }\AttributeTok{nsmall =} \DecValTok{1}\NormalTok{),}
    \StringTok{"\% var, P="}\NormalTok{, }\FunctionTok{format}\NormalTok{(dist\_paov}\SpecialCharTok{$}\StringTok{\textasciigrave{}}\AttributeTok{Pr(\textgreater{}F)}\StringTok{\textasciigrave{}}\NormalTok{[}\FunctionTok{rownames}\NormalTok{(dist\_paov) }\SpecialCharTok{==} \StringTok{"line:condition"}\NormalTok{], }\AttributeTok{digits =} \DecValTok{2}\NormalTok{)}
\NormalTok{  )}
\NormalTok{  p0 }\OtherTok{\textless{}{-}}\NormalTok{ phyloseq}\SpecialCharTok{::}\FunctionTok{plot\_ordination}\NormalTok{(phy\_sub, CAP\_bray, }\AttributeTok{title =}\NormalTok{ title, }\AttributeTok{axes =} \FunctionTok{c}\NormalTok{(}\DecValTok{1}\NormalTok{, }\DecValTok{2}\NormalTok{))}
\NormalTok{  p0}\SpecialCharTok{$}\NormalTok{layers }\OtherTok{\textless{}{-}}\NormalTok{ p0}\SpecialCharTok{$}\NormalTok{layers[}\SpecialCharTok{{-}}\DecValTok{1}\NormalTok{]}
\NormalTok{  p0 }\OtherTok{\textless{}{-}}\NormalTok{ p0 }\SpecialCharTok{+}
    \FunctionTok{stat\_ellipse}\NormalTok{(}
      \AttributeTok{geom =} \StringTok{"polygon"}\NormalTok{, }\AttributeTok{type =} \StringTok{"t"}\NormalTok{, }\AttributeTok{level =} \FloatTok{0.95}\NormalTok{, }\AttributeTok{size =} \FloatTok{0.4}\NormalTok{,}
      \FunctionTok{aes}\NormalTok{(}\AttributeTok{group =}\NormalTok{ condition, }\AttributeTok{color =}\NormalTok{ condition, }\AttributeTok{fill =}\NormalTok{ condition), }
      \AttributeTok{alpha =} \FloatTok{0.01}
\NormalTok{    ) }\SpecialCharTok{+}
    \FunctionTok{scale\_fill\_manual}\NormalTok{(}\AttributeTok{values =}\NormalTok{ colors) }\SpecialCharTok{+}
    \FunctionTok{scale\_color\_manual}\NormalTok{(}\AttributeTok{values =}\NormalTok{ colors) }\SpecialCharTok{+}
    \FunctionTok{new\_scale\_color}\NormalTok{() }\SpecialCharTok{+}
    \FunctionTok{geom\_point}\NormalTok{(}\FunctionTok{aes}\NormalTok{(}\AttributeTok{shape =} \FunctionTok{factor}\NormalTok{(line), }\AttributeTok{color =}\NormalTok{ condition), }\AttributeTok{size =} \FloatTok{1.0}\NormalTok{) }\SpecialCharTok{+}
    \FunctionTok{scale\_shape\_manual}\NormalTok{(}\AttributeTok{values =} \FunctionTok{c}\NormalTok{(}\DecValTok{16}\NormalTok{, }\DecValTok{2}\NormalTok{)) }\SpecialCharTok{+}
    \FunctionTok{scale\_color\_manual}\NormalTok{(}\AttributeTok{values =}\NormalTok{ colors) }\SpecialCharTok{+}
    \FunctionTok{labs}\NormalTok{(}
      \AttributeTok{subtitle =}\NormalTok{ score,}
      \AttributeTok{color =} \StringTok{"Microbiome"}\NormalTok{, }\AttributeTok{fill =} \StringTok{"Microbiome"}\NormalTok{, }
      \AttributeTok{shape =} \StringTok{"Line"}\NormalTok{, }\AttributeTok{alpha =} \StringTok{"Line"}
\NormalTok{    ) }\SpecialCharTok{+}
    \FunctionTok{theme\_bw}\NormalTok{() }\SpecialCharTok{+}
    \FunctionTok{theme}\NormalTok{(}
      \AttributeTok{text =} \FunctionTok{element\_text}\NormalTok{(}\AttributeTok{size =} \DecValTok{7}\NormalTok{),}
      \AttributeTok{axis.title =} \FunctionTok{element\_text}\NormalTok{(}\AttributeTok{size =} \DecValTok{7}\NormalTok{),}
      \AttributeTok{axis.text =} \FunctionTok{element\_text}\NormalTok{(}\AttributeTok{size =} \DecValTok{7}\NormalTok{),}
      \AttributeTok{legend.title =} \FunctionTok{element\_text}\NormalTok{(}\AttributeTok{size =} \DecValTok{7}\NormalTok{),}
      \AttributeTok{legend.text =} \FunctionTok{element\_text}\NormalTok{(}\AttributeTok{size =} \DecValTok{7}\NormalTok{),}
      \AttributeTok{plot.title =} \FunctionTok{element\_text}\NormalTok{(}\AttributeTok{size =} \DecValTok{7}\NormalTok{),}
      \AttributeTok{plot.subtitle =} \FunctionTok{element\_text}\NormalTok{(}\AttributeTok{size =} \DecValTok{7}\NormalTok{)}
\NormalTok{    )}
  \FunctionTok{return}\NormalTok{(p0)}
\NormalTok{\}}

\NormalTok{fig\_fun }\OtherTok{\textless{}{-}} \FunctionTok{process\_CAP}\NormalTok{(fPHYSEQ, fdist\_paov, }\StringTok{"Fungi"}\NormalTok{, level\_cols\_condition)}
\NormalTok{fig\_oom }\OtherTok{\textless{}{-}} \FunctionTok{process\_CAP}\NormalTok{(oPHYSEQ, odist\_paov, }\StringTok{"Oomycetes"}\NormalTok{, level\_cols\_condition)}

\CommentTok{\# ==============================}
\CommentTok{\# Combine and Save Plots {-}{-}{-}{-}{-}{-}{-}{-}{-}{-}{-}{-}{-}{-}{-}{-}{-}{-}{-}{-}{-}{-}{-}{-}{-}{-}{-}}
\CommentTok{\# ==============================}

\CommentTok{\# Combine all three figures}
\NormalTok{fig\_5}\FloatTok{.2} \OtherTok{\textless{}{-}} \FunctionTok{plot\_grid}\NormalTok{(}
\NormalTok{  fig\_bac }\SpecialCharTok{+} \FunctionTok{theme}\NormalTok{(}\AttributeTok{legend.position =} \StringTok{"none"}\NormalTok{),}
\NormalTok{  fig\_fun }\SpecialCharTok{+} \FunctionTok{theme}\NormalTok{(}\AttributeTok{legend.position =} \StringTok{"none"}\NormalTok{),}
\NormalTok{  fig\_oom }\SpecialCharTok{+} \FunctionTok{theme}\NormalTok{(}\AttributeTok{legend.position =} \StringTok{"none"}\NormalTok{),}
  \AttributeTok{align =} \StringTok{"hv"}\NormalTok{, }
  \AttributeTok{nrow =} \DecValTok{1}
\NormalTok{)}

\CommentTok{\# Extract legend from the Bacteria plot}
\NormalTok{legend }\OtherTok{\textless{}{-}} \FunctionTok{get\_legend}\NormalTok{(fig\_bac }\SpecialCharTok{+} \FunctionTok{theme}\NormalTok{(}\AttributeTok{legend.position =} \StringTok{"bottom"}\NormalTok{) }\SpecialCharTok{+} \FunctionTok{labs}\NormalTok{(}\AttributeTok{fill =} \StringTok{"line"}\NormalTok{))}

\CommentTok{\# Add legend to the combined plot}
\NormalTok{CAP\_plot }\OtherTok{\textless{}{-}} \FunctionTok{plot\_grid}\NormalTok{(fig\_5}\FloatTok{.2}\NormalTok{, legend, }\AttributeTok{nrow =} \DecValTok{2}\NormalTok{, }\AttributeTok{rel\_heights =} \FunctionTok{c}\NormalTok{(}\DecValTok{1}\NormalTok{, }\FloatTok{0.1}\NormalTok{))}

\FunctionTok{print}\NormalTok{(CAP\_plot)}
\CommentTok{\# Save the combined plot}
\FunctionTok{ggsave}\NormalTok{(}
  \StringTok{"CAP\_condition\_color.pdf"}\NormalTok{, }
  \AttributeTok{plot =}\NormalTok{ CAP\_plot, }
  \AttributeTok{width =} \FloatTok{17.5}\NormalTok{, }
  \AttributeTok{height =} \DecValTok{8}\NormalTok{, }
  \AttributeTok{units =} \StringTok{"cm"}\NormalTok{, }
  \AttributeTok{path =} \StringTok{"Output/Figures/"}
\NormalTok{)}
\end{Highlighting}
\end{Shaded}

\begin{figure}[H]

\centering{

\pandocbounded{\includegraphics[keepaspectratio]{05_Microbiome_analysis_files/figure-pdf/fig-CAP_plot-1.pdf}}

}

\caption{\label{fig-CAP_plot}CAP plot for bacteria, fungi and
oomycetes.}

\end{figure}%

\section{Session information}\label{session-information}

\textbf{Purpose}: Print the session information, including details about
the R version, loaded packages, and the system environment. This is
useful for reproducibility, as it documents the computational
environment in which the analysis was performed.

\begin{Shaded}
\begin{Highlighting}[]
\CommentTok{\# Get session information}
\NormalTok{SI }\OtherTok{\textless{}{-}} \FunctionTok{sessionInfo}\NormalTok{()}

\CommentTok{\# Extract the attached packages and their versions}
\NormalTok{attached\_packages }\OtherTok{\textless{}{-}}\NormalTok{ SI}\SpecialCharTok{$}\NormalTok{otherPkgs}

\CommentTok{\# Print the R version}
\FunctionTok{cat}\NormalTok{(}\StringTok{"R version:"}\NormalTok{, SI}\SpecialCharTok{$}\NormalTok{R.version}\SpecialCharTok{$}\NormalTok{version.string, }\StringTok{"}\SpecialCharTok{\textbackslash{}n\textbackslash{}n}\StringTok{"}\NormalTok{)}
\end{Highlighting}
\end{Shaded}

\begin{verbatim}
R version: R version 4.4.2 (2024-10-31) 
\end{verbatim}

\begin{Shaded}
\begin{Highlighting}[]
\CommentTok{\# Print the attached packages and their versions}
\FunctionTok{cat}\NormalTok{(}\StringTok{"Attached packages:}\SpecialCharTok{\textbackslash{}n}\StringTok{"}\NormalTok{)}
\end{Highlighting}
\end{Shaded}

\begin{verbatim}
Attached packages:
\end{verbatim}

\begin{Shaded}
\begin{Highlighting}[]
\FunctionTok{invisible}\NormalTok{(}\FunctionTok{lapply}\NormalTok{(attached\_packages, }\ControlFlowTok{function}\NormalTok{(pkg) \{}
  \FunctionTok{cat}\NormalTok{(pkg}\SpecialCharTok{$}\NormalTok{Package, }\StringTok{"{-}"}\NormalTok{, pkg}\SpecialCharTok{$}\NormalTok{Version, }\StringTok{"}\SpecialCharTok{\textbackslash{}n}\StringTok{"}\NormalTok{)}
\NormalTok{\}))}
\end{Highlighting}
\end{Shaded}

\begin{verbatim}
lubridate - 1.9.4 
purrr - 1.0.2 
readr - 2.1.5 
tidyverse - 2.0.0 
ggpubr - 0.6.0 
ggnewscale - 0.5.0 
knitr - 1.49 
sjstats - 0.19.0 
metagenomeSeq - 1.48.1 
glmnet - 4.1-8 
Matrix - 1.7-1 
limma - 3.62.1 
Biobase - 2.66.0 
Maaslin2 - 1.20.0 
ANCOMBC - 2.8.0 
ALDEx2 - 1.38.0 
latticeExtra - 0.6-30 
zCompositions - 1.5.0-4 
truncnorm - 1.0-9 
NADA - 1.6-1.1 
multcomp - 1.4-26 
TH.data - 1.1-2 
MASS - 7.3-64 
survival - 3.8-3 
mvtnorm - 1.3-2 
emmeans - 1.10.6 
hilldiv - 1.5.1 
ggh4x - 0.3.0 
RColorBrewer - 1.1-3 
pander - 0.6.5 
cowplot - 1.1.3 
ggpattern - 1.1.3 
stringr - 1.5.1 
forcats - 1.0.0 
tibble - 3.2.1 
tidyr - 1.3.1 
readxl - 1.4.3 
DECIPHER - 3.2.0 
Biostrings - 2.74.1 
GenomeInfoDb - 1.42.1 
XVector - 0.46.0 
IRanges - 2.40.1 
S4Vectors - 0.44.0 
BiocGenerics - 0.52.0 
vegan - 2.6-8 
lattice - 0.22-6 
permute - 0.9-7 
phyloseq - 1.50.0 
compositions - 2.0-8 
ggplot2 - 3.5.1 
dplyr - 1.1.4 
nlme - 3.1-166 
rstudioapi - 0.17.1 
\end{verbatim}

\begin{Shaded}
\begin{Highlighting}[]
\NormalTok{knitr}\SpecialCharTok{::}\FunctionTok{write\_bib}\NormalTok{(}\AttributeTok{file =} \StringTok{\textquotesingle{}packages\_norm.bib\textquotesingle{}}\NormalTok{)}
\end{Highlighting}
\end{Shaded}





\end{document}
